% ============================================================
%   Bachelor Thesis (PFG) -- Universidad de Deusto
%   Template: memoriaTFG.cls (2024\textendash 25)
%   Project : CALA-SafeRL
% ============================================================

\documentclass{TFG-Deusto-LaTeX-main/plantilla/memoriaTFG}

% ------------------------------------------------------------
% Extra packages required by the body of this thesis
% (amsmath: align/gather/cases; booktabs: professional tables; siunitx: physical units)
% ------------------------------------------------------------
\usepackage{amsmath}
\usepackage{booktabs}
\usepackage{multirow}   % grouped row labels in the ablation tables (Ch. 8)
\usepackage{siunitx}

% ------------------------------------------------------------
% Metadata (also consumed by hyperref for PDF bookmarks)
% ------------------------------------------------------------
\title{Safe Reinforcement Learning for Autonomous Driving: A Shield-Based Framework on the CARLA Simulator}
\autor{Isaac Esteban Herbozo}
\titulacion{Grado en Ingenier\'ia Inform\'atica}
\director{[Project Director]}

\resumen{%
This work presents CALA-SafeRL, a modular framework for Safe Reinforcement Learning applied to
autonomous driving in the CARLA simulator. The framework couples a Proximal Policy Optimization
(PPO) agent with two interchangeable safety shields, a dense reward shaper and a performance-driven
curriculum, and is designed as a composable chain of Gymnasium wrappers. A central finding is that
shielded agents become \emph{dependent} on the shield: they drive competently while it is active
but fail when it is removed. We show that loosening the shield with a Large-Language-Model
meta-controller does not reduce this dependence, whereas using the shield as a behavioural-cloning
\emph{teacher} of its steering corrections does, to the point where the trained agent drives as
well without the shield as with it. Finally, we close the residual collision gap left by weaning
with a dense reward learned offline from a language model's safety preferences, delivered as
policy-invariant potential-based shaping; once perception and exposure are de-confounded, this
reward improves lane-keeping but does not, on its own, teach braking: a clean negative result
that locates the remaining gap in longitudinal control.%
}
\descriptores{Safe Reinforcement Learning, Autonomous Driving, Safety Shield, Proximal Policy Optimization, CARLA Simulator}

% ------------------------------------------------------------
% Document
% ------------------------------------------------------------
\begin{document}

\selectlanguage{english}

% ============================================================
%   FRONTMATTER  (roman page numbers)
% ============================================================
\frontmatter

% Official cover pages (unsigned + signed)
% Replace with your own portada PDFs if needed (PFG_IngenieriaInformatica.pdf for the unsigned one)
\includepdf[pages=1]{imgs/portada.pdf}
\includepdf[pages=1]{imgs/portada.pdf}

% ============================================================
%   Abstract and Keywords  (English)
% ============================================================
%
% The class macros \hacerresumen and \resumenPersonalizado hard-code
% the Spanish word ``Resumen''. Because this document is written in
% English, the abstract is laid out manually here while the metadata
% macros \resumen and \descriptores (defined in main.tex) are still
% picked up by hyperref for PDF keywords.
% ------------------------------------------------------------

\cleardoublepage
\pagenumbering{roman}
\pagestyle{headings}

\vspace*{3cm}

{\bfseries\Large Abstract\par}
\vspace{0.5cm}

\noindent
This work presents CALA-SafeRL, a modular framework for Safe Reinforcement Learning applied to
autonomous driving in the CARLA simulator. The central challenge addressed is how to train a
Proximal Policy Optimization (PPO) agent that learns to drive under realistic traffic conditions
without catastrophic failures during the learning process itself. The framework is structured as a
composable chain of Gymnasium wrappers, allowing the agent, two interchangeable safety shields, a
dense reward shaper and the simulator to be combined with minimal coupling. The first shield
performs a lightweight threshold-based verification on semantic LIDAR readings together with
CARLA's Waypoint API; the second implements an adaptive-horizon strategy, predicting the vehicle's
future trajectory with a kinematic bicycle model calibrated for a Tesla Model 3 and selecting the
least invasive corrective action from a prioritised candidate set. The agent's observation space
has 739 dimensions, combining three semantic LIDAR channels with lane, lane-marking, vehicle and
route features derived from the Waypoint API. Training stability is improved through running-mean observation
normalisation, return scaling, Generalised Advantage Estimation, clipped surrogate updates with a
target-KL early stop, and a performance-driven curriculum that progressively introduces traffic
density. The framework is accompanied by a live metrics pipeline built on SQLite and a Streamlit
dashboard. Experiments are structured around the Town04 highway map and systematically compare the
three shielding configurations against baseline PPO. The resulting system is open, reproducible and
can serve as a foundation for further research into safety-aware policy learning.

\vspace{2cm}

\cleardoublepage


\tableofcontents
\listoffigures
\listoftables

% ============================================================
%   MAINMATTER  (arabic page numbers)
% ============================================================
\mainmatter

% ============================================================
%   Chapter 1  \textemdash\  Introduction
% ============================================================

\chapter{Introduction}\label{ch:introduction}

\section{Context and motivation}\label{sec:intro-context}

Road transport is one of the dominant causes of avoidable injury and death worldwide. According
to the World Health Organization, approximately 1.19 million people die every year as a direct
consequence of road traffic crashes, with human error being the leading contributing
factor~\cite{who2023road}. This figure has motivated a sustained effort, both in industry and in
academia, to automate vehicle control up to the point where the driver can be completely
displaced from the control loop. The SAE J3016 standard formalises this progression in six
levels, from no automation (Level 0) to full self-driving with no operational design domain
restrictions (Level 5)~\cite{sae2021j3016}. While most commercially deployed systems remain at
Level 2, research platforms increasingly target Levels 4 and 5, where the vehicle must cope with
the full variety of traffic situations, weather conditions and road topologies encountered in
real-world operation.

Among the techniques that have received growing attention for high-autonomy driving,
Reinforcement Learning (RL)~\cite{sutton2018rl} occupies a distinctive position. Classical
model-predictive and rule-based controllers require a hand-crafted description of every
situation and a correspondingly explicit control law; in contrast, an RL agent learns a policy
end-to-end by interacting with an environment and maximising a scalar reward signal. Recent
surveys of deep RL for autonomous driving report promising results on lane keeping, overtaking,
intersection handling and dense urban scenarios~\cite{kiran2022drlsurvey,aradi2022survey}.

The main obstacle to deploying these systems, however, is not their peak performance but their
behaviour \emph{during learning}. Vanilla RL relies on stochastic exploration, which is
incompatible with the safety requirements of a vehicle operating on public roads: even a single
exploratory collision is unacceptable. The research field generally referred to as
\emph{Safe Reinforcement Learning} (Safe RL) addresses this tension through multiple
strategies~\cite{garcia2015comprehensive}: encoding constraints in the reward function,
constraining the policy optimisation problem itself, or intercepting unsafe actions with a
runtime verifier known as a \emph{safety shield}~\cite{alshiekh2018shielding}. This project
belongs to the latter family: a PPO agent~\cite{schulman2017ppo} is trained with two alternative
shields that monitor its proposed actions, reason about their near-future consequences, and
substitute them when the resulting trajectory is judged unsafe. The shield layer runs entirely
during training, so safety violations are not merely discouraged by the reward function but
actively prevented at the actuator.

\section{Problem statement}\label{sec:intro-problem}

The central problem addressed in this thesis can be stated as follows. A vehicle, represented as
an RL agent, drives inside the CARLA simulator~\cite{dosovitskiy2017carla}. The agent observes a
partial representation of the environment through virtual sensors and a limited number of
kinematic variables, and must produce low-level control commands (steering, throttle and brake).
The goal is to learn a policy that

\begin{enumerate}
  \item \emph{maximises} a task reward that encodes progress, speed compliance and lane keeping,
  \item \emph{minimises} unsafe events such as collisions, off-road excursions and stalls, and
  \item \emph{tolerates} the shift in traffic density that occurs as the scenario is progressively
        complicated during training, and
  \item \emph{generalises} beyond the shield, so that the learned policy remains safe even when the
        shield is removed at deployment.
\end{enumerate}

The research question that guides the work is then: \emph{can a lightweight, interchangeable
safety-shield layer, combined with a dense reward shaper and a performance-driven curriculum,
train a PPO agent that is competitive with unconstrained PPO in terms of reward while keeping
crash and off-road rates close to zero during learning?} A second question, which emerges from the
first, is whether the very mechanism that enforces safety during training, the shield, induces
a \emph{dependence} that the agent cannot shed, and, if so, how that dependence can be measured and
reduced; this question is taken up in Chapter~\ref{ch:dependence}.

\section{Objectives}\label{sec:intro-objectives}

\subsection*{General objective}

To design, implement and empirically analyse a modular Safe RL framework for autonomous driving
on the CARLA simulator, combining a PPO agent with interchangeable safety shields, a dense
reward shaper and a curriculum manager.

\subsection*{Specific objectives}

\begin{itemize}
  \item[\textbf{SO1.}] Build a Gymnasium-compatible wrapper of the CARLA Python API that exposes a
        739-dimensional observation space (three semantic LIDAR channels plus lane-marking, lane,
        vehicle and route features) and a continuous two-dimensional action space, running in
        synchronous mode at 20\,Hz.
  \item[\textbf{SO2.}] Implement a PPO agent with generalised advantage estimation, clipped
        surrogate updates, value clipping, a target-KL early stop, cosine learning-rate annealing
        and online normalisation of the vector observation block.
  \item[\textbf{SO3.}] Develop a basic safety shield that combines LIDAR thresholds with
        information from CARLA's Waypoint API and applies a proportional correction to steering
        and brake whenever the agent's action is deemed unsafe.
  \item[\textbf{SO4.}] Develop an adaptive-horizon safety shield that classifies the current
        situation into three risk levels, predicts the vehicle's trajectory with a kinematic
        bicycle model, and selects the least invasive corrective action from a priority-ordered
        candidate set.
  \item[\textbf{SO5.}] Design a dense reward shaper that combines the sparse CARLA rewards with
        components for velocity tracking, lane centring, heading alignment, smoothness, progress,
        idle avoidance and lane-change tolerance.
  \item[\textbf{SO6.}] Build a performance-driven curriculum that progressively introduces
        non-playable vehicles (NPCs) into the scenario, with automatic rollback to earlier
        stages on performance collapse.
  \item[\textbf{SO7.}] Provide a live metrics pipeline based on SQLite (WAL mode) and a Streamlit
        dashboard so that each training run can be inspected while it is still in progress.
  \item[\textbf{SO8.}] Measure and reduce the agent's \emph{dependence} on the shield: instrument a
        shield-off evaluation probe, and study two strategies: a Large-Language-Model
        meta-controller that loosens the shield over training, and a behavioural-cloning scheme
        that turns the shield into a teacher, to train a policy that remains safe when the
        shield is removed at deployment.
\end{itemize}

\section{Scope and limitations}\label{sec:intro-scope}

The scope of this project is deliberately bounded to keep the contribution verifiable and
reproducible within the duration of a Bachelor's thesis.

\begin{itemize}
  \item The simulator version is \emph{CARLA 0.9.16}. Compatibility with earlier or future
        versions is not guaranteed.
  \item The main benchmark map is \emph{Town04}, a highway-oriented layout that is well suited
        for lane-keeping and overtaking experiments. Urban maps (Town01--Town03) and the maximum
        difficulty map (Town05) are supported by the framework but not systematically evaluated.
  \item Two shielding strategies are developed and compared: \emph{basic} (threshold + Waypoint
        API correction) and \emph{adaptive-horizon} (trajectory prediction with a kinematic
        bicycle model). Alternatives based on reachability analysis, model-predictive safety
        filters or formal verification are left for future work.
  \item The vehicle model used by the adaptive shield is the kinematic bicycle
        model~\cite{rajamani2012vehicle}. A dynamic model with tyre slip is not considered; at
        the speeds and lateral accelerations encountered in the experiments the kinematic
        approximation is a reasonable trade-off between fidelity and computational cost.
  \item Perception is based on semantic LIDAR plus CARLA's Waypoint API. Camera-based perception,
        sensor fusion and learned occupancy grids are not part of the observation space.
  \item The framework trains a \emph{single ego vehicle}. Multi-agent safe RL is an explicit
        direction for future extensions.
\end{itemize}

\section{Methodology overview}\label{sec:intro-methodology}

The project follows an iterative engineering methodology structured around four phases: design,
implementation, training and analysis, which are revisited in successive loops as new empirical
evidence becomes available.

\begin{enumerate}
  \item \textbf{Design.} The software architecture is expressed as a layered wrapper chain
        on top of Gymnasium. Each layer (environment, shield, reward shaper) has a single
        responsibility and a well-defined interface, so that components can be swapped or
        reconfigured without touching the others.
  \item \textbf{Implementation.} The Python code base is organised into self-contained modules
        (\texttt{src/CARLA/Env}, \texttt{src/PPO}, \texttt{src/Adaptative\_Shield},
        \texttt{src/reward\_shaper.py}, \texttt{src/curriculumManager.py}, \texttt{src/Metrics}).
        Automated style checks with Ruff and targeted unit tests on the reward shaper prevent
        regressions.
  \item \textbf{Training.} Each experimental configuration is launched via \texttt{main\_train.py}
        with explicit command-line flags. The agent writes metrics to a per-run SQLite database
        and checkpoints every $200$ episodes.
  \item \textbf{Analysis.} The \texttt{dataVisualizer.py} Streamlit dashboard is used to inspect
        reward curves, safety rates, curriculum transitions and shield activations during and
        after training. Observed pathologies are translated into new design iterations.
\end{enumerate}

\section{Document structure}\label{sec:intro-structure}

The remainder of this document is organised as follows.

\begin{description}
  \item[Chapter~\ref{ch:sota}] reviews the theoretical background. It introduces Markov Decision
        Processes, derives the PPO objective from the policy gradient, describes the Generalised
        Advantage Estimator, surveys the main families of Safe RL, formalises the kinematic
        bicycle model and summarises the features of the CARLA simulator that are relevant to
        this project.
  \item[Chapter~\ref{ch:requirements}] captures the functional and non-functional requirements
        that the framework must satisfy, together with the technical constraints imposed by the
        simulator, the hardware and the Bachelor-thesis timeline.
  \item[Chapter~\ref{ch:design}] describes the system design: the wrapper chain, the structure
        of the observation and action spaces, the Actor\textendash Critic network, the two
        safety shields, the reward shaper and the curriculum manager.
  \item[Chapter~\ref{ch:implementation}] details the implementation: technology stack, project
        layout, key algorithms in pseudocode, training and evaluation pipelines, and
        reproducibility considerations.
  \item[Chapter~\ref{ch:experimentation}] presents the experimental protocol and a structured
        template to report the empirical results obtained on Town04.
  \item[Chapter~\ref{ch:dependence}] studies the agent's dependence on the shield: it introduces
        the shield-off evaluation probe, reports a negative result for a Large-Language-Model
        meta-controller that loosens the shield, and a positive result for a behavioural-cloning
        scheme that turns the shield into a teacher of unaided safe driving.
  \item[Chapter~\ref{ch:conclusions}] summarises the contributions, evaluates the fulfilment of
        the objectives enumerated above, and proposes future work.
  \item[Appendices~\ref{app:hyperparameters}, \ref{app:observation} and \ref{app:reproducibility}]
        provide the complete list of default hyperparameters, a full breakdown of the observation
        space and a reproducibility guide, respectively.
\end{description}

% ============================================================
%   Chapter 2  --  Project Planning and Management
% ============================================================

\chapter{Project Planning and Management}\label{ch:planning}

The previous chapter framed CALA-SafeRL as a research problem. This chapter frames it as a
\emph{managed engineering project}. It maps the work onto the five canonical project-management
stages (Section~\ref{sec:plan-approach}), decomposes it into a work breakdown structure with an
effort estimate (Section~\ref{sec:plan-wbs}), records the milestones
(Section~\ref{sec:plan-milestones}), the task precedences and the resulting schedule and Gantt
chart (Section~\ref{sec:plan-schedule}), the human-resources plan
(Section~\ref{sec:plan-hr}) and, finally, the risk register and the deviations that actually
occurred (Section~\ref{sec:plan-risk}). The plan presented here is the \emph{baseline} against
which execution was monitored; because the engineering methodology is iterative
(Section~\ref{sec:intro-methodology}), the baseline was re-planned at the end of each
design\,$\to$\,implementation\,$\to$\,training\,$\to$\,analysis loop rather than followed rigidly.

\section{Project-management approach}\label{sec:plan-approach}

The project is organised around the five standard project-management stages (definition,
planning, execution, monitoring and control, and closure), adapted to a research-and-development
context in which the requirements are partly discovered during execution.
Table~\ref{tab:stages} maps each stage to the concrete activities and artefacts of CALA-SafeRL.
A point worth highlighting is the \emph{monitoring and control} stage: in a conventional software
project it is realised through status reports, whereas here it is realised through engineering
instrumentation: the live metrics dashboard (Section~\ref{sec:design-metrics}), the
shield-off probe (Section~\ref{sec:dep-probe}) and the KL and curriculum gates, which provide
continuous, quantitative feedback on whether the work is converging.

\begin{table}[H]
  \centering
  \caption{The five project-management stages mapped to concrete activities and artefacts of
           CALA-SafeRL.}
  \label{tab:stages}
  \small
  \begin{tabular}{@{}lp{6.0cm}p{4.6cm}@{}}
    \toprule
    \textbf{Stage} & \textbf{Activities in this project} & \textbf{Artefacts} \\
    \midrule
    Definition            & Problem framing, research question, scope and objectives
                          & Sections~\ref{sec:intro-context}--\ref{sec:intro-scope} \\
    Planning              & Work breakdown, schedule, resource and risk plans
                          & This chapter \\
    Execution             & Iterative design, implementation, training and analysis loops
                          & Chapters~\ref{ch:design}--\ref{ch:dependence} \\
    Monitoring \& Control & Live metrics dashboard, shield-off probe, KL and curriculum gates,
                            rolling re-planning
                          & Sections~\ref{sec:design-metrics}, \ref{sec:dep-probe} \\
    Closure               & Final evaluation, memoir, defense, reproducibility package
                          & Chapter~\ref{ch:conclusions}, Appendix~\ref{app:reproducibility} \\
    \bottomrule
  \end{tabular}
\end{table}

\section{Work breakdown structure and effort estimation}\label{sec:plan-wbs}

The project is decomposed into eleven work packages (WP0--WP10). WP0 (management) and WP10
(writing) run continuously alongside the technical packages; the remaining packages map directly
onto the specific objectives SO1--SO8 of Section~\ref{sec:intro-objectives}.
Table~\ref{tab:wbs} lists each package, its main tasks, the objective(s) it serves and the
estimated effort. The total of \(550\,\mathrm{h}\) is the author's post-hoc reconstruction of the
effort actually invested; it is consistent with, and somewhat above, the nominal workload of a
\(12\)-ECTS Bachelor's thesis (\(\approx 300\text{--}360\,\mathrm{h}\)), the difference reflecting
the research-heavy nature of the shield-dependence study (WP9) and the large number of training
runs required (WP8).

\begin{table}[H]
  \centering
  \caption{Work breakdown structure and effort estimation. The linked specific objectives (SO)
           are those of Section~\ref{sec:intro-objectives}.}
  \label{tab:wbs}
  \small
  \begin{tabular}{@{}llp{6.0cm}lr@{}}
    \toprule
    \textbf{ID} & \textbf{Work package} & \textbf{Main tasks} & \textbf{Obj.} & \textbf{Hours} \\
    \midrule
    WP0  & Project management        & Planning, scheduling, supervision meetings, risk tracking, version control & -- & 35 \\
    WP1  & Research \& state of the art & RL and Safe-RL literature, shielding, CARLA, bicycle model & -- & 50 \\
    WP2  & Environment \& wrappers   & \texttt{CarlaEnv}, synchronous mode, semantic LIDAR, Waypoint features, 739-d observation & SO1 & 60 \\
    WP3  & PPO agent                 & Actor--Critic network, GAE, clipped surrogate, KL early stop, normalisation & SO2 & 55 \\
    WP4  & Safety shields            & Basic shield, adaptive-horizon shield, bicycle model, shield credit assignment & SO3, SO4 & 70 \\
    WP5  & Reward shaper             & Dense reward components, speed/lane gates, unit tests & SO5 & 30 \\
    WP6  & Curriculum manager        & Five-stage NPC ramp, rollback rule, in-place update & SO6 & 22 \\
    WP7  & Metrics \& dashboard      & SQLite WAL logger, Streamlit/Plotly dashboard & SO7 & 25 \\
    WP8  & Experimentation           & Training runs, shield ablations, evaluation protocol, analysis & -- & 55 \\
    WP9  & Shield-dependence study   & Shield-off probe, LLM meta-controller, behavioural-cloning teacher, failure analysis & SO8 & 68 \\
    WP10 & Thesis writing \& defense & Memoir, figures, reproducibility guide, defense preparation & -- & 80 \\
    \midrule
    \multicolumn{4}{r}{\textbf{Total}} & \textbf{550} \\
    \bottomrule
  \end{tabular}
\end{table}

\section{Milestones}\label{sec:plan-milestones}

Six milestones mark the closure of a major block of work and gate the start of the next.
Table~\ref{tab:milestones} lists them with their target dates and deliverables. Milestones M3--M5
correspond to the three empirical achievements of the project: a working shielded agent, the main
shield ablation, and the resolution of the shield-dependence problem.

\begin{table}[H]
  \centering
  \caption{Project milestones.}
  \label{tab:milestones}
  \small
  \begin{tabular}{@{}llp{7.4cm}@{}}
    \toprule
    \textbf{ID} & \textbf{Target} & \textbf{Milestone / deliverable} \\
    \midrule
    M1 & Sep 2025 & Project definition and baseline plan approved \\
    M2 & Nov 2025 & State of the art consolidated; core bibliography fixed \\
    M3 & Jan 2026 & Full wrapper chain and PPO agent operational; first shielded training run \\
    M4 & Mar 2026 & Main experimentation complete (shield ablation and curriculum) \\
    M5 & May 2026 & Shield-dependence study complete (LLM negative result, BC positive result) \\
    M6 & Jun 2026 & Thesis submitted and defended \\
    \bottomrule
  \end{tabular}
\end{table}

\section{Task scheduling and Gantt chart}\label{sec:plan-schedule}

The technical work packages form a mostly linear dependency chain: the environment (WP2) must
exist before the agent (WP3), both are needed before the shields (WP4), and all of the
components must be in place before systematic experimentation (WP8) and the dependence study
(WP9). Table~\ref{tab:precedence} records the precedences; the resulting \emph{critical path} is
\(\text{WP1}\to\text{WP2}\to\text{WP3}\to\text{WP4}\to\text{WP8}\to\text{WP9}\to\text{WP10}\).
The reward shaper, curriculum manager and metrics pipeline (WP5--WP7) are off the critical path
and were developed in parallel with the shields.

\begin{table}[H]
  \centering
  \caption{Task precedence. WP0 and WP10 run continuously alongside the technical packages.}
  \label{tab:precedence}
  \small
  \begin{tabular}{@{}lll@{}}
    \toprule
    \textbf{WP} & \textbf{Predecessors} & \textbf{Relationship} \\
    \midrule
    WP0  & --                           & Continuous \\
    WP1  & --                           & Start \\
    WP2  & WP1                          & Finish-to-start \\
    WP3  & WP2                          & Finish-to-start \\
    WP4  & WP2, WP3                     & Finish-to-start \\
    WP5  & WP2                          & Finish-to-start (parallel) \\
    WP6  & WP2                          & Finish-to-start (parallel) \\
    WP7  & WP2                          & Finish-to-start (parallel) \\
    WP8  & WP3, WP4, WP5, WP6, WP7      & Finish-to-start \\
    WP9  & WP8                          & Finish-to-start \\
    WP10 & WP1--WP9                     & Continuous \\
    \bottomrule
  \end{tabular}
\end{table}

Figure~\ref{fig:gantt} shows the baseline schedule as a Gantt chart over the ten months of the
project. The diamonds mark the milestones of Table~\ref{tab:milestones}. The chart makes the
front-loading of research (WP1) and the back-loading of writing (WP10) explicit, and shows the
two-month overlap between the close of experimentation (WP8) and the start of the dependence
study (WP9) that occurred in spring 2026.

\begin{figure}[H]
  \centering
  \includegraphics[width=\textwidth]{imgs/gantt_chart}
  \caption{Baseline Gantt chart of the project (September 2025 to June 2026). Bars are the
           work packages of Table~\ref{tab:wbs}; diamonds are the milestones of
           Table~\ref{tab:milestones}. Management (WP0) and writing (WP10) span the whole
           project; the technical critical path runs WP2, WP3, WP4, WP8 and WP9.}
  \label{fig:gantt}
\end{figure}

\section{Human-resources plan}\label{sec:plan-hr}

Although the project was carried out by a single author, the work spanned several professional
\emph{roles}, each with a distinct skill profile and, in the budget of Chapter~\ref{ch:budget}, a
distinct cost rate. Table~\ref{tab:hr-roles} distributes the \(550\,\mathrm{h}\) of
Table~\ref{tab:wbs} across these roles. The RL/ML Engineer role dominates, as expected for a
reinforcement-learning thesis, but the Software Engineer role is also substantial because the
CARLA wrapper, the metrics pipeline and the dashboard are non-trivial systems-engineering tasks.

\begin{table}[H]
  \centering
  \caption{Roles assumed by the author and their effort. The hourly rates are used in the budget
           (Chapter~\ref{ch:budget}).}
  \label{tab:hr-roles}
  \small
  \begin{tabular}{@{}lp{5.8cm}rr@{}}
    \toprule
    \textbf{Role} & \textbf{Main responsibilities} & \textbf{Hours} & \textbf{Rate (€/h)} \\
    \midrule
    Project Manager   & Planning, scheduling, risk management, supervision liaison      & 45  & 40 \\
    RL/ML Engineer    & RL research, PPO, shield logic, reward design, results analysis & 256 & 35 \\
    Software Engineer & CARLA wrapper, metrics pipeline, dashboard, tooling             & 134 & 30 \\
    QA/Test Engineer  & Unit tests, deterministic evaluation, training-run validation   & 45  & 25 \\
    Technical Writer  & Memoir, figures, documentation                                 & 70  & 22 \\
    \midrule
    \multicolumn{2}{@{}l}{\textbf{Total effort}} & \textbf{550} & \\
    \bottomrule
  \end{tabular}
\end{table}

The project also relied on a \emph{support team} that is not part of the author's effort but is
costed in Chapter~\ref{ch:budget}: the project Director, who supervised the work and reviewed the
memoir; the Bachelor-thesis Coordinator, responsible for the administrative process; and the
three-member defense tribunal.

\section{Risk management and deviations}\label{sec:plan-risk}

Table~\ref{tab:risks} is the project risk register. Because the plan is presented retrospectively,
the final column records whether each risk actually materialised. Four of the six risks did, and
the way they were handled is documented in the lessons learned of
Section~\ref{sec:concl-lessons}: the curriculum-change deadlock was resolved by updating the NPC
count in place rather than rebuilding the environment; the \texttt{log\_std} saturation was
resolved by the tanh-squashed parameterisation; the shortage of compute for multi-seed runs was
\emph{accepted} and reported honestly (a single seed with seed-aligned cross-configuration
evaluation); and the late-emerging LLM reward-exposure line was de-scoped to future work to
protect the core objectives. The \(5\%\) contingency reserve in the budget
(Chapter~\ref{ch:budget}) is the financial counterpart of this risk exposure.

\begin{table}[H]
  \centering
  \caption{Risk register. The final column records whether each risk materialised during
           execution.}
  \label{tab:risks}
  \small
  \begin{tabular}{@{}p{4.4cm}ccp{4.4cm}c@{}}
    \toprule
    \textbf{Risk} & \textbf{Lik.} & \textbf{Imp.} & \textbf{Mitigation} & \textbf{Status} \\
    \midrule
    CARLA server instability / crashes               & High & Med  & Synchronous headless mode; checkpoint every 200 episodes & Managed \\
    Deadlock when rebuilding the env on curriculum change & Med  & High & In-place \texttt{num\_npc} update, no reconnection & Resolved \\
    Training instability (\texttt{log\_std} saturation)   & Med  & High & tanh-squashed \texttt{log\_std} parameterisation & Resolved \\
    Insufficient compute for multi-seed runs         & High & Med  & Seed-aligned cross-configuration evaluation; honest reporting & Accepted \\
    Scope creep (LLM reward-exposure line)           & Med  & Med  & De-scoped to future work; core objectives protected & De-scoped \\
    Single-workstation hardware failure              & Low  & High & Frequent checkpointing; Git version control & Not hit \\
    \bottomrule
  \end{tabular}
\end{table}

With the schedule, resources and risks established, Chapter~\ref{ch:budget} turns this plan into a
cost estimate.

\clearpage

% ============================================================
%   Chapter 3  --  Project Budget
% ============================================================

\chapter{Project Budget}\label{ch:budget}

This chapter estimates the cost of CALA-SafeRL \emph{as if it had been carried out as a
professional engineering contract}, rather than as an unpaid academic exercise. The estimate
builds directly on the planning of Chapter~\ref{ch:planning}: the labour cost follows the
role-and-effort distribution of the human-resources plan (Table~\ref{tab:hr-roles}), and the
contingency reserve follows the risk register (Table~\ref{tab:risks}).
Section~\ref{sec:budget-assumptions} fixes the cost model and its assumptions, the subsequent
sections cost labour, equipment, software, energy and travel, and
Section~\ref{sec:budget-total} aggregates them into the total budget.

\section{Cost model and assumptions}\label{sec:budget-assumptions}

The estimate rests on the following assumptions, chosen to be defensible for a project executed
in Spain during the 2025\,--\,2026 academic year:

\begin{itemize}
  \item \textbf{Labour} is charged per role at the rates of Table~\ref{tab:hr-roles}
        (\(40\)\,€/h Project Manager, \(35\)\,€/h RL/ML Engineer, \(30\)\,€/h Software Engineer,
        \(25\)\,€/h QA/Test Engineer, \(22\)\,€/h Technical Writer), reflecting Spanish market
        rates for a junior-to-mid professional in each speciality.
  \item \textbf{Equipment} is \emph{pre-owned} and charged by linear amortisation,
        \(\text{cost}\times(\text{months used}/\text{useful life})\times\text{dedication}\),
        with a \(60\)-month useful life for the workstation, an \(11\)-month project span and an
        \(80\%\) dedication factor.
  \item \textbf{Electricity} is priced at \(0.18\)\,€/kWh, a representative Spanish domestic
        tariff for the period.
  \item \textbf{Indirect costs} (workspace, internet, administration) are taken as \(15\%\) of
        direct costs. A \(5\%\) \textbf{contingency} reserve covers the risks of
        Table~\ref{tab:risks}.
  \item \textbf{Value-added tax} (IVA) is applied at the Spanish general rate of \(21\%\) on the
        base budget.
\end{itemize}

\section{Labour costs}\label{sec:budget-labor}

Labour is by far the dominant cost, as is typical for a research-and-development project whose
main input is engineering time. Table~\ref{tab:labor-internal} costs the author's own effort by
role, and Table~\ref{tab:labor-support} costs the support team (Director, Coordinator and tribunal),
whose time is a real cost of the project even though it is not borne by the author.

\begin{table}[H]
  \centering
  \caption{Internal labour cost (project team). Hours taken from the human-resources plan
           (Table~\ref{tab:hr-roles}).}
  \label{tab:labor-internal}
  \small
  \begin{tabular}{@{}lccr@{}}
    \toprule
    \textbf{Role} & \textbf{Hours} & \textbf{Rate (€/h)} & \textbf{Cost (€)} \\
    \midrule
    Project Manager   & 45  & 40 & 1\,800.00 \\
    RL/ML Engineer    & 256 & 35 & 8\,960.00 \\
    Software Engineer & 134 & 30 & 4\,020.00 \\
    QA/Test Engineer  & 45  & 25 & 1\,125.00 \\
    Technical Writer  & 70  & 22 & 1\,540.00 \\
    \midrule
    \textbf{Subtotal} & \textbf{550} & & \textbf{17\,445.00} \\
    \bottomrule
  \end{tabular}
\end{table}

\begin{table}[H]
  \centering
  \caption{Support-team labour cost.}
  \label{tab:labor-support}
  \small
  \begin{tabular}{@{}lccr@{}}
    \toprule
    \textbf{Member} & \textbf{Hours} & \textbf{Rate (€/h)} & \textbf{Cost (€)} \\
    \midrule
    Project Director (supervisor) & 30 & 60 & 1\,800.00 \\
    TFG Coordinator               & 5  & 50 & 250.00 \\
    Defense tribunal (3 members)  & 12 & 50 & 600.00 \\
    \midrule
    \textbf{Subtotal} & \textbf{47} & & \textbf{2\,650.00} \\
    \bottomrule
  \end{tabular}
\end{table}

The combined labour cost is therefore \(17\,445.00 + 2\,650.00 = 20\,095.00\)~€.

\section{Equipment and materials}\label{sec:budget-equipment}

The project ran on a single pre-owned desktop workstation (the machine described in
Section~\ref{sec:req-constraints} and used for the experiments of
Chapter~\ref{ch:experimentation}). Table~\ref{tab:equipment} charges only the fraction of its
useful life consumed by the project. No consumables of note were required, and, as discussed
in Section~\ref{sec:budget-software}, the entire scientific software stack is open-source, so
its licensing cost is zero.

\begin{table}[H]
  \centering
  \caption{Amortised equipment cost.
           Amortisation \(=\text{cost}\times(\text{months used}/\text{useful life})\times\text{dedication}\).}
  \label{tab:equipment}
  \small
  \begin{tabular}{@{}lccccr@{}}
    \toprule
    \textbf{Item} & \textbf{Cost (€)} & \textbf{Life (mo)} & \textbf{Used (mo)} & \textbf{Ded.} & \textbf{Amort. (€)} \\
    \midrule
    Workstation (i9-9900K, RTX 2080 Ti, 32\,GB) & 2\,000.00 & 60 & 11 & 80\% & 293.33 \\
    Monitor and peripherals                     & 350.00    & 72 & 11 & 80\% & 42.78 \\
    \midrule
    \textbf{Subtotal} & & & & & \textbf{336.11} \\
    \bottomrule
  \end{tabular}
\end{table}

\section{Software and AI services}\label{sec:budget-software}

A deliberate design decision was to build the framework entirely on open-source software, which
keeps the licensing cost at zero and the work fully reproducible (Table~\ref{tab:software}). The
only paid software is the AI coding-and-writing assistant whose use is disclosed in
Appendix~\ref{app:ai-usage}. The local Gemma model (served through Ollama) and the free tier of
Gemini incur no per-call cost, with their compute accounted for under energy
(Section~\ref{sec:budget-energy}).

\begin{table}[H]
  \centering
  \caption{Software and AI-service costs. The scientific stack is fully open-source.}
  \label{tab:software}
  \small
  \begin{tabular}{@{}p{5.8cm}p{5.0cm}r@{}}
    \toprule
    \textbf{Item} & \textbf{Licence / basis} & \textbf{Cost (€)} \\
    \midrule
    CARLA, Gymnasium, PyTorch, NumPy, Matplotlib, Streamlit, Plotly, SQLite, Ruff, pytest
      & MIT / BSD-3 / Apache-2.0 / PSF / public domain & 0.00 \\
    AI coding/writing assistant (\(\approx 20\)\,€/mo \(\times\) 8 mo)
      & Commercial subscription (Appendix~\ref{app:ai-usage}) & 160.00 \\
    Local LLM (Ollama + Gemma), Gemini free tier
      & Local / free tier & 0.00 \\
    \midrule
    \textbf{Subtotal} & & \textbf{160.00} \\
    \bottomrule
  \end{tabular}
\end{table}

\section{Energy consumption}\label{sec:budget-energy}

Reinforcement-learning training is energy-intensive: each configuration requires hundreds of
episodes of CARLA simulation with the policy on the GPU. Table~\ref{tab:energy} estimates the
electricity consumed, separating heavy GPU training from lighter development and analysis.

\begin{table}[H]
  \centering
  \caption{Energy consumption. Electricity price: \(0.18\)\,€/kWh.}
  \label{tab:energy}
  \small
  \begin{tabular}{@{}lccr@{}}
    \toprule
    \textbf{Use} & \textbf{Hours} & \textbf{Power (kW)} & \textbf{Energy (kWh)} \\
    \midrule
    GPU training (CARLA + policy) & 350 & 0.45 & 157.5 \\
    Development and analysis      & 250 & 0.20 & 50.0 \\
    \midrule
    \textbf{Total} & & & \textbf{207.5} \\
    \bottomrule
  \end{tabular}
\end{table}

At \(0.18\)\,€/kWh, the \(207.5\,\mathrm{kWh}\) consumed amount to \(37.35\)~€. While small in
absolute terms, the figure is reported for completeness and because it scales linearly with the
number of training runs, which is the project's main computational cost driver.

\section{Travel and accommodation}\label{sec:budget-travel}

The project was developed locally and supervised through a combination of on-campus and online
meetings, so no travel or accommodation was required. This line is therefore \(0.00\)~€.

\section{Total budget}\label{sec:budget-total}

Table~\ref{tab:budget-summary} aggregates the previous sections. Indirect costs are applied at
\(15\%\) of direct costs, the contingency reserve at \(5\%\) of the direct-plus-indirect subtotal,
and IVA at \(21\%\) of the base budget. The resulting total cost of the project is
\(\mathbf{30\,139.73}\)~€.

\begin{table}[H]
  \centering
  \caption{Total project budget.}
  \label{tab:budget-summary}
  \small
  \begin{tabular}{@{}lr@{}}
    \toprule
    \textbf{Concept} & \textbf{Cost (€)} \\
    \midrule
    Labour, project team (Table~\ref{tab:labor-internal})  & 17\,445.00 \\
    Labour, support team (Table~\ref{tab:labor-support})   & 2\,650.00 \\
    Equipment, amortised (Table~\ref{tab:equipment})          & 336.11 \\
    Software and AI services (Table~\ref{tab:software})       & 160.00 \\
    Energy (Table~\ref{tab:energy})                           & 37.35 \\
    Travel and accommodation                                  & 0.00 \\
    \midrule
    \textbf{Direct costs}                                     & \textbf{20\,628.46} \\
    Indirect costs / overhead (15\%)                          & 3\,094.27 \\
    Contingency reserve (5\%)                                 & 1\,186.14 \\
    \midrule
    \textbf{Base budget (before tax)}                         & \textbf{24\,908.87} \\
    VAT (IVA, 21\%)                                            & 5\,230.86 \\
    \midrule
    \textbf{Total project budget}                             & \textbf{30\,139.73} \\
    \bottomrule
  \end{tabular}
\end{table}

Two observations close the chapter. First, labour accounts for \(20\,095\)~€ of the
\(20\,628\)~€ of direct cost (roughly \(97\%\)), which confirms that, like most
research-and-development efforts, the project's value is overwhelmingly in engineering time rather
than in hardware or licences. Second, the open-source software stack contributes a licensing cost
of exactly zero. Had the framework been built on commercial simulation or RL tooling, this line
alone could have rivalled the hardware cost. The low marginal cost of an additional experiment
(essentially only electricity, \(\approx 0.18\)\,€/kWh) is precisely what makes a reproducible,
open Safe-RL framework such as CALA-SafeRL economically attractive as a research platform.

\clearpage

% ============================================================
%   Chapter 2  --  State of the Art and Theoretical Framework
% ============================================================

\chapter{State of the Art and Theoretical Framework}\label{ch:sota}

This chapter reviews the concepts that underpin the rest of the document. It first formalises the
Reinforcement Learning problem (Section~\ref{sec:sota-rl}), then introduces Proximal Policy
Optimization and Generalised Advantage Estimation
(Sections~\ref{sec:sota-ppo}--\ref{sec:sota-gae}), surveys the Safe RL literature with a focus on
shielding (Sections~\ref{sec:sota-saferl}--\ref{sec:sota-shielding}), introduces the kinematic
bicycle model used by the adaptive shield (Section~\ref{sec:sota-bicycle}), and finishes with an
overview of the CARLA simulator and related work
(Sections~\ref{sec:sota-carla}--\ref{sec:sota-related}).

\section{Reinforcement Learning fundamentals}\label{sec:sota-rl}

The standard formulation of Reinforcement Learning relies on Markov Decision Processes
(MDPs)~\cite{sutton2018rl,puterman1994mdp}. A (discrete-time, discounted) MDP is a tuple
\(\mathcal{M} = (\mathcal{S}, \mathcal{A}, P, r, \gamma, \rho_0)\) where \(\mathcal{S}\) is the
state space, \(\mathcal{A}\) is the action space, \(P(s'\mid s,a)\) is the transition kernel,
\(r(s,a)\) is the reward function, \(\gamma \in [0,1)\) is the discount factor and \(\rho_0\)
is the initial-state distribution.

A stochastic policy \(\pi_\theta: \mathcal{S}\to\Delta(\mathcal{A})\), parameterised by
\(\theta\), induces a distribution over trajectories \(\tau = (s_0, a_0, s_1, a_1, \ldots)\).
The \emph{return} is the discounted cumulative reward:
\begin{equation}
  G(\tau) \;=\; \sum_{t=0}^{\infty} \gamma^t\, r(s_t, a_t).
  \label{eq:return}
\end{equation}
The training objective is to maximise the expected return:
\begin{equation}
  J(\theta) \;=\; \mathbb{E}_{\tau \sim \pi_\theta}\!\left[ G(\tau) \right].
  \label{eq:objective}
\end{equation}

Two central quantities guide optimisation. The \emph{state-value function} \(V^{\pi_\theta}(s)\)
measures the expected discounted return starting from state \(s\) under policy \(\pi_\theta\).
The \emph{action-value function} \(Q^{\pi_\theta}(s,a)\) additionally conditions on taking
action \(a\). The \emph{advantage} \(A^{\pi_\theta}(s,a) = Q^{\pi_\theta}(s,a) - V^{\pi_\theta}(s)\)
measures how much better than average a given action is at state \(s\).

\subsection{Policy gradient theorem}\label{sec:sota-pg}

Policy gradient methods directly optimise Equation~\eqref{eq:objective} by gradient ascent on
\(\theta\). The policy gradient theorem~\cite{sutton1999policygradient} provides a tractable
expression for the gradient:
\begin{equation}
  \nabla_\theta J(\theta) \;=\; \mathbb{E}_{\tau \sim \pi_\theta}\!\left[
     \sum_{t=0}^{\infty} \nabla_\theta \log \pi_\theta(a_t \mid s_t)\; A^{\pi_\theta}(s_t, a_t)
  \right].
  \label{eq:pg}
\end{equation}
In practice the advantage is replaced by a sample-based estimator \(\hat{A}_t\). Modern on-policy
methods such as A2C~\cite{mnih2016asynchronous}, TRPO~\cite{schulman2015trpo} and PPO~\cite{schulman2017ppo} all build on
Equation~\eqref{eq:pg}, differing mainly in how they constrain the update step size.

\section{Proximal Policy Optimization}\label{sec:sota-ppo}

Vanilla policy gradient is sensitive to the step size: a single large update can move the policy
so far from the data distribution that subsequent rollouts are unreliable. TRPO~\cite{schulman2015trpo}
resolves this by constraining the Kullback-Leibler divergence between old and new policy at each
update, but doing so requires second-order information. PPO~\cite{schulman2017ppo} achieves a
similar trust-region effect with a simple first-order surrogate.

Let \(r_t(\theta) = \pi_\theta(a_t \mid s_t)\,/\,\pi_{\theta_{\text{old}}}(a_t \mid s_t)\)
be the importance-sampling ratio. The \emph{clipped surrogate objective} is
\begin{equation}
  \mathcal{L}^{\mathrm{CLIP}}(\theta)
   \;=\; \mathbb{E}_t\!\left[
     \min\!\Big(
        r_t(\theta)\,\hat{A}_t,\;
        \mathrm{clip}\!\big(r_t(\theta),\, 1-\epsilon,\, 1+\epsilon\big)\,\hat{A}_t
     \Big)
   \right],
  \label{eq:ppo-clip}
\end{equation}
where \(\epsilon = 0.2\). As Figure~\ref{fig:ppo-clip} illustrates, the clipping prevents the
ratio \(r_t(\theta)\) from drifting far from one: for a positive advantage the policy cannot
over-exploit a good action beyond \(1+\epsilon\), and for a negative advantage it cannot
over-penalise beyond \(1-\epsilon\). This keeps updates within an approximate trust region
without ever computing a Hessian.

The full PPO loss also includes a squared-error value-function regression term weighted by
\(c_V = 0.25\) and an entropy bonus \(c_H \cdot \mathcal{H}[\pi_\theta]\) to encourage
exploration.

\begin{figure}[H]
  \centering
  \includegraphics[width=0.85\textwidth]{imgs/ppo_surrogate}
  \caption{PPO clipped surrogate objective. For a positive advantage (solid) the ratio is
           capped at \(1+\epsilon=1.2\), preventing over-exploitation of a good action. For a
           negative advantage (dashed) it is capped at \(1-\epsilon=0.8\). Both caps limit how
           far the new policy can deviate from the data-collecting policy.}
  \label{fig:ppo-clip}
\end{figure}

\subsection{Stabilisers and heuristics}\label{sec:sota-ppo-stabilisers}

Empirical studies~\cite{andrychowicz2020matters,engstrom2020implementation} have shown that
several implementation details have a first-order effect on PPO's sample efficiency and
stability. The ones relevant to this project are:

\begin{description}
  \item[\textbf{Value clipping.}] The value-function update is clipped so the critic cannot make
        overly aggressive updates within a single optimisation batch. The clipping magnitude
        \(\xi\) mirrors the PPO ratio clipping.
  \item[\textbf{KL early stopping.}] An unbiased estimator~\cite{schulman2020approxkl} of the
        KL divergence between old and new policy is computed after every minibatch. If the
        rolling average exceeds a target \(D^\star_{\mathrm{KL}} = 0.08\), the remaining epochs
        of the current update are skipped to prevent destructive updates.
  \item[\textbf{Observation normalisation.}] The 19-dimensional lane/marking/vehicle/route
        block is rescaled online to zero mean and unit variance using a running estimator. The
        rationale for normalising only the vector block (not the 720 LIDAR dimensions) is
        discussed in Section~\ref{sec:design-obsnorm}.
  \item[\textbf{Cosine learning-rate annealing.}] The learning rate decays from its initial
        value \(\eta_{\max}\) to a minimum \(\eta_{\min}\) following a cosine curve over the
        expected total number of PPO updates.
\end{description}

\section{Generalised Advantage Estimation}\label{sec:sota-gae}

The advantage \(A^{\pi_\theta}(s_t, a_t)\) in Equation~\eqref{eq:pg} must be estimated from
sampled rollouts. A plain Monte Carlo estimate has zero bias but high variance, whereas a one-step
TD residual \(\delta_t^V = r_t + \gamma V_\theta(s_{t+1}) - V_\theta(s_t)\) has low variance
but high bias. Generalised Advantage Estimation (GAE)~\cite{schulman2016gae} interpolates between
the two via an exponentially-weighted sum of TD residuals:
\begin{equation}
  \hat{A}_t^{\mathrm{GAE}(\gamma,\lambda)}
    \;=\; \sum_{\ell=0}^{\infty} (\gamma\lambda)^\ell\, \delta_{t+\ell}^V,
  \label{eq:gae}
\end{equation}
where \(\lambda \in [0,1]\) is the GAE parameter and \(\delta_{t+\ell}^V\) is the one-step TD
residual evaluated \(\ell\) steps ahead. Figure~\ref{fig:gae-lambda} shows how \(\lambda\) controls
the weight \((\gamma\lambda)^\ell\) placed on each residual, and therefore the bias-variance
trade-off: \(\lambda = 0\) keeps only the first residual (the one-step TD advantage: low variance,
biased), \(\lambda = 1\) weights all residuals by \(\gamma^\ell\) alone (the Monte Carlo advantage:
unbiased, high variance), and the value \(\lambda = 0.95\) used in this work sits between the two.
When an episode terminates due to a collision or going off-road, the critic bootstrap value is set to zero. When an episode is truncated by a step-count timeout, the critic value at the final observed state is
used to approximate the missing return tail.

\begin{figure}[H]
  \centering
  \includegraphics[width=0.85\textwidth]{imgs/gae_lambda}
  \caption{Generalised Advantage Estimation weights the one-step TD residual
           \(\delta_{t+\ell}^V\) by \((\gamma\lambda)^\ell\) (here \(\gamma = 0.99\)). The
           parameter \(\lambda\) interpolates between the low-variance, biased one-step TD estimate
           (\(\lambda = 0\), only \(\ell = 0\) contributes) and the unbiased, high-variance Monte
           Carlo estimate (\(\lambda = 1\)). The choice \(\lambda = 0.95\) gives a practical
           balance.}
  \label{fig:gae-lambda}
\end{figure}

\section{Safe Reinforcement Learning}\label{sec:sota-saferl}

The comprehensive survey by Garc\'ia and Fern\'andez~\cite{garcia2015comprehensive} organises
Safe RL into two broad categories: \emph{risk-sensitive} approaches that modify the optimality
criterion, and \emph{exploration-modifying} approaches that constrain the set of actions available
during learning. Four families are particularly relevant to autonomous driving:

\begin{itemize}
  \item \textbf{Constrained MDPs and CPO.} Safety is expressed as one or more auxiliary cost
        functions, and the policy is trained to maximise reward subject to cost-rate constraints.
        Achiam et al.~\cite{achiam2017cpo} introduced Constrained Policy Optimization (CPO), a
        TRPO-style algorithm that provides approximately monotonic constraint satisfaction.
  \item \textbf{Lagrangian and primal-dual methods.} A more scalable approach forms the dual
        problem with an adaptive multiplier for each constraint and solves it via alternating
        updates on the policy and the multiplier~\cite{stooke2020responsive}. Variants underpin
        much of modern Safe RL.
  \item \textbf{Safety-aware reward shaping.} Penalty terms in the reward function discourage
        unsafe behaviour. This is lightweight but provides no hard guarantees: the agent can
        still attempt unsafe actions during exploration. The reward shaper of
        Section~\ref{sec:design-reward} belongs to this family.
  \item \textbf{Runtime shielding.} A verifier inspects each proposed action and substitutes a
        safe fall-back if a safety specification would be violated. Shielding is orthogonal to
        the learning algorithm and can be composed with any RL backbone. Alshiekh et
        al.~\cite{alshiekh2018shielding} formalise the approach for finite systems with
        temporal-logic specifications, and Jansen et al.~\cite{jansen2020shielded} extend it to
        probabilistic settings. The shields in this project reason about predicted trajectories
        rather than formal logic.
\end{itemize}

\section{Safety shields}\label{sec:sota-shielding}

A runtime safety shield can be formalised as a function
\begin{equation}
  \Sigma: \mathcal{S}\times\mathcal{A} \to \mathcal{A},
  \qquad
  a^{\star}_t \;=\; \Sigma(s_t,\, a_t),
  \label{eq:shield-fn}
\end{equation}
that maps the current state and the agent's proposed action to a safe executed action. Two design
principles recur throughout the literature:

\begin{description}
  \item[\textbf{Minimal interference.}] The shield should override the agent's action only when
        strictly necessary, and the override should be as close as possible to the proposed
        action. In practice this is implemented by projecting onto a set of certified-safe
        actions, or by searching through a priority-ordered list of candidates.
  \item[\textbf{Credit assignment through the executed action.}] When the shield intervenes,
        the environment sees \(a^\star_t\), not the policy's \(a_t\). Storing the proposed
        action in the replay buffer would push the policy towards an action that was never
        executed. The framework addresses this by always storing \(a^\star_t\) and
        re-evaluating its log probability under the current policy
        (Section~\ref{sec:impl-shield-credit}).
\end{description}

Both shields in this project are \emph{post-posed}: they are inserted between the policy and the
simulator as a \texttt{gymnasium.Wrapper} and inspect the output of the policy after it has been
sampled.

\section{Kinematic bicycle model}\label{sec:sota-bicycle}

The adaptive shield must predict the vehicle's future position to verify whether a candidate
action will keep it within the lane. It uses a \emph{kinematic bicycle model}~\cite{rajamani2012vehicle,kong2015kinematic}, a well-established approximation that collapses each axle
into a single virtual wheel and tracks the rear axle position \((x,y)\), heading \(\theta\) and
longitudinal speed \(v\) (Figure~\ref{fig:bicycle-model}). The front-wheel steering angle
\(\delta\) and a longitudinal acceleration command drive the state forward in time steps of
\(\Delta t = 0.05\,\mathrm{s}\).

\begin{figure}[H]
  \centering
  \includegraphics[width=0.65\textwidth]{imgs/bicycle_model}
  \caption{Kinematic bicycle model geometry used by the adaptive shield. The rear axle is at
           \((x,y)\) with heading \(\theta\). Control inputs are the front-wheel steering
           \(\delta\) and a longitudinal acceleration. The model advances one step
           \(\Delta t = 0.05\,\mathrm{s}\) using a closed-form arc integration to avoid
           Euler-step drift at large steering angles.}
  \label{fig:bicycle-model}
\end{figure}

For small \(|\delta|\) the state advances with a simple forward-Euler step. For non-negligible
\(|\delta|\), the exact arc geometry is used to avoid integration drift. The discrete equations
are implemented in \texttt{src/Adaptative\_Shield/BicycleModel.py}. Kong et
al.~\cite{kong2015kinematic} show that, at the speeds and curvatures typical of urban driving,
the kinematic model closely approximates the full dynamic bicycle model for planning horizons
below one second, which covers the adaptive shield's longest horizon of ten steps
(\(0.5\,\mathrm{s}\)). The vehicle parameters for CARLA's Tesla Model~3 are in
Table~\ref{tab:bicycle-params} (Chapter~\ref{ch:design}).

\section{CARLA simulator}\label{sec:sota-carla}

CARLA~\cite{dosovitskiy2017carla} is an open-source simulator for autonomous driving research
built on Unreal Engine 4. Its features relevant to this thesis are summarised below.

\begin{itemize}
  \item \textbf{Client\textendash server architecture.} The simulator runs as a server process
        (\texttt{CarlaUE4.exe}) that holds the world state and renders the scene, while the
        Python client connects to it over a TCP port. This decoupling makes it straightforward
        to run training and the simulator on different machines.
  \item \textbf{Synchronous mode.} By default CARLA runs asynchronously. When
        \texttt{synchronous=True} the server waits for a \texttt{world.tick()} call from the
        client before advancing one fixed time step (\(\Delta t = 0.05\,\mathrm{s}\), i.e.\
        \(20\,\mathrm{Hz}\)). This is a prerequisite for reproducible experiments.
  \item \textbf{Waypoint API.} The road network is represented as an OpenDRIVE map from which
        CARLA exposes a lattice of \emph{waypoints} along lane centres. Each waypoint carries
        lane width, lane type, road curvature, lane-marking type and the allowed lane-change
        directions. The Waypoint API is the primary source of geometric information for both the
        reward shaper and the safety shields in this project.
  \item \textbf{Sensors.} CARLA provides a library of virtual sensors. The framework uses three:
        a semantic LIDAR, a collision detector and a lane-invasion detector. The semantic LIDAR
        tags every ray hit with the \texttt{carla.CityObjectLabel} of the struck primitive,
        enabling the three-channel decomposition of Section~\ref{sec:design-obs-lidar}.
  \item \textbf{Traffic Manager.} Non-playable vehicles (NPCs) are governed by CARLA's internal
        Traffic Manager, which follows the simulator's own traffic rules. The framework uses it
        to instantiate the NPC density required by each curriculum stage.
  \item \textbf{Predefined towns.} Town04 is a highway loop with multi-lane segments and
        entrance ramps, and it is the main benchmarking map used in
        Chapter~\ref{ch:experimentation}. Towns 01--03 cover urban scenarios, and Town05 offers
        maximum difficulty.
\end{itemize}

\section{Related work}\label{sec:sota-related}

\paragraph{Deep RL for autonomous driving.}
The surveys by Kiran et al.~\cite{kiran2022drlsurvey} and Aradi~\cite{aradi2022survey} offer
broad panoramas of the field. PPO is the most common backbone for continuous-control experiments
on CARLA, with notable examples including Chen et al.~\cite{chen2019carlarl} on urban driving and
CARLA Leaderboard submissions combining PPO with learned perception~\cite{chitta2023transfuser}.

\paragraph{Safety shields in simulated driving.}
Shielding has been applied to benchmarks such as HighwayEnv~\cite{leurent2018highway} and CommonRoad~\cite{althoff2017commonroad}. Krasowski et
al.~\cite{krasowski2020safe} use reachability analysis to certify lane-change manoeuvres,
Wabersich and Zeilinger~\cite{wabersich2021filter} combine learning-based control with a
model-predictive safety filter, and MetaDrive~\cite{li2023metadrive} provides a lighter-weight
shielding interface that inspired the early prototypes of this project. The present work
contributes a \emph{pair} of interchangeable shields in the CARLA ecosystem.

\paragraph{Curriculum learning.}
The general idea of presenting tasks in order of increasing difficulty dates back to Bengio et
al.~\cite{bengio2009curriculum}. Applications to RL range from procedural domain
randomisation~\cite{akkaya2019solving} to self-play curricula. The curriculum manager of
Section~\ref{sec:design-curriculum} includes a \emph{rollback} mechanism driven by empirical
instability on the 0--20~NPC progression, which is uncommon in the curriculum-learning
literature.

\paragraph{Bicycle model for trajectory prediction.}
Kong et al.~\cite{kong2015kinematic} show that, at speeds and curvatures typical of urban
driving, the kinematic model closely approximates the dynamic bicycle model for sub-second
horizons. This justifies its use inside the adaptive shield
(Section~\ref{sec:sota-bicycle}).

\clearpage

% ============================================================
%   Chapter 3  \textemdash\  Requirements Analysis
% ============================================================

\chapter{Requirements Analysis}\label{ch:requirements}

This chapter captures the requirements that guided the design of the framework. It distinguishes
functional requirements (what the system has to \emph{do}, Section~\ref{sec:req-functional}) from
non-functional requirements (how \emph{well} it has to do it, Section~\ref{sec:req-nonfunctional}).
Section~\ref{sec:req-constraints} lists the technical constraints imposed by the simulator, the
hardware and the academic timeline of the project. Section~\ref{sec:req-acceptance} converts
these requirements into measurable acceptance criteria that can be evaluated with the metrics
logged by the training pipeline.

\section{Functional requirements}\label{sec:req-functional}

Functional requirements describe the observable behaviour of the system. Each requirement is
given a short identifier of the form \texttt{FR-XX}, a descriptive title and a succinct
rationale. Requirements that depend on a specific implementation choice are justified in
Chapter~\ref{ch:design}.

\begin{table}[H]
  \centering
  \caption{Functional requirements of CALA-SafeRL.}
  \label{tab:functional-requirements}
  \small
  \begin{tabular}{@{}llp{7.8cm}@{}}
    \toprule
    \textbf{ID} & \textbf{Title} & \textbf{Description} \\
    \midrule
    FR-01 & Gymnasium environment    & The simulator shall be exposed through a
        \texttt{gymnasium.Env}-compatible wrapper supporting \texttt{reset}, \texttt{step},
        \texttt{render} and \texttt{close}. \\
    FR-02 & Continuous control       & The action space shall be a two-dimensional continuous
        box \([-1,1]^2\) corresponding to steering and joint throttle\textendash brake. \\
    FR-03 & Semantic LIDAR           & The observation shall include three LIDAR channels
        (combined, dynamic, static) of 240 rays each, normalised to \([0,1]\). \\
    FR-04 & Lane and route features  & The observation shall include lane features (lateral
        offset, heading error, edge distances, curvature, lane-change direction) and route
        features (waypoint angles, progress, speed limit) derived from the Waypoint API. \\
    FR-05 & PPO training loop        & A PPO agent with GAE, clipped surrogate, value loss,
        entropy bonus and KL early stopping shall be provided. \\
    FR-06 & Interchangeable shields  & The training and evaluation entry points shall accept a
        \texttt{--shield\_type} argument with values \texttt{none}, \texttt{basic} or
        \texttt{adaptive}, and switch shield implementations without code changes. \\
    FR-07 & Dense reward shaping     & A wrapper shall augment CARLA's sparse reward with dense
        components (speed, lane centring, heading, smoothness, progress, idle, shield
        interaction). \\
    FR-08 & Performance-driven curriculum & A curriculum manager shall progressively add NPC
        traffic as the agent's performance improves and shall roll back when performance
        collapses. \\
    FR-09 & Checkpoint persistence   & The agent shall save the best checkpoint on reward
        improvement and periodic checkpoints at a configurable frequency. \\
    FR-10 & Live metrics             & Every episode and every PPO update shall be logged to a
        SQLite database in WAL mode so that an external dashboard can read it while training is
        still in progress. \\
    FR-11 & Evaluation dashboard     & A headless evaluation mode and a Matplotlib dashboard
        mode shall be available, the latter showing LIDAR polar plot, speed gauge, lateral
        offset, shield status and intervention count. \\
    FR-12 & Lane-change awareness    & Both the reward shaper and the shield shall detect the
        condition \emph{lane change is permitted by road markings} and relax penalties that
        would otherwise block the manoeuvre. \\
    FR-13 & Reproducibility          & Training and evaluation shall be seeded independently
        (\(\mathrm{seed}=42\) and \(\mathrm{seed}=100\) respectively) and the seeds shall be
        exposed through command-line flags. \\
    \bottomrule
  \end{tabular}
\end{table}

\section{Non-functional requirements}\label{sec:req-nonfunctional}

Non-functional requirements describe qualities of the system: how fast, how portable, how
maintainable it must be. They are grouped below by ISO/IEC~25010 quality dimension.

\begin{table}[H]
  \centering
  \caption{Non-functional requirements of CALA-SafeRL.}
  \label{tab:nonfunctional-requirements}
  \small
  \begin{tabular}{@{}llp{7.8cm}@{}}
    \toprule
    \textbf{ID} & \textbf{Dimension} & \textbf{Description} \\
    \midrule
    NFR-01 & Performance             & The simulator shall run in synchronous mode at
        \(20\,\mathrm{Hz}\) (\(\Delta t = 0.05\,\mathrm{s}\)); all wrappers shall introduce an
        overhead small enough to sustain this rate on a single GPU workstation. \\
    NFR-02 & Determinism             & Given a fixed seed and a fixed hardware/software
        configuration, two training runs shall produce the same episode outcomes up to
        floating-point non-associativity. \\
    NFR-03 & Modularity              & The wrapper chain shall obey the Open/Closed principle:
        adding a new shield or a new reward component shall require no modifications to the
        existing layers. \\
    NFR-04 & Maintainability         & Python source files shall pass \texttt{ruff check} with
        the project's default configuration; public classes shall carry short docstrings. \\
    NFR-05 & Portability             & The framework shall run on Linux and Windows and shall
        automatically select the best PyTorch device among CUDA, Apple MPS and CPU. \\
    NFR-06 & Observability           & Every component that modifies actions or rewards shall
        expose an auxiliary \texttt{info} dictionary key (\texttt{executed\_action},
        \texttt{shield\_activated}, \texttt{risk\_level}, reward sub-components) so that
        training behaviour can be diagnosed post-hoc. \\
    NFR-07 & Extensibility           & The live-metrics schema shall admit new keys without
        breaking existing dashboards; both the SQLite logger and the Streamlit reader shall
        treat unknown metrics as first-class citizens. \\
    NFR-08 & Testability             & The reward shaper shall ship with deterministic unit
        tests that do not require a running CARLA server. \\
    NFR-09 & Documentation           & The project shall include a \texttt{CLAUDE.md} file that
        documents the wrapper chain, the observation space layout and the commands to reproduce
        training and evaluation. \\
    \bottomrule
  \end{tabular}
\end{table}

\section{Technical constraints}\label{sec:req-constraints}

The requirements above must be satisfied inside the technical envelope described in
Table~\ref{tab:technical-constraints}. These constraints are not design decisions but external
boundary conditions: the CARLA version is fixed by the research group's infrastructure, the
hardware is the author's workstation, and the timeline is set by the Bachelor's thesis calendar.

\begin{table}[H]
  \centering
  \caption{Technical constraints of the project.}
  \label{tab:technical-constraints}
  \small
  \begin{tabular}{@{}lp{10cm}@{}}
    \toprule
    \textbf{Category} & \textbf{Constraint} \\
    \midrule
    Simulator   & CARLA 0.9.16 running on Windows with the \texttt{-RenderOffScreen} flag
                  and fixed port \(2000\). \\
    Language    & Python \(\ge 3.10\). \\
    Core libs   & PyTorch (GPU when available), Gymnasium, NumPy, Matplotlib, Streamlit,
                  SQLite. Licences shall be compatible with open academic
                  distribution (MIT, BSD-3, Apache-2.0). \\
    Hardware    & Single desktop workstation; GPU with at least \(6\,\mathrm{GB}\) of VRAM is
                  required to run CARLA and the PyTorch policy concurrently. Training on CPU is
                  supported but orders of magnitude slower. \\
    Maps        & The main benchmark shall be Town04. Additional towns (Town01\textendash Town07,
                  Town10HD) may be used for qualitative evaluation but are not guaranteed to
                  converge with the default hyperparameters. \\
    Timeline    & The project shall be delivered within the Bachelor's thesis calendar of the
                  Universidad de Deusto. This precludes large hyperparameter sweeps that would
                  require multiple weeks of simulator time per configuration. \\
    Safety      & The framework shall only run inside simulation; deployment on a physical
                  vehicle is explicitly out of scope. \\
    \bottomrule
  \end{tabular}
\end{table}

\section{Acceptance criteria}\label{sec:req-acceptance}

Acceptance criteria translate the previous requirements into measurable conditions over the
metrics produced by \texttt{main\_train.py} and \texttt{main\_eval.py}. They are used in
Chapter~\ref{ch:experimentation} as the reference against which the three shielding
configurations are compared. The baseline corresponds to the unconstrained PPO agent
(\texttt{--shield\_type none}).

\begin{table}[H]
  \centering
  \caption{Acceptance criteria derived from the functional and non-functional requirements.
           Thresholds are placeholders to be replaced by the empirical numbers obtained in
           Chapter~\ref{ch:experimentation}.}
  \label{tab:acceptance-criteria}
  \small
  \begin{tabular}{@{}llp{7.5cm}@{}}
    \toprule
    \textbf{ID} & \textbf{Metric} & \textbf{Acceptance threshold} \\
    \midrule
    AC-01 & Reward convergence     & The 100-episode rolling average reward of the shielded
        agent shall be no worse than the baseline after an equal amount of training. \\
    AC-02 & Crash rate             & With the adaptive shield, the 100-episode rolling
        \texttt{crash\_rate} shall drop by at least a factor of \(\alpha=4\) with respect to the
        unshielded probe (empirically: \(16\%\to 2\%\), a factor of \(8\times\); the threshold
        is set conservatively at \(4\times\)). \\
    AC-03 & Off-road rate          & Under the adaptive shield, the 100-episode rolling
        \texttt{offroad\_rate} shall remain below \(\beta=0.15\) (\(15\%\)).
        Empirically: \(0.0\%\) at training end; \(12\%\) in the \(50\)-episode evaluation. \\
    AC-04 & Success rate           & On \(50\) evaluation episodes with seed \(100\), the agent
        shall reach the \texttt{success\_distance} target in at least \(\gamma=0.60\)
        (\(60\%\)) of them. Empirically: \(70\%\) with adaptive shield. \\
    AC-05 & Shield interference    & The fraction of all steps in which the adaptive shield
        activates shall not exceed \(\tau=0.40\) (\(40\%\)). The \emph{safe}, \emph{warning}
        and \emph{critical} risk-level fractions shall sum to \(1\). Empirically: shield
        activates in \(38.7\%\) of steps; \(65.8\%\) safe, \(25.6\%\) warning, \(8.6\%\)
        critical. \\
    AC-06 & Speed compliance       & At least \(90\%\) of steps with a valid speed limit shall
        satisfy \(\text{speed\_kmh} \le 1.05\cdot\text{speed\_limit\_kmh}\). \\
    AC-07 & Curriculum progression & The curriculum manager shall advance from stage 0 to at
        least stage 3 during a standard \(2\,500\)-episode run in all shielded
        configurations. \\
    AC-08 & Dashboard liveness     & The Streamlit dashboard shall refresh training metrics
        with an end-to-end latency below \(2\,\mathrm{s}\) from the moment the logger writes a
        row to the SQLite database. \\
    AC-09 & Reproducibility        & Two independent runs with the same seed and configuration
        shall produce identical reward curves in the first \(50\) episodes (before floating-point
        divergence becomes significant). \\
    \bottomrule
  \end{tabular}
\end{table}

The criteria with placeholder thresholds (\(\alpha, \beta, \gamma, \tau\)) are explicitly left to
be fixed once the experiments of Chapter~\ref{ch:experimentation} have been executed. This keeps
the requirements traceable while avoiding the pitfall of committing to numeric targets before the
empirical evidence is available.

\clearpage

% ============================================================
%   Chapter 4  \textemdash\  System Design
% ============================================================

\chapter{System Design}\label{ch:design}

This chapter describes how the concepts introduced in Chapter~\ref{ch:sota} are instantiated into
a concrete software architecture. Section~\ref{sec:design-architecture} presents the layered
wrapper chain, Sections~\ref{sec:design-obs}--\ref{sec:design-action} formalise the observation
and action spaces, Section~\ref{sec:design-agent} describes the PPO agent,
Section~\ref{sec:design-shields} the two safety shields,
Section~\ref{sec:design-reward} the reward shaper, Section~\ref{sec:design-curriculum} the
curriculum manager and Section~\ref{sec:design-metrics} the live metrics pipeline.

\section{Architectural overview}\label{sec:design-architecture}

The framework is organised as a vertical stack of components that communicate through the
standard \texttt{gymnasium.Env} interface. Each layer consumes the observation produced by the
layer below and (possibly) transforms both the action travelling downwards and the reward
travelling upwards. Figure~\ref{fig:wrapper-chain} shows the chain.

\begin{figure}[h]
  \centering
  \small
  \begin{tabular}{c}
    \hline
    \textbf{PPOAgent}                            \\[2pt]
    \(\downarrow\) proposed action $\; a_t$ \qquad \(\uparrow\) observation $\; s_t$, reward $\; r_t$ \\[4pt]
    \textbf{Safety Shield}\; (\texttt{none}\,/\,\texttt{basic}\,/\,\texttt{adaptive}) \\[2pt]
    \(\downarrow\) executed action $\; a_t^\star$ \qquad \(\uparrow\) observation, base reward \\[4pt]
    \textbf{CarlaRewardShaper}                  \\[2pt]
    \(\downarrow\) executed action $\; a_t^\star$ \qquad \(\uparrow\) shaped reward $\; \tilde{r}_t$ \\[4pt]
    \textbf{CarlaEnv}\; (\texttt{gymnasium.Env} wrapping the CARLA Python API) \\[2pt]
    \(\downarrow\) vehicle command \qquad \(\uparrow\) sensor readings \\[4pt]
    \textbf{CARLA Simulator}\; (\texttt{CarlaUE4.exe}) \\[2pt]
    \hline
  \end{tabular}
  \caption{Wrapper chain of the framework. Each layer is a \texttt{gymnasium.Wrapper} that can be
           composed, replaced or removed without touching its neighbours.}
  \label{fig:wrapper-chain}
\end{figure}

A subtlety of this layout deserves an explicit comment. The safety shield is placed \emph{above}
the reward shaper, not below it. This ordering has a functional reason: the shaper reads the
\texttt{info} keys \texttt{executed\_action} and \texttt{proposed\_action}\textemdash both
produced by the shield\textemdash to compute components such as the shield-intervention penalty,
the smoothness penalty during an override and the lane-change transition guard. If the shaper
were placed above the shield, those keys would not yet exist when the shaper inspects the
\texttt{info} dictionary, and the shield interaction could not be accounted for in the reward.

The environment runs in \emph{synchronous mode} with a fixed time step of \(\Delta t =
0.05\,\mathrm{s}\). Every call to \texttt{env.step(action)} sends the action to CARLA, requests a
single \texttt{world.tick()} and collects the resulting sensor readings. This guarantees
determinism given a fixed seed, at the cost of not being able to use CARLA's asynchronous
rendering speed-ups.

\section{Observation space}\label{sec:design-obs}

The observation vector \(s_t \in \mathbb{R}^{735}\) concatenates three semantic LIDAR channels and
a small block of geometric and route features. Appendix~\ref{app:observation} contains the full
index-by-index table; this section describes the design rationale.

\subsection{Semantic LIDAR}\label{sec:design-obs-lidar}

The framework uses a simulated LIDAR with \(N = 240\) rays uniformly distributed over
\(360^\circ\) (angular resolution \(1.5^\circ/\text{ray}\), range \(50\,\mathrm{m}\)). Each ray
is normalised to the \([0, 1]\) interval, where \(1\) denotes a free ray up to the maximum range
and values close to \(0\) denote an obstacle next to the vehicle. The semantic channel of each
ray is obtained from CARLA's \texttt{SemanticLidarSensor}, which tags every hit with the
\texttt{carla.CityObjectLabel} of the struck primitive.

Three derived channels are packed into the observation:

\begin{itemize}
  \item \(\text{obs}[0:240]\) \textemdash\ \textbf{combined channel}. Includes every obstacle
        (dynamic and static) plus road edges.
  \item \(\text{obs}[240:480]\) \textemdash\ \textbf{dynamic-only channel}. Keeps only rays that
        struck a vehicle or a pedestrian.
  \item \(\text{obs}[480:720]\) \textemdash\ \textbf{static-only channel}. Keeps only rays that
        struck a wall, a barrier, a pole or a similar static object.
\end{itemize}

The decomposition is motivated by an empirical observation: a single-channel LIDAR induces a
\emph{zig-zag} pattern when the vehicle drives along a guardrail, because the policy interprets
the low returns from the guardrail as a moving obstacle and keeps trying to evade it laterally.
Separating dynamic from static obstacles allows the policy to learn that static obstacles can be
treated as soft boundaries while dynamic obstacles require anticipation.

\subsection{Lane, vehicle and route features}\label{sec:design-obs-vector}

A \(15\)-dimensional vector block is appended to the LIDAR. Its components are:

\begin{itemize}
  \item \(\text{obs}[720:728]\) \textemdash\ \textbf{lane features (8 dims)}:
        \(\texttt{lateral\_offset\_norm}\),
        \(\texttt{heading\_error\_norm}\),
        \(\texttt{on\_edge\_warning}\),
        \(\texttt{lane\_width\_norm}\),
        \(\texttt{dist\_left\_edge\_norm}\),
        \(\texttt{dist\_right\_edge\_norm}\),
        \(\texttt{lane\_change\_permitted}\),
        \(\texttt{road\_curvature\_norm}\).
  \item \(\text{obs}[728:730]\) \textemdash\ \textbf{vehicle state (2 dims)}:
        \(\texttt{speed\_norm}\), \(\texttt{steering}\).
  \item \(\text{obs}[730:735]\) \textemdash\ \textbf{route features (5 dims)}:
        \(\texttt{next\_wp\_angle\_norm}\) (\(5\,\mathrm{m}\) ahead),
        \(\texttt{wp\_angle\_20m\_norm}\),
        \(\texttt{progress\_norm}\),
        \(\texttt{speed\_limit\_norm}\),
        \(\texttt{speed\_ratio}\).
\end{itemize}

All features are extracted from CARLA's Waypoint API and therefore carry centimetre-level
accuracy. The two offsets \(\texttt{dist\_left\_edge\_norm}\) and
\(\texttt{dist\_right\_edge\_norm}\) are normalised by the lane width, so that both take the
value \(0.5\) when the vehicle is at the centre of the lane and approach \(0\) as it nears the
corresponding edge.

\section{Action space}\label{sec:design-action}

The action space is the continuous box \(\mathcal{A} = [-1, 1]^2\). Its two components are
\begin{itemize}
  \item \(a^{(1)} \in [-1, 1]\) \textemdash\ \emph{steering}, mapped linearly to
        \(\delta \in [-\delta_{\max}, +\delta_{\max}]\) with \(\delta_{\max} = 0.6\,\mathrm{rad}\).
  \item \(a^{(2)} \in [-1, 1]\) \textemdash\ \emph{joint throttle\textendash brake}: positive
        values drive the throttle, negative values drive the brake, zero is coast. The
        acceleration saturation values are \(+3.0\,\mathrm{m/s^2}\) and \(-7.0\,\mathrm{m/s^2}\).
\end{itemize}

Merging throttle and brake into a single signed channel has two practical advantages: it reduces
the dimensionality of the policy output, and it prevents the agent from learning the
counter-productive behaviour of activating throttle and brake simultaneously.

Continuous actions are sampled from a diagonal Gaussian and squashed through a hyperbolic
tangent to keep them inside the box:
\begin{equation}
  u_t \sim \mathcal{N}\!\big(\mu_\theta(s_t),\, \sigma_\theta^2\big),
  \qquad
  a_t = \tanh(u_t).
  \label{eq:policy-squash}
\end{equation}
The log probability of the squashed sample requires the log-determinant of the
tanh Jacobian~\cite{haarnoja2018sac}:
\begin{equation}
  \log \pi_\theta(a_t \mid s_t)
    \;=\; \log \mathcal{N}\!\big(u_t;\, \mu_\theta(s_t),\, \sigma_\theta\big)
        \;-\; \sum_{i} \log\!\big(1 - a_{t,i}^2 + \varepsilon\big),
  \label{eq:logprob-tanh}
\end{equation}
where \(\varepsilon\) is a small constant added for numerical stability. The same correction is
used to evaluate the log probability of an executed action, which is a subtle but important
point for credit assignment under a shield (Section~\ref{sec:impl-shield-credit}).

\section{Agent design}\label{sec:design-agent}

\subsection{Actor\textendash Critic network}\label{sec:design-agent-net}

The actor and the critic share a representation of the observation but have separate heads.
Feeding the raw \(735\)-dimensional observation directly into a shallow MLP dilutes the signal:
the \(720\) LIDAR dimensions dominate the input while the \(15\) lane/vehicle/route features
carry most of the task-relevant information. The network therefore includes a dedicated
\emph{LIDAR encoder}:

\begin{equation}
  \mathrm{enc}(s_t)
    \;=\; \mathrm{Tanh}\!\Big( W_2 \, \mathrm{Tanh}\big(W_1\, s_t^{\mathrm{LIDAR}} + b_1\big) + b_2 \Big)
    \;\in\; \mathbb{R}^{64},
  \label{eq:lidar-encoder}
\end{equation}
with \(W_1 \in \mathbb{R}^{256\times 720}\) and \(W_2 \in \mathbb{R}^{64 \times 256}\). The
encoder output is concatenated with the \(15\)-dimensional vector block to produce a
\(79\)-dimensional feature \(\phi(s_t)\):

\begin{equation}
  \phi(s_t) \;=\; \big[\, \mathrm{enc}(s_t)\; ;\; s_t^{\mathrm{vector}} \,\big] \in \mathbb{R}^{79}.
  \label{eq:features}
\end{equation}

The critic is a two-layer MLP with \(256\) units per hidden layer and a scalar output:
\begin{equation}
  V_\theta(s_t) \;=\; W^{V}\, \mathrm{Tanh}\!\Big( \mathrm{MLP}_{\mathrm{critic}}\big(\phi(s_t)\big) \Big) + b^{V}.
  \label{eq:critic-head}
\end{equation}

The actor shares the same trunk layout but, instead of a scalar output, produces a
two-dimensional mean vector and has a separate log-standard-deviation parameter that does not
depend on the state:
\begin{equation}
  \mu_\theta(s_t) \;=\; W^{\mu}\, \mathrm{Tanh}\!\Big( \mathrm{MLP}_{\mathrm{actor}}\big(\phi(s_t)\big) \Big) + b^{\mu},
  \qquad
  \log \sigma_\theta \in \mathbb{R}^2.
  \label{eq:actor-head}
\end{equation}

The log-standard-deviation is clamped to the range \([-3.0,\, -0.2]\). Its initial value is set to
\(-0.7\), which corresponds to a pre-\(\tanh\) standard deviation of approximately \(0.50\). This
gives broad but non-saturating exploration at the beginning of training.

\subsection{PPO update}\label{sec:design-agent-ppo}

At each environment step the agent appends the tuple
\(\big(s_t,\, a^\star_t,\, \log \pi_\theta(a^\star_t\mid s_t),\, r_t,\, d_t,\, \mathrm{trunc}_t,\, V^{\text{boot}}_t\big)\)
to the rollout buffer. Two details are specific to this project:

\begin{enumerate}
  \item The buffer stores the \emph{executed} action \(a^\star_t\) (the action that actually
        reached the simulator), not the action sampled from the policy. The log probability
        stored in the buffer is recomputed by \texttt{agent.evaluate\_action} with
        Equation~\eqref{eq:logprob-tanh}. This is what ties the gradient computation to the
        action the environment saw, avoiding the credit-assignment error that would otherwise be
        introduced by the shield.
  \item When an episode is truncated by a timeout but not terminated (no crash, no off-road),
        the value \(V^{\text{boot}}_t = V_\theta(s_{t+1})\) is stored and used to bootstrap GAE
        at that step. When the episode terminates naturally, \(V^{\text{boot}}_t = 0\) is stored.
\end{enumerate}

When the buffer reaches \(M = 2048\) samples, GAE is run backwards through the rollout
(Equation~\eqref{eq:gae}), advantages are standardised to zero mean and unit variance, and the
return target is divided by the running standard deviation described in
Section~\ref{sec:design-returnnorm}. A loop of up to \(K = 10\) epochs of PPO is then performed,
with an early stop triggered by Equation~\eqref{eq:kl-schulman} once the moving average of the
approximate KL divergence exceeds \(D^{\star}_{\mathrm{KL}}\). The gradient is clipped to an
\(\ell_2\) norm of \(1.0\) and the learning rate follows the cosine schedule of
Equation~\eqref{eq:cosine-lr}.

\subsection{Observation normalisation}\label{sec:design-obsnorm}

Observation normalisation is performed online with Welford's incremental mean/variance
estimator~\cite{welford1962notes}. Given a running mean \(\hat{\mu}\), running variance
\(\hat{\sigma}^2\) and count \(N\), the update for a batch \(\mathbf{x}\) of size \(B\) is
\begin{align}
  \Delta \;&=\; \bar{\mathbf{x}} - \hat{\mu}, \\
  \hat{\mu} \;&\leftarrow\; \hat{\mu} + \Delta\, \frac{B}{N+B}, \\
  \hat{\sigma}^2 \;&\leftarrow\; \frac{N\hat{\sigma}^2 + B \, \mathrm{Var}(\mathbf{x}) + \Delta^2\, \frac{NB}{N+B}}{N+B}, \\
  N \;&\leftarrow\; N + B.
  \label{eq:welford}
\end{align}
Normalised observations are clipped to \([-5, 5]\) to contain outliers.

An implementation detail that turned out to be important is that normalisation is applied
\emph{only} to the \(15\)-dimensional vector block, not to the \(720\) LIDAR dimensions. The
LIDAR is already in \([0,1]\) by construction, so normalisation would not change its scale but
would artificially introduce outliers when the curriculum advances from stage 0 (no NPCs, the
dynamic channel is the constant \(1.0\)) to stage 1 (dynamic channel suddenly has non-unit
entries). Keeping the LIDAR un-normalised avoids this distribution shift and therefore avoids
the gradient spikes observed in early prototypes.

\subsection{Return normalisation}\label{sec:design-returnnorm}

The return targets passed to the critic loss in Equation~\eqref{eq:ppo-aux} can span several
orders of magnitude, which makes MSE training unstable. To counter this, the returns are divided
by the running standard deviation
\begin{equation}
  \hat{R}_t \;\leftarrow\; \frac{\hat{R}_t}{\hat{\sigma}_R + \varepsilon},
  \label{eq:return-norm}
\end{equation}
where \(\hat{\sigma}_R\) is computed by the same Welford estimator applied to the rollout returns.
Note that the mean is \emph{not} subtracted; subtracting a running mean of returns would bias the
advantages, which are already standardised separately.

\section{Safety shields}\label{sec:design-shields}

Both shields implement the function \(\Sigma\) of Equation~\eqref{eq:shield-fn}. They differ in
the semantic information they consume and in the strategy by which they generate the corrected
action.

\subsection{Basic shield}\label{sec:design-basic-shield}

The basic shield (\texttt{CarlaSafetyShield}) operates in two layers. The first layer analyses
the LIDAR scan for imminent obstacles; the second layer uses the Waypoint API to detect
unsafe lane-keeping conditions.

\paragraph{Obstacle analysis.}
Let \(\mathrm{scan}_t \in [0,1]^{240}\) be the combined LIDAR channel at time \(t\). The shield
partitions it into three angular sectors:
\begin{align}
  \mathrm{front}_t \;&=\; \mathrm{scan}_t\big[225^\circ\!:\!315^\circ\big] \;\cup\; \mathrm{scan}_t\big[0^\circ\!:\!22.5^\circ\big] \cup \mathrm{scan}_t\big[337.5^\circ\!:\!360^\circ\big], \label{eq:shield-front}\\
  \mathrm{side}^{\mathrm{R}}_t \;&=\; \mathrm{scan}_t\big[60^\circ\!:\!120^\circ\big], \label{eq:shield-rside}\\
  \mathrm{side}^{\mathrm{L}}_t \;&=\; \mathrm{scan}_t\big[240^\circ\!:\!300^\circ\big], \label{eq:shield-lside}
\end{align}
and declares the situation unsafe whenever any of the following conditions holds:
\begin{equation}
  \min \mathrm{front}_t < \tau_{\mathrm{f}}, \quad
  \min \mathrm{side}^{\mathrm{R}}_t < \tau_{\mathrm{s}}, \quad
  \min \mathrm{side}^{\mathrm{L}}_t < \tau_{\mathrm{s}}.
  \label{eq:shield-obstacle}
\end{equation}
The default thresholds are \(\tau_{\mathrm{f}} = 0.15\) and \(\tau_{\mathrm{s}} = 0.04\) in
normalised LIDAR units.

\paragraph{Lane-keeping analysis.}
The shield reads the lateral offset \(\ell_t = \texttt{lateral\_offset\_norm}\) and the heading
error \(h_t = \texttt{heading\_error\_norm}\) from the \texttt{info} dict and declares the
situation unsafe when
\begin{equation}
  |\ell_t| > \tau_\ell
  \quad \text{or} \quad
  |h_t| > \tau_h,
  \label{eq:shield-lane}
\end{equation}
with defaults \(\tau_\ell = 0.82\) (expressed as a fraction of the lane half-width) and
\(\tau_h = 0.60\).

\paragraph{Correction law.}
When the situation is unsafe the shield overrides the action with a proportional controller.
Let \(k_{\mathrm{s}}\) and \(k_{\mathrm{b}}\) be the steering and brake gains. The correction is
\begin{equation}
  \delta^\star_t \;=\; \mathrm{clip}\!\Big(
      -k_{\mathrm{s}}\, \ell_t \;-\; 0.45\, k_{\mathrm{s}}\, h_t \;+\; \mathrm{sidepush}_t,\;
      -1,\, 1\Big),
  \label{eq:shield-steer-correction}
\end{equation}
\begin{equation}
  (tb)^\star_t \;=\; \mathrm{clip}\!\Big(
      -\, k_{\mathrm{b}}\, \max\!\big(0,\; \tau_{\mathrm{f}} - \min\mathrm{front}_t\big) / \tau_{\mathrm{f}},\;
      -1,\, 0\Big),
  \label{eq:shield-brake-correction}
\end{equation}
where \(\mathrm{sidepush}_t\) contains the lateral correction triggered by close obstacles in
\(\mathrm{side}^{\mathrm{R}}\) or \(\mathrm{side}^{\mathrm{L}}\). Equations
\eqref{eq:shield-steer-correction}--\eqref{eq:shield-brake-correction} follow the
\emph{minimal-interference principle}: if \(\ell_t\), \(h_t\) and \(\min\mathrm{front}_t\) are
all far from their thresholds, all three contributions vanish and the shield produces the
zero-override action. Conversely, the closer the system is to an unsafe region, the stronger the
corrective action.

\subsection{Adaptive-horizon shield}\label{sec:design-adaptive-shield}

The adaptive-horizon shield (\texttt{CarlaAdaptiveHorizonShield}) extends the basic design with
three mechanisms: a risk classifier, a trajectory prediction with the bicycle model, and a
priority-ordered candidate search.

\paragraph{Risk classifier.}
Three risk levels are defined: \(\mathrm{safe}\), \(\mathrm{warning}\) and \(\mathrm{critical}\).
Each level fixes a verification horizon \(H\) (number of bicycle-model steps) and a threshold
multiplier \(\kappa\) applied to the basic thresholds. The configuration is
\[
  \mathrm{safe}:\; H=1,\; \kappa = 1.0;
  \qquad
  \mathrm{warning}:\; H=5,\; \kappa = 1.5;
  \qquad
  \mathrm{critical}:\; H=10,\; \kappa = 2.0.
\]
The level is chosen as the worst of two independent criteria. The first is the
\emph{frontal-distance} criterion: let \(d^{\mathrm{dyn}}_t = \min\text{front}^{\mathrm{dynamic}}_t\)
be the minimum of the dynamic-only LIDAR front sector. Then
\begin{equation}
  \mathrm{level}^{\mathrm{frontal}}_t \;=\;
     \begin{cases}
       \mathrm{safe}     & \text{if } d^{\mathrm{dyn}}_t > 0.50,\\
       \mathrm{warning}  & \text{if } 0.20 < d^{\mathrm{dyn}}_t \le 0.50,\\
       \mathrm{critical} & \text{if } d^{\mathrm{dyn}}_t \le 0.20.
     \end{cases}
  \label{eq:shield-risk-front}
\end{equation}
The second criterion, \emph{lateral position}, thresholds \(|\ell_t|\) at \(0.70\) and \(0.85\)
(warning and critical, respectively). If the Waypoint API reports that a lane change is
permitted at the current location, a \emph{critical} lateral level is downgraded to
\emph{warning} so that the shield does not block a legitimate lane-change manoeuvre.

\paragraph{Trajectory check.}
Given a candidate action \(a = (\delta, tb)\), the shield asks the bicycle model of
Section~\ref{sec:sota-bicycle} for the trajectory \(\{(x_k, y_k, \theta_k)\}_{k=0}^{H}\). Each
predicted position \((x_k, y_k)\) is then projected onto the road using
\texttt{carla.Map.get\_waypoint} restricted to driving lanes. The candidate is declared safe
when all of the following hold:

\begin{equation}
  \min \mathrm{front}^{\mathrm{combined}}_t \;\ge\; \tau_{\mathrm{f}} / \kappa,
  \qquad
  \min \mathrm{side}^{\mathrm{R/L,\,static}}_t \;\ge\; \tau_{\mathrm{s}} / \kappa,
  \label{eq:shield-lidar-check}
\end{equation}
\begin{equation}
  \mathrm{nearest\_pedestrian\_m}_t \;\ge\; 4.0\,\mathrm{m},
  \label{eq:shield-ped-check}
\end{equation}
\begin{equation}
  \forall\, k \in \{1,\ldots,H\}:\;
    |\ell_k| \;\le\; \max\!\big( \tau_\ell / \kappa,\; 0.45 \big),
  \label{eq:shield-lane-check}
\end{equation}
where \(\ell_k = \texttt{lat\_offset}(x_k, y_k) / (\texttt{lane\_width}/2)\) is the normalised
lateral offset of the predicted position with respect to the nearest driving lane. The floor of
\(0.45\) in Equation~\eqref{eq:shield-lane-check} prevents the lateral threshold from becoming
so tight in the \emph{critical} level that no correction action can satisfy it. When the
\emph{lane-change permitted} flag is set, the lateral threshold is further raised to \(1.2\) so
that the transition band between lanes is admissible.

\paragraph{Priority-ordered candidate search.}
If the agent's action \(a_t\) does not pass the trajectory check, the shield inspects an
ordered list of candidate actions and returns the first one that does. The ordering depends on
the dominant threat, which is identified in three mutually exclusive buckets:

\begin{itemize}
  \item \textbf{Pedestrian emergency}\; (\(\mathrm{nearest\_pedestrian\_m}_t < 4\,\mathrm{m}\)).
        The override is always \((\delta, tb) = (0, -1)\) (maximum brake, no search).
  \item \textbf{Dynamic threat}\; (nearby vehicle in front, predicted time-to-contact below
        \(\max(5\,\mathrm{s},\, 1.5\,v_t)\) metres). The first candidates privilege braking;
        steering is adjusted afterwards.
  \item \textbf{Static threat}\; (static obstacle in the side sector). The first candidates
        privilege a steering correction back to the lane centre, accompanied by a mild brake.
  \item \textbf{Lateral-only threat}\; (\(|\ell_t| > 0.65\) without any dynamic or static
        threat). The first candidates apply a lane-centring steering with a small throttle to
        keep enough speed for directional control.
  \item \textbf{Balanced fallback}\; (default): a symmetric list of brake/steer combinations.
\end{itemize}

If none of the candidates pass the trajectory check, the shield returns the \emph{fallback}
action
\begin{equation}
  a_t^{\mathrm{fallback}} \;=\; \big(\mathrm{clip}(-1.2\,\ell_t,\, -0.8,\, 0.8),\; -0.4\big).
  \label{eq:shield-fallback}
\end{equation}

\subsection{Bicycle-model parameters}\label{sec:design-bicycle-params}

The bicycle model described in Section~\ref{sec:sota-bicycle} is instantiated with the
parameters of CARLA's Tesla Model 3, summarised in Table~\ref{tab:bicycle-params}.

\begin{table}[h]
  \centering
  \caption{Parameters of the kinematic bicycle model used by the adaptive shield.}
  \label{tab:bicycle-params}
  \small
  \begin{tabular}{@{}lll@{}}
    \toprule
    \textbf{Symbol} & \textbf{Value} & \textbf{Meaning} \\
    \midrule
    \(L\)               & \(2.87\,\mathrm{m}\)          & Wheelbase of the Tesla Model 3 \\
    \(\delta_{\max}\)   & \(0.60\,\mathrm{rad}\) (\({\approx}34^\circ\)) & Maximum steering angle \\
    \(\Delta t\)        & \(0.05\,\mathrm{s}\)          & Time step, matched to the simulator \\
    \(a_{\max}\)        & \(3.0\,\mathrm{m/s^2}\)       & Maximum forward acceleration \\
    \(-a_{\min}\)       & \(7.0\,\mathrm{m/s^2}\)       & Maximum deceleration (brake) \\
    \bottomrule
  \end{tabular}
\end{table}

\section{Reward shaping}\label{sec:design-reward}

The reward shaper \texttt{CarlaRewardShaper} produces a shaped reward \(\tilde{r}_t\) by
augmenting the sparse base reward \(r^{\mathrm{base}}_t\) of \texttt{CarlaEnv} (success
\(+30\), collision \(-10\), off-road \(-20\)) with twelve dense components. The shaped reward is
\begin{equation}
  \tilde{r}_t \;=\; r^{\mathrm{base}}_t + r^{\mathrm{alive}}_t + r^{\mathrm{speed}}_t + r^{\mathrm{lane}}_t + r^{\mathrm{head}}_t + r^{\mathrm{prog}}_t + r^{\mathrm{no-shield}}_t
             - p^{\mathrm{smooth}}_t - p^{\mathrm{inv}}_t - p^{\mathrm{edge}}_t - p^{\mathrm{wrong-h}}_t - p^{\mathrm{shield}}_t - p^{\mathrm{idle}}_t - p^{\mathrm{drift}}_t.
  \label{eq:shaped-reward}
\end{equation}

Each component is described below. Unless stated otherwise, all variables read from the
\texttt{info} dictionary come from the Waypoint API.

\paragraph{Speed gate.}
A common factor \(g_t \in [0,1]\) modulates the components that only make sense when the vehicle
is actually moving. It is a linear ramp of the speed \(v_t\) between the two
thresholds \(v_{\min}\) and \(v_g\):
\begin{equation}
  g_t \;=\; \mathrm{clip}\!\left( \frac{v_t - v_{\min}}{v_g - v_{\min}},\, 0,\, 1 \right).
  \label{eq:speed-gate}
\end{equation}
With defaults \(v_{\min} = 5\,\mathrm{km/h}\) and \(v_g = 10\,\mathrm{km/h}\).

\paragraph{Alive bonus.}
\begin{equation}
  r^{\mathrm{alive}}_t \;=\; c_{\mathrm{alive}}\, g_t \;\cdot\; \mathbb{1}[\,\text{on\_road}_t\,].
  \label{eq:r-alive}
\end{equation}
Gating by \(g_t\) is essential to prevent the local optimum \emph{``stop and collect
0.15/step''} that plagued earlier prototypes of the shaper.

\paragraph{Speed bonus.}
Let \(v^{\mathrm{lim}}_t\) be the effective speed limit (from CARLA's speed limits or the
project-level target, whichever is defined). A curvature correction reduces the target speed in
sharp bends:
\begin{equation}
  v^{\mathrm{target}}_t \;=\; v^{\mathrm{lim}}_t \cdot \max\!\Big(\,1 - c_{\mathrm{curv}}\cdot \min\!\big(|\kappa_t|/0.6,\, 1\big),\; 0.4\,\Big).
  \label{eq:target-speed}
\end{equation}
The speed bonus is a Gaussian centred on this curvature-adjusted target, with a speed-ceiling
penalty when the vehicle exceeds the raw limit by more than the margin:
\begin{equation}
  r^{\mathrm{speed}}_t \;=\; w_{\mathrm{v}}\, \exp\!\left( -\frac{(v_t - v^{\mathrm{target}}_t)^2}{2\sigma^2} \right)
     \;-\; w_{\mathrm{v}} \cdot 0.8\, \max\!\Big(\,0,\, \frac{v_t - (1+\eta)\,v^{\mathrm{lim}}_t}{v^{\mathrm{lim}}_t}\,\Big),
  \label{eq:r-speed}
\end{equation}
with \(\sigma = 0.35\, v^{\mathrm{target}}_t\) and \(\eta = 0.05\) the tolerance margin.

\paragraph{Lane centring and heading.}
\begin{equation}
  r^{\mathrm{lane}}_t \;=\; w_{\mathrm{lc}}\, \max(g_t,\, 0.3)\, \mathrm{clip}\!\big(\min(d^{\mathrm{L}}_t, d^{\mathrm{R}}_t)/0.5,\, 0,\, 1\big),
  \label{eq:r-lane}
\end{equation}
\begin{equation}
  r^{\mathrm{head}}_t \;=\; w_{\mathrm{h}}\, g_t \, \exp\!\left( -\frac{h_t^2}{2\cdot 0.40^2} \right),
  \label{eq:r-head}
\end{equation}
where \(d^{\mathrm{L}}_t,\, d^{\mathrm{R}}_t\) are the normalised distances to the left and right
lane edges. The floor of \(0.3\) in Equation~\eqref{eq:r-lane} provides a small centring signal
even when the vehicle is stopped, which helps the policy learn to align before accelerating.

\paragraph{Smoothness penalty.}
\begin{equation}
  p^{\mathrm{smooth}}_t \;=\; w_{\mathrm{sm}}\, \big| \delta_t - \delta_{t-1} \big|.
  \label{eq:p-smooth}
\end{equation}

\paragraph{Invasion, edge and wrong-heading penalties.}
The lane-invasion penalty uses the signal of CARLA's lane-invasion sensor; it is weighted by the
invasion severity (a function of \(|\ell_t|\)) and is suppressed when the invasion is
intentional (large steering magnitude) or when a lane change is permitted. The edge penalty is a
quadratic ramp triggered when the closest edge distance drops below \(0.4\); the wrong-heading
penalty triggers when \(|\text{heading\_error\_deg}_t|\) exceeds \(90^\circ\).

\paragraph{Progress bonus.}
A fixed bonus \(w_{\mathrm{prog}}\) is awarded every time the vehicle crosses a distance
milestone of \(\Delta_{\mathrm{prog}} = 25\,\mathrm{m}\). This avoids the pathological behaviour
of a speed bonus that rewards high speed regardless of whether the vehicle is actually moving
along the route.

\paragraph{Shield interaction.}
Let \(a_t\) be the action proposed by the agent and \(a^\star_t\) the action executed by the
environment. Define the shield-divergence measure
\begin{equation}
  \Delta^{\mathrm{shield}}_t \;=\; \frac{\|a^\star_t - a_t\|_2}{2\sqrt{2}}.
  \label{eq:shield-div}
\end{equation}
The normalisation by \(2\sqrt{2}\) maps the divergence to the unit interval (the maximum
distance between two points of \([-1,1]^2\) is \(2\sqrt{2}\)). The shield-interaction components
are
\begin{equation}
  p^{\mathrm{shield}}_t \;=\; \Delta^{\mathrm{shield}}_t\, w_{\mathrm{sh}}\;
    \mathbb{1}[\,\text{shield\_activated}_t\,],
  \qquad
  r^{\mathrm{no-shield}}_t \;=\; w_{\mathrm{no-sh}}\,
    \mathbb{1}[\,\neg \text{shield\_activated}_t\,].
  \label{eq:shield-reward}
\end{equation}
Either \(w_{\mathrm{sh}}\) or \(w_{\mathrm{no-sh}}\) can be zero, which makes it possible to
penalise shield activations, to reward the agent for not needing the shield, or to run in a
neutral configuration.

\paragraph{Idle penalty and shield grace period.}
An idle penalty discourages the agent from stopping. To prevent the penalty from punishing
legitimate shield-induced stalls, a \emph{grace period} of \(G\) steps is opened the first time
the shield activates in an episode. The counter is re-armed \emph{only} when the shield fires
and the grace period has already expired, which prevents an always-active shield from
suppressing the idle penalty forever.

\paragraph{Drift penalty.}
Asymmetric driving (e.g.\ hugging the right edge) is penalised by
\begin{equation}
  p^{\mathrm{drift}}_t \;=\; w_{\mathrm{dr}}\, \max\!\big(0,\, |d^{\mathrm{L}}_t - d^{\mathrm{R}}_t| - 0.3\big)\;
     \mathbb{1}\!\left[\min(d^{\mathrm{L}}_t, d^{\mathrm{R}}_t) < 0.35\right].
  \label{eq:p-drift}
\end{equation}

\paragraph{Lane-change transition guard.}
When the Waypoint API reports \texttt{lane\_change\_permitted} and \(|\ell_t| > 0.5\), the
penalties \(p^{\mathrm{inv}}_t\), \(p^{\mathrm{edge}}_t\) and \(p^{\mathrm{drift}}_t\) are all
set to zero. This keeps the cost of a legitimate lane-change manoeuvre from blocking the
behaviour during training.

The default weights of all components are listed in
Appendix~\ref{app:hyperparameters}, Table~\ref{tab:hp-reward}.

\section{Curriculum manager}\label{sec:design-curriculum}

The curriculum manager divides the interval \([0,\, N_{\max}]\) of NPC vehicles uniformly in
five stages: \((0,\, N_{\max}/4,\, N_{\max}/2,\, 3N_{\max}/4,\, N_{\max})\). For the default
\(N_{\max}=40\) this yields \((0,\,10,\,20,\,30,\,40)\). The manager keeps two counters: the
number of episodes at the current stage \(E_{\mathrm{stage}}\) and a decay-weighted rollback
counter \(C_{\mathrm{rb}}\).

\paragraph{Advancement rule.}
After \(E_{\mathrm{stage}} \ge E^{\star} = 100\) episodes, if the stage-specific condition of
Table~\ref{tab:curriculum-rules} is satisfied, the manager advances one stage and resets both
counters to zero.

\paragraph{Rollback rule.}
After \(E_{\mathrm{stage}} > E^{\star}\), on every episode the counter \(C_{\mathrm{rb}}\)
is incremented by one if the stage is \emph{collapsing} and decremented by one otherwise (but
not below zero). When \(C_{\mathrm{rb}} \ge C^{\star} = 50\), the manager rolls back one stage
and resets all counters. The decay by \(-1\) on non-collapsing episodes prevents a streak of
bad luck from triggering a rollback that is not warranted by a sustained performance drop.

\begin{table}[h]
  \centering
  \caption{Advancement and rollback thresholds per curriculum stage. All rates are 100-episode
           rolling averages.}
  \label{tab:curriculum-rules}
  \small
  \begin{tabular}{@{}clll@{}}
    \toprule
    \textbf{Stage} & \textbf{NPCs} & \textbf{Advance if} & \textbf{Rollback if} \\
    \midrule
    0 & \(0\)                 & \(\text{offroad\_rate} < 0.20 \wedge \text{avg\_reward} > 0\)
                              & \textemdash\ \\
    1 & \(N_{\max}/4\)        & \(\text{crash\_rate} < 0.25 \wedge \text{offroad\_rate} < 0.15\)
                              & \(\text{crash\_rate} > 0.45 \vee \text{offroad\_rate} > 0.40\) \\
    2 & \(N_{\max}/2\)        & \(\text{crash\_rate} < 0.15\)
                              & \(\text{crash\_rate} > 0.40\) \\
    3 & \(3N_{\max}/4\)       & \(\text{crash\_rate} < 0.10\)
                              & \(\text{crash\_rate} > 0.35\) \\
    4 & \(N_{\max}\)          & \textemdash\ (terminal)
                              & \(\text{crash\_rate} > 0.30\) \\
    \bottomrule
  \end{tabular}
\end{table}

\paragraph{Non-disruptive NPC update.}
When the stage changes, the manager does \emph{not} rebuild the environment. It updates the
attribute \texttt{CarlaEnv.num\_npc\_vehicles} in place; the change takes effect at the next
\texttt{env.reset()}. This avoids the deadlock observed in earlier prototypes that had to tear
down the wrapper chain every time the stage changed.

\section{Live metrics pipeline}\label{sec:design-metrics}

Training metrics are written to a per-run SQLite database at
\texttt{runs/\textless run\_name\textgreater/metrics.sqlite}. The logger
(\texttt{src/Metrics/live\_metrics.py}) opens the database in \emph{Write-Ahead Logging}
(WAL)~\cite{sqlitewal} mode, which allows concurrent reads and writes: the training process
writes while the Streamlit dashboard reads.

Two tables are used, corresponding to the two time axes of the experiment:

\begin{itemize}
  \item \texttt{metrics\_episode}: one row per episode, indexed by episode number. Columns
        include the reward components, the safety rates, the curriculum stage, and the
        LIDAR-derived minimum distances.
  \item \texttt{metrics\_update}: one row per PPO update, indexed by the cumulative timestep.
        Columns include the policy loss, the value loss, the approximate KL, the number of
        epochs that actually ran, the gradient norm and the current learning rate.
\end{itemize}

The dashboard (\texttt{dataVisualizer.py}) reads both tables every few seconds and renders
interactive Plotly charts. Unknown metric keys are rendered as soft-typed line plots, which
keeps the dashboard forward-compatible with future additions to the logger.

% ============================================================
%   Chapter 5  \textemdash\  Implementation
% ============================================================

\chapter{Implementation}\label{ch:implementation}

This chapter describes how the design of Chapter~\ref{ch:design} is materialised in Python code.
Section~\ref{sec:impl-stack} introduces the technology stack and
Section~\ref{sec:impl-layout} the project layout. The subsequent sections cover each component:
the CARLA wrapper (\S\ref{sec:impl-env}), the PPO agent (\S\ref{sec:impl-agent}), the two shields
(\S\ref{sec:impl-shields}), the reward shaper (\S\ref{sec:impl-reward}), the curriculum manager
(\S\ref{sec:impl-curriculum}), the entry-point pipelines (\S\ref{sec:impl-pipelines}) and
reproducibility concerns (\S\ref{sec:impl-repro}).

\section{Technology stack}\label{sec:impl-stack}

The framework is written in pure Python (\(\ge 3.10\)) and relies on the libraries summarised in
Table~\ref{tab:stack}. Licences are reported because the framework is meant to be distributed as
open academic software.

\begin{table}[h]
  \centering
  \caption{Main third-party dependencies.}
  \label{tab:stack}
  \small
  \begin{tabular}{@{}lll@{}}
    \toprule
    \textbf{Library} & \textbf{Role}                              & \textbf{Licence} \\
    \midrule
    \texttt{carla}         & Python binding for CARLA 0.9.16            & MIT \\
    \texttt{gymnasium}     & Standard RL environment API                & MIT \\
    \texttt{torch}         & Deep learning backend                      & BSD-3 \\
    \texttt{numpy}         & Numerical arrays                           & BSD-3 \\
    \texttt{matplotlib}    & Evaluation dashboard                       & PSF-based \\
    \texttt{streamlit}     & Live training dashboard                    & Apache-2.0 \\
    \texttt{plotly}        & Interactive charts in the dashboard        & MIT \\
    \texttt{sqlite3} (stdlib) & Metrics database                        & Public domain \\
    \texttt{ruff}          & Linter and formatter                       & MIT \\
    \texttt{pytest}        & Unit testing of the reward shaper          & MIT \\
    \bottomrule
  \end{tabular}
\end{table}

\section{Project layout}\label{sec:impl-layout}

The repository root is organised as follows (directories without relevance to the thesis are
omitted):

\begin{itemize}
  \item \texttt{main\_train.py} \textemdash\ training entry point.
  \item \texttt{main\_eval.py} \textemdash\ evaluation entry point.
  \item \texttt{dataVisualizer.py} \textemdash\ Streamlit dashboard that reads the SQLite
        metrics database.
  \item \texttt{src/CARLA/Env/carla\_env.py} \textemdash\ base Gymnasium environment.
  \item \texttt{src/PPO/} \textemdash\ \texttt{ActorCritic.py}, \texttt{ppo\_agent.py},
        \texttt{RunningMeanStd.py}.
  \item \texttt{src/safety\_shield.py} \textemdash\ basic safety shield.
  \item \texttt{src/Adaptative\_Shield/} \textemdash\ \texttt{adaptive\_horizon\_shield.py},
        \texttt{BicycleModel.py}.
  \item \texttt{src/reward\_shaper.py} \textemdash\ dense reward shaper.
  \item \texttt{src/curriculumManager.py} \textemdash\ performance-driven curriculum.
  \item \texttt{src/Metrics/} \textemdash\ live metrics logger and offline evaluation reporter.
  \item \texttt{tests/test\_reward\_balance.py} \textemdash\ unit tests on the reward shaper.
  \item \texttt{data/models/} \textemdash\ serialised checkpoints.
  \item \texttt{runs/} \textemdash\ per-run SQLite databases and TensorBoard logs.
\end{itemize}

Every module exposes either a class that inherits from \texttt{gymnasium.Wrapper} (the shields,
the reward shaper) or a small self-contained helper (the curriculum, the running-mean estimator).
There is no inter-module state other than what flows through the Gymnasium interface, which makes
individual components trivially testable.

\section{Environment wrapper}\label{sec:impl-env}

The class \texttt{CarlaEnv} inherits from \texttt{gymnasium.Env} and encapsulates the connection
to the CARLA server. At construction time it performs the following setup:

\begin{enumerate}
  \item Opens a TCP connection on \texttt{(host, port)} and loads the requested map.
  \item Enables synchronous mode with \texttt{fixed\_delta\_seconds = 0.05}.
  \item Instantiates the Traffic Manager on \texttt{tm\_port} and spawns the requested number
        of NPC vehicles.
  \item Sets the requested weather preset.
  \item Registers the collision, lane-invasion and semantic-LIDAR sensors.
  \item Defines \texttt{observation\_space} as \(\text{Box}([0,1]^{739})\) and
        \texttt{action\_space} as \(\text{Box}([-1,1]^2)\).
\end{enumerate}

On every \texttt{step(action)} call, the class converts the normalised action into a
\texttt{carla.VehicleControl}, ticks the simulator exactly once, reads the sensors, computes the
geometric features from the Waypoint API, builds the 739-dimensional observation, and derives the
base reward together with the \texttt{done} and \texttt{truncated} flags. The \texttt{info}
dictionary contains every scalar used by downstream wrappers.

\section{PPO agent}\label{sec:impl-agent}

The PPO agent is split into two files: \texttt{src/PPO/ActorCritic.py} contains the network
described in Figure~\ref{fig:actor-critic-arch}, and \texttt{src/PPO/ppo\_agent.py} contains
the training loop. The three public methods used by \texttt{main\_train.py} are
\texttt{select\_action}, \texttt{evaluate\_action} and \texttt{update}.
Figure~\ref{fig:select-action} shows the pseudocode of \texttt{select\_action} and
\texttt{evaluate\_action}; Figure~\ref{fig:ppo-update} shows the pseudocode of \texttt{update}.

\begin{figure}[h]
  \centering
  \includegraphics[width=0.95\linewidth]{imgs/code/select_action}
  \caption{Action selection and evaluation methods. \texttt{select\_action} normalises the
           vector block of the observation, samples from the policy (or returns the mean for
           deterministic evaluation), and returns the action, log probability and value estimate.
           \texttt{evaluate\_action} is called with \texttt{info["executed\_action"]} to
           recompute the log probability of the action that actually reached the environment.}
  \label{fig:select-action}
\end{figure}

\begin{figure}[h]
  \centering
  \includegraphics[width=0.95\linewidth]{imgs/code/ppo_update}
  \caption{PPO update loop with GAE backward pass, return normalisation by running standard
           deviation, and up to \(K=10\) epochs of clipped surrogate loss with optional
           value clipping and KL early stop.}
  \label{fig:ppo-update}
\end{figure}

\subsection{Credit assignment under a shield}\label{sec:impl-shield-credit}

The interplay between the shield and the PPO update is the single most error-prone point of the
framework. Let \(a_t\) be the action sampled from the policy and
\(a^\star_t = \Sigma(s_t, a_t)\) the action produced by the shield. The shield is a
\texttt{gymnasium.Wrapper}, so the environment sees \(a^\star_t\); the agent only receives it
through the \texttt{executed\_action} key of the \texttt{info} dictionary.

The training loop in \texttt{main\_train.py} stores \(a^\star_t\) (not \(a_t\)) in the rollout
buffer and recomputes the log probability of \(a^\star_t\) under the current policy via
\texttt{evaluate\_action}. Without this step, the gradient estimator would update the policy
towards the action that \emph{was not executed}, effectively decoupling the policy from the
environment and making training degenerate. With this step, the policy learns to produce actions
that are close to what the shield considered safe, which is precisely the desired behaviour: the
shield teaches the agent what \emph{not} to do.

The approximate KL used by the early-stop mechanism (Figure~\ref{fig:ppo-update}) is the
unbiased estimator of~\cite{schulman2020approxkl}. The KL target and the value-clip parameter
\(\xi\) are both exposed as command-line flags.

\section{Shield implementations}\label{sec:impl-shields}

Both shields inherit from \texttt{gymnasium.Wrapper} and implement the same public interface
(\texttt{reset}, \texttt{step}, \texttt{get\_statistics}, \texttt{reset\_statistics}). The
interchangeability is what makes \texttt{--shield\_type} a drop-in parameter.

\subsection{Basic shield}\label{sec:impl-basic}

The implementation of \texttt{CarlaSafetyShield} follows the correction law of
Section~\ref{sec:design-basic-shield} with steering gain \(k_s = 2.5\) and brake gain
\(k_b = 8.0\). The class keeps per-episode statistics of the sub-reason that triggered each
intervention (\texttt{front}, \texttt{side\_right}, \texttt{side\_left}, \texttt{lane\_right},
\texttt{lane\_left}, \texttt{heading}), surfaced by \texttt{get\_statistics()} and logged once
per episode.

\subsection{Adaptive-horizon shield}\label{sec:impl-adaptive}

The adaptive shield is more involved. Figure~\ref{fig:adaptive-step} shows the structure of its
\texttt{step} method, which classifies the current risk level, queries the bicycle model for
the appropriate horizon, checks the proposed action, searches for a safe alternative if needed,
and decorates the \texttt{info} dictionary before returning.

\begin{figure}[h]
  \centering
  \includegraphics[width=0.95\linewidth]{imgs/code/adaptive_shield_step}
  \caption{Adaptive-horizon shield \texttt{step} method. The semantic analysis extracts LIDAR
           and Waypoint features; the risk classifier selects the horizon \(H\); the trajectory
           check decides whether to pass the proposed action through or invoke the candidate
           search; the \texttt{info} dictionary is updated with shield diagnostics.}
  \label{fig:adaptive-step}
\end{figure}

Figure~\ref{fig:adaptive-check} shows the trajectory-check routine: it validates LIDAR
thresholds, checks for a pedestrian emergency, and then runs the bicycle model over \(H\) steps,
projecting each predicted position onto the road network via CARLA's Waypoint API.

\begin{figure}[h]
  \centering
  \includegraphics[width=0.95\linewidth]{imgs/code/adaptive_trajectory_check}
  \caption{Trajectory safety check of the adaptive shield. The threshold multiplier \(\kappa\)
           scales the LIDAR and lateral thresholds according to the current risk level. The
           bicycle model predicts \(H\) positions; each is projected onto the nearest driving
           lane and its lateral offset is checked. The lane-change relaxation (\(\tau_l = 1.2\))
           prevents the shield from blocking legitimate manoeuvres.}
  \label{fig:adaptive-check}
\end{figure}

The candidate-search routine follows the priority-ordered policy of
Section~\ref{sec:design-adaptive-shield}. Each candidate is a hard-coded two-dimensional vector;
the first that passes \texttt{\_check\_trajectory\_safety} is returned, and if none passes, the
fallback action is returned.

\subsection{Bicycle model}\label{sec:impl-bicycle}

\texttt{BicycleModel.predict\_trajectory} is a loop of \(H\) discrete steps that implements
the forward-Euler integration with closed-form arc update when \(|\delta| \ge 10^{-4}\) rad.
The class is a plain Python object with no PyTorch or CARLA dependency and is therefore
trivially unit-testable.

\section{Reward shaper}\label{sec:impl-reward}

\texttt{CarlaRewardShaper} reads approximately 20 scalars from the \texttt{info} dictionary and
produces the components of Equation~\eqref{eq:shaped-reward} from Chapter~\ref{ch:design}. It
also writes every sub-component back into \texttt{info} under clearly named keys
(\texttt{speed\_bonus}, \texttt{lane\_center\_bonus}, \texttt{heading\_bonus}, etc.) so the
metrics logger can record the per-step reward decomposition. Three minor state variables are kept
across steps:

\begin{itemize}
  \item \texttt{\_last\_steering}: needed by the smoothness penalty.
  \item \texttt{\_last\_milestone}: needed by the progress bonus to avoid double-crediting the
        same \(25\,\mathrm{m}\) segment.
  \item \texttt{\_shield\_grace\_steps}: count-down of the grace period that temporarily
        disables the idle penalty after a shield activation.
\end{itemize}

The unit test \texttt{tests/test\_reward\_balance.py} feeds a crafted \texttt{info} dictionary
to the shaper and asserts that each component matches its analytical expression. The test runs
in \(< 1\,\mathrm{s}\) without a CARLA server and is used as a pre-commit safeguard against
regressions.

\section{Curriculum manager}\label{sec:impl-curriculum}

\texttt{CurriculumManager} is a small self-contained class. Its only state is the tuple
\((i,\, E_{\mathrm{stage}},\, C_{\mathrm{rb}})\) where \(i\) is the current stage index. The
main routine is \texttt{step(offroad\_rate, crash\_rate, avg\_reward)}, called from
\texttt{main\_train.py} after each episode; it returns a pair
\((\text{target\_npcs},\, \text{event} \in \{\texttt{none},\, \texttt{advance},\, \texttt{rollback}\})\).
The training loop reacts to the event by updating \texttt{CarlaEnv.num\_npc\_vehicles} in place,
so no reconnection to CARLA is required (Section~\ref{sec:design-curriculum}).

\section{Training and evaluation pipelines}\label{sec:impl-pipelines}

Both entry points share the function \texttt{build\_env(args)} that constructs the wrapper chain
of Figure~\ref{fig:wrapper-chain} with the shield type, reward weights, curriculum NPC count and
map chosen by the command line. Keeping the same builder on both sides guarantees that the
evaluation environment matches the training environment exactly.

\subsection{Training pipeline}\label{sec:impl-training}

Figure~\ref{fig:train-algorithm} summarises the training loop.

\begin{figure}[h]
  \centering
  \includegraphics[width=0.95\linewidth]{imgs/code/training_loop}
  \caption{High-level training loop (\texttt{main\_train.py}). The inner loop collects
           transitions and appends the \emph{executed} action (after shield correction) to the
           rollout buffer. A PPO update is triggered every 2048 steps. After each episode, the
           curriculum manager is polled; if it advances or rolls back a stage, the NPC count
           is updated in place without restarting the environment.}
  \label{fig:train-algorithm}
\end{figure}

The key correctness invariant of the training loop is step~11 of Figure~\ref{fig:train-algorithm}:
the log probability stored in the buffer corresponds to the \emph{executed} action
\(a^\star_t\), not the sampled \(a_t\). This is what makes the PPO gradient unbiased in the
presence of shield interventions (Section~\ref{sec:impl-shield-credit}).

\subsection{Evaluation pipeline}\label{sec:impl-evaluation}

The evaluation pipeline (\texttt{main\_eval.py}) rebuilds the wrapper chain with the same
\texttt{build\_env} function but sets all reward weights to zero so that the reported episode
reward is the pure CARLA base reward. The shield is still attached so that action modifications
are visible. An optional Matplotlib dashboard displays the LIDAR polar plot, speed gauge, lateral
offset marker, and a monospaced panel with live information (risk level, shield activation,
steering, throttle/brake, total distance, collisions). At the end of the run a
\texttt{SafetyMetricsReporter} prints a summary table aggregating success rate, crash rate,
off-road rate, shield activations and LIDAR-derived minimum distances.

\section{Persistence and reproducibility}\label{sec:impl-repro}

\subsection{Checkpoint format}\label{sec:impl-checkpoints}

Checkpoints are saved with \texttt{torch.save} as a dictionary of two entries: the
\texttt{state\_dict} of the \texttt{ActorCritic} network and the state of the
\texttt{RunningMeanStd} normaliser. The normaliser state is essential for evaluation: if it were
missing, the running statistics would be rebuilt from scratch and the input distribution would not
match the one the policy trained on. A shape check on load tolerates the legacy format (state
dictionary only) for backwards compatibility with the earliest checkpoints.

\subsection{Seeds and determinism}\label{sec:impl-seeds}

Two distinct seeds are used:

\begin{itemize}
  \item \texttt{seed=42} for training. Propagated to NumPy, Python's \texttt{random} module,
        PyTorch, CARLA's Traffic Manager and the environment's episode-counter random state.
  \item \texttt{seed=100} for evaluation, intentionally different from the training seed so
        that evaluation trajectories are not a subset of those seen during learning.
\end{itemize}

Full bit-exact determinism across hardware is not achievable because of PyTorch's cuDNN
non-reproducibility and CARLA's floating-point physics divergence. However, determinism is
sufficient for the first fifty episodes of a run, which is long enough to diagnose regressions
introduced by a code change (acceptance criterion~AC-09, Table~\ref{tab:acceptance-criteria}).

\subsection{Tooling}\label{sec:impl-tooling}

The source code is linted and formatted with \texttt{ruff}. Reward regressions are caught by
\texttt{pytest tests/test\_reward\_balance.py}. The utility scripts under \texttt{utils/}
provide helpers to list, inspect and compare checkpoints (\texttt{ModelManager}), to analyse
the metrics of finished runs (\texttt{ExperimentAnalyzer}) and to generate configuration files
from a template (\texttt{ConfigurationTemplate}). None of these utilities are required for
training; they are purely auxiliary.

% ============================================================
%   Chapter 6  \textemdash\  Experimentation and Results
% ============================================================
%
% NOTE TO THE AUTHOR: this chapter is written as a STRUCTURAL TEMPLATE.
% Every empirical number appears as "TBD" and is marked with a comment
% explaining the exact source (metrics.sqlite column, dashboard panel,
% etc.). Replace each TBD with the final figures from your runs.
% ============================================================

\chapter{Experimentation and Results}\label{ch:experimentation}

This chapter reports the empirical evaluation of the framework. It is organised around the four
configurations that matter: the baseline agent with no shield (Section~\ref{sec:exp-baseline}),
the basic shield (Section~\ref{sec:exp-basic}), the adaptive-horizon shield
(Section~\ref{sec:exp-adaptive}), and the effect of enabling or disabling the curriculum
(Section~\ref{sec:exp-curriculum}). Section~\ref{sec:exp-protocol} fixes the experimental
protocol and Section~\ref{sec:exp-metrics} formally defines every metric reported hereafter;
Section~\ref{sec:exp-discussion} offers a cross-configuration discussion.

All numerical entries marked as \emph{TBD} are to be filled in with the exact values obtained
from the SQLite metrics database of each run. The SQL queries used to compute the reported
statistics are documented in Appendix~\ref{app:reproducibility}.

\section{Experimental protocol}\label{sec:exp-protocol}

\paragraph{Hardware and software.}
All experiments are carried out on the workstation described in
Section~\ref{sec:req-constraints}:
\begin{itemize}
  \item CPU: \emph{TBD}.% replace with e.g. "AMD Ryzen 7 7800X3D, 8 cores / 16 threads"
  \item GPU: \emph{TBD}.% e.g. "NVIDIA RTX 4070, 12 GB VRAM"
  \item System RAM: \emph{TBD}.% e.g. "32 GB DDR5"
  \item Operating system: Windows 11 Pro.
  \item Python: \emph{TBD}.% e.g. "3.11.9"
  \item PyTorch: \emph{TBD} with CUDA \emph{TBD}.% e.g. "2.3.0 with CUDA 12.1"
  \item CARLA: 0.9.16, launched with \texttt{.\textbackslash CarlaUE4.exe -RenderOffScreen
        -carla-port=2000}.
\end{itemize}

\paragraph{Benchmark scenario.}
The main benchmark is the Town04 highway loop with clear weather (\texttt{ClearNoon}). The
vehicle spawns at a random waypoint each episode and the goal is to travel at least
\(250\,\mathrm{m}\) along the map's route graph within \(1000\) simulator steps
(\(50\,\mathrm{s}\)).

\paragraph{Training budget.}
Every configuration is trained for \(T = 2\,500\) episodes, with PPO updates every
\(2\,048\) steps (\(\sim\!\emph{TBD}\) updates per run).% fill in updates-per-run
The learning rate starts at \(10^{-4}\) and anneals cosine-style to \(10^{-5}\); the target KL
for the early stop is \(0.08\); the entropy bonus is \(0.02\); the value-loss coefficient is
\(0.25\); all other hyperparameters take the defaults of Appendix~\ref{app:hyperparameters}.

\paragraph{Random seeds.}
Every training run uses \(\mathrm{seed}=42\); every evaluation uses \(\mathrm{seed}=100\). Three
independent runs per configuration are carried out when time permits, by varying the seed to
\(\{42, 43, 44\}\); otherwise the single \(\mathrm{seed}=42\) run is used. The reported aggregate
statistics (mean, standard deviation, median) are computed across these three runs when
applicable.

\paragraph{Curriculum.}
Unless explicitly stated otherwise, the curriculum is enabled with \(N_{\max} = 40\). The ablation
of Section~\ref{sec:exp-curriculum} disables it.

\paragraph{Evaluation protocol.}
After training, each agent is evaluated on \(10\) episodes with \(\mathrm{seed}=100\). The shield
configuration used at evaluation is identical to the one used during training: removing the
shield at evaluation would compromise the comparability of the results because a policy trained
with a shield has learned to rely on its interventions.

\section{Metrics}\label{sec:exp-metrics}

The metrics reported throughout this chapter are exactly those produced by the training and
evaluation pipelines. Their operational definitions are collected here to avoid ambiguity.

\begin{description}
  \item[\textbf{Reward/Raw\_Episode.}] Undiscounted sum of shaped rewards of a single episode.
  \item[\textbf{Reward/Average\_100\_Episodes.}] Rolling mean of \emph{Reward/Raw\_Episode} over
        the last \(100\) episodes.
  \item[\textbf{Training/Success\_Rate.}] Fraction of the last \(100\) episodes that ended with
        \texttt{info["arrive\_dest"] = True}.
  \item[\textbf{Training/Crash\_Rate.}] Fraction of the last \(100\) episodes that ended with
        \texttt{info["collision"] = True}.
  \item[\textbf{Training/Offroad\_Rate.}] Fraction of the last \(100\) episodes that ended with
        \texttt{info["out\_of\_road"] = True}.
  \item[\textbf{Safety/Shield\_Activations.}] Number of steps in which the shield overrode the
        agent during the episode.
  \item[\textbf{Safety/Shield\_Rate.}] \emph{Safety/Shield\_Activations} divided by the number of
        steps in the episode.
  \item[\textbf{Safety/Semantic/\{Dynamic,Static,Pedestrian\}\_Interventions.}] Per-episode
        intervention count by threat category, as produced by
        \texttt{CarlaAdaptiveHorizonShield.get\_statistics}.
  \item[\textbf{Safety/Semantic/\{Safe,Warning,Critical\}\_Step\_Rate.}] Fractions of steps spent
        in each risk level.
  \item[\textbf{CARLA/Mean\_Speed\_kmh.}] Arithmetic mean of \texttt{speed\_kmh} across the steps
        of the episode.
  \item[\textbf{CARLA/Speed\_Compliance\_Rate.}] Fraction of steps (with a positive speed limit)
        that satisfy \(\texttt{speed\_kmh} \le 1.05\,\texttt{speed\_limit\_kmh}\).
  \item[\textbf{Safety/Min\_\{Vehicle,Pedestrian\}\_Distance\_m.}] Minimum distance over the
        episode from the ego vehicle to the nearest vehicle / pedestrian.
  \item[\textbf{Loss/Policy\_Loss, Loss/Value\_Loss, Loss/Grad\_Norm, Training/Entropy,
        Training/Approx\_KL, Training/Epochs\_Run.}] Per-update diagnostics produced by the PPO
        loop.
\end{description}

\section{Baseline\,: \texttt{shield\_type=none}}\label{sec:exp-baseline}

The baseline corresponds to an unconstrained PPO agent that follows exactly the same training
protocol but without any shield. Its purpose is twofold: (i) it quantifies what a raw
RL agent can learn on Town04 given the curriculum and the dense reward, and (ii) it provides
the reference numbers against which every shielded configuration is compared.

\paragraph{Training curves.}
Figure~\ref{fig:baseline-reward} shows the rolling reward, the success rate and the safety
rates over the \(2\,500\) training episodes. The agent reaches a final rolling reward of
\emph{TBD} with a success rate of \emph{TBD} and a crash rate of \emph{TBD}.% pull from metrics_episode

\begin{figure}[h]
  \centering
  \fbox{\parbox{0.9\textwidth}{\centering\vspace{1cm} \emph{TBD: figure file, e.g.}
        \texttt{figures/baseline\_reward.pdf} \vspace{1cm}}}
  \caption{Baseline run (\texttt{--shield\_type none}). Top: rolling reward over \(100\)
           episodes. Bottom left: success, crash and off-road rates. Bottom right: number of
           shield activations per episode (identically zero for this configuration).}
  \label{fig:baseline-reward}
\end{figure}

\paragraph{Curriculum trajectory.}
Figure~\ref{fig:baseline-curriculum} shows the evolution of the curriculum stage. Rollbacks, if
any, are visible as downward steps.

\begin{figure}[h]
  \centering
  \fbox{\parbox{0.9\textwidth}{\centering\vspace{1cm} \emph{TBD: figure file, e.g.}
        \texttt{figures/baseline\_curriculum.pdf}\vspace{1cm}}}
  \caption{Baseline run. Curriculum stage (0 to 4) across training episodes and associated NPC
           count.}
  \label{fig:baseline-curriculum}
\end{figure}

\paragraph{Evaluation.}
Table~\ref{tab:baseline-eval} summarises the \(10\) evaluation episodes.

\begin{table}[h]
  \centering
  \caption{Baseline evaluation, \(10\) episodes, \(\mathrm{seed}=100\).}
  \label{tab:baseline-eval}
  \small
  \begin{tabular}{@{}lc@{}}
    \toprule
    \textbf{Metric} & \textbf{Baseline} \\
    \midrule
    Average reward        & TBD $\pm$ TBD \\
    Success rate          & TBD \\
    Crash rate            & TBD \\
    Off-road rate         & TBD \\
    Timeout rate          & TBD \\
    Mean speed (km/h)     & TBD \\
    Speed-compliance rate & TBD \\
    \bottomrule
  \end{tabular}
\end{table}

\section{Basic shield}\label{sec:exp-basic}

The second configuration enables the basic shield (\texttt{--shield\_type basic}). Shield
interventions are now logged per step, per episode and per semantic reason
(\texttt{front}, \texttt{side\_right}, \texttt{side\_left}, \texttt{lane\_right},
\texttt{lane\_left}, \texttt{heading}).

\paragraph{Training curves.}
Figure~\ref{fig:basic-training} mirrors Figure~\ref{fig:baseline-reward}; the additional panel
reports the shield-activation rate per episode. The final rolling reward is \emph{TBD};
the final crash rate is \emph{TBD}, compared to \emph{TBD} for the
baseline.

\begin{figure}[h]
  \centering
  \fbox{\parbox{0.9\textwidth}{\centering\vspace{1cm} \emph{TBD: figure file, e.g.}
        \texttt{figures/basic\_training.pdf}\vspace{1cm}}}
  \caption{Basic-shield run. Top: rolling reward. Bottom left: safety rates.
           Bottom right: shield-activation rate per episode.}
  \label{fig:basic-training}
\end{figure}

\paragraph{Intervention breakdown.}
Table~\ref{tab:basic-breakdown} gives the mean number of interventions per episode by sub-reason,
averaged over the last \(100\) training episodes.

\begin{table}[h]
  \centering
  \caption{Basic shield: interventions per episode by sub-reason (last \(100\) episodes).}
  \label{tab:basic-breakdown}
  \small
  \begin{tabular}{@{}lc@{}}
    \toprule
    \textbf{Sub-reason}  & \textbf{Mean interventions / episode} \\
    \midrule
    \texttt{front}       & TBD \\
    \texttt{side\_right} & TBD \\
    \texttt{side\_left}  & TBD \\
    \texttt{lane\_right} & TBD \\
    \texttt{lane\_left}  & TBD \\
    \texttt{heading}     & TBD \\
    \bottomrule
  \end{tabular}
\end{table}

\paragraph{Evaluation.}
Table~\ref{tab:basic-eval} reports the \(10\) evaluation episodes.

\begin{table}[h]
  \centering
  \caption{Basic shield evaluation, \(10\) episodes, \(\mathrm{seed}=100\).}
  \label{tab:basic-eval}
  \small
  \begin{tabular}{@{}lc@{}}
    \toprule
    \textbf{Metric} & \textbf{Basic shield} \\
    \midrule
    Average reward           & TBD $\pm$ TBD \\
    Success rate             & TBD \\
    Crash rate               & TBD \\
    Off-road rate            & TBD \\
    Timeout rate             & TBD \\
    Mean speed (km/h)        & TBD \\
    Speed-compliance rate    & TBD \\
    Shield activations / ep  & TBD \\
    \bottomrule
  \end{tabular}
\end{table}

\section{Adaptive-horizon shield}\label{sec:exp-adaptive}

The third configuration enables the adaptive-horizon shield
(\texttt{--shield\_type adaptive}). In addition to the metrics reported for the basic shield,
the adaptive shield exposes the time spent in each risk level and the breakdown of
interventions by threat category (dynamic, static, pedestrian).

\paragraph{Training curves.}
Figure~\ref{fig:adaptive-training} mirrors Figure~\ref{fig:basic-training} with an extra panel
for the risk-level time fractions.

\begin{figure}[h]
  \centering
  \fbox{\parbox{0.9\textwidth}{\centering\vspace{1cm} \emph{TBD: figure file, e.g.}
        \texttt{figures/adaptive\_training.pdf}\vspace{1cm}}}
  \caption{Adaptive-shield run. Top-left: rolling reward. Top-right: safety rates.
           Bottom-left: shield-activation rate by threat category. Bottom-right: fraction of
           steps per risk level.}
  \label{fig:adaptive-training}
\end{figure}

\paragraph{Risk distribution.}
Table~\ref{tab:adaptive-risk} reports the average time spent in each risk level over the last
\(100\) training episodes.

\begin{table}[h]
  \centering
  \caption{Adaptive shield: average time per risk level (last \(100\) episodes).}
  \label{tab:adaptive-risk}
  \small
  \begin{tabular}{@{}lc@{}}
    \toprule
    \textbf{Risk level}  & \textbf{Fraction of steps} \\
    \midrule
    safe     & TBD \\
    warning  & TBD \\
    critical & TBD \\
    \bottomrule
  \end{tabular}
\end{table}

\paragraph{Intervention breakdown.}
Table~\ref{tab:adaptive-breakdown} reports the intervention count per episode by threat
category.

\begin{table}[h]
  \centering
  \caption{Adaptive shield: interventions per episode by threat category (last \(100\)
           episodes).}
  \label{tab:adaptive-breakdown}
  \small
  \begin{tabular}{@{}lc@{}}
    \toprule
    \textbf{Category}    & \textbf{Mean interventions / episode} \\
    \midrule
    dynamic    & TBD \\
    static     & TBD \\
    pedestrian & TBD \\
    \bottomrule
  \end{tabular}
\end{table}

\paragraph{Evaluation.}
Table~\ref{tab:adaptive-eval} reports the \(10\) evaluation episodes.

\begin{table}[h]
  \centering
  \caption{Adaptive shield evaluation, \(10\) episodes, \(\mathrm{seed}=100\).}
  \label{tab:adaptive-eval}
  \small
  \begin{tabular}{@{}lc@{}}
    \toprule
    \textbf{Metric} & \textbf{Adaptive shield} \\
    \midrule
    Average reward                   & TBD $\pm$ TBD \\
    Success rate                     & TBD \\
    Crash rate                       & TBD \\
    Off-road rate                    & TBD \\
    Timeout rate                     & TBD \\
    Mean speed (km/h)                & TBD \\
    Speed-compliance rate            & TBD \\
    Shield activations / ep          & TBD \\
    Mean time in \emph{critical} (\%)& TBD \\
    \bottomrule
  \end{tabular}
\end{table}

\paragraph{Cross-configuration comparison.}
Table~\ref{tab:config-comparison} aggregates the evaluation of the three configurations
side-by-side.

\begin{table}[h]
  \centering
  \caption{Side-by-side evaluation comparison (\(10\) episodes, seed \(100\)).}
  \label{tab:config-comparison}
  \small
  \begin{tabular}{@{}lccc@{}}
    \toprule
    \textbf{Metric}            & \textbf{None} & \textbf{Basic} & \textbf{Adaptive} \\
    \midrule
    Average reward             & TBD & TBD & TBD \\
    Success rate               & TBD & TBD & TBD \\
    Crash rate                 & TBD & TBD & TBD \\
    Off-road rate              & TBD & TBD & TBD \\
    Speed-compliance rate      & TBD & TBD & TBD \\
    Min. vehicle distance (m)  & TBD & TBD & TBD \\
    Shield activations / ep    & 0   & TBD & TBD \\
    \bottomrule
  \end{tabular}
\end{table}

\section{Effect of the curriculum}\label{sec:exp-curriculum}

To quantify the contribution of the curriculum we repeat the adaptive-shield configuration with
\texttt{--no-curriculum}, i.e.\ with the manager disabled and the agent exposed to the maximum
NPC count from episode~\(1\). Everything else is held fixed.

\paragraph{Training curves.}
Figure~\ref{fig:curriculum-ablation} compares the rolling reward of both runs.
Without the curriculum, the agent takes \emph{TBD} episodes to reach the reward level that the
curriculum-enabled run reaches after \emph{TBD} episodes; the final rolling crash rate without
curriculum is \emph{TBD} against \emph{TBD} with curriculum.% pull from metrics_episode

\begin{figure}[h]
  \centering
  \fbox{\parbox{0.9\textwidth}{\centering\vspace{1cm} \emph{TBD: figure file, e.g.}
        \texttt{figures/curriculum\_ablation.pdf}\vspace{1cm}}}
  \caption{Curriculum ablation: adaptive shield with curriculum enabled (solid) vs.\ disabled
           (dashed). Top: rolling reward. Bottom: rolling crash rate.}
  \label{fig:curriculum-ablation}
\end{figure}

\paragraph{Rollback analysis.}
When the curriculum is enabled, Table~\ref{tab:curriculum-events} enumerates the stage
transitions and rollbacks recorded during training.

\begin{table}[h]
  \centering
  \caption{Curriculum events during the adaptive-shield run.}
  \label{tab:curriculum-events}
  \small
  \begin{tabular}{@{}llll@{}}
    \toprule
    \textbf{Episode} & \textbf{Event} & \textbf{From stage} & \textbf{To stage} \\
    \midrule
    TBD & advance  & TBD & TBD \\
    TBD & advance  & TBD & TBD \\
    TBD & rollback & TBD & TBD \\
    \emph{\ldots} & \emph{\ldots} & \emph{\ldots} & \emph{\ldots} \\
    \bottomrule
  \end{tabular}
\end{table}

\section{Discussion}\label{sec:exp-discussion}

This section will discuss the trade-offs observed across configurations. The analysis will cover
at least the following four axes; each is introduced here as a question that the empirical
numbers of Sections~\ref{sec:exp-baseline}--\ref{sec:exp-curriculum} will answer.

\begin{enumerate}
  \item \textbf{Safety vs.\ performance.} How much reward does a shield cost in exchange for a
        lower crash rate? Is the adaptive shield strictly Pareto-better than the basic shield,
        or does one of them dominate on one axis while the other dominates on the complementary
        axis?
  \item \textbf{Intervention regime.} What fraction of a typical episode is spent with the
        shield active, and how does this fraction evolve during training? A decreasing
        shield-activation rate is consistent with the hypothesis that the policy is internalising
        the shield's behaviour.
  \item \textbf{Threat composition.} Which fraction of interventions is due to dynamic threats
        (other vehicles), static threats (barriers, walls) or pedestrians? How do these
        fractions change across curriculum stages?
  \item \textbf{Curriculum value.} Does the performance-driven curriculum accelerate learning
        and, more importantly, does it reduce the worst-case transient crash rate of the
        adaptive-shield run? The empirical answer to this question informs the recommendation
        of whether to enable the curriculum by default.
\end{enumerate}

Each enumerated item will be supported by a specific cross-reference to the tables and figures
above. Once the TBD entries have been filled in, a concise summary (bulleted or tabular) will be
inserted at the end of this section to provide a reader-friendly overview of the
findings.

% ============================================================
%   Chapter 7  --  Reducing Shield Dependence
% ============================================================

\chapter{Reducing Shield Dependence}\label{ch:dependence}

The evaluation protocol of Chapter~\ref{ch:experimentation} makes an explicit concession: the
shield is kept active at evaluation time because ``a policy trained with a shield has learned to
rely on its interventions''. This chapter takes that concession as its research question. It
first \emph{measures} how dependent the trained agent actually is
(Section~\ref{sec:dep-problem}--\ref{sec:dep-probe}), then studies two strategies to reduce that
dependence: a Large-Language-Model (LLM) meta-controller that dynamically loosens the shield
(Section~\ref{sec:dep-llm}), and a behavioural-cloning scheme that turns the shield into a
\emph{teacher} (Section~\ref{sec:dep-teacher}). The first strategy is a carefully evaluated
\emph{negative} result; the second succeeds to the point where the agent drives \emph{better}
without the shield than with it. Section~\ref{sec:dep-collapse} analyses an instructive failure
of the teacher scheme, and Section~\ref{sec:dep-discussion} distils the central lesson: shield
dependence is a \emph{learning} problem, not a \emph{scheduling} problem.

\section{The shield-dependence problem}\label{sec:dep-problem}

Two design choices, each reasonable in isolation, conspire to make dependence the optimal
behaviour for the agent.

\paragraph{No cost for being shielded (R1).} The reward shaper of
Section~\ref{sec:design-reward} contains no term that penalises shield activation. When the
shield overrides a reckless action it keeps the vehicle on the road and moving, so the agent
collects \emph{more} return than if its own unsafe action had reached the simulator. Being
shielded is therefore free and, in fact, beneficial: the reward landscape makes the dependent
policy reward-optimal.

\paragraph{No gradient where it matters (R2).} The masked PPO objective
(Section~\ref{sec:design-agent-ppo}) computes the policy gradient only over unshielded steps, weighting
each transition by \((1-\alpha_t)^2\), where \(\alpha_t\in[0,1]\) is the shield intensity. The
shielded steps are precisely the dangerous, near-boundary states. The actor therefore receives a
learning signal only on the easy states the shield never touched, and is \emph{structurally
prevented} from learning how to act at the boundary. The critic does observe those states (its
loss is unmasked), but the actor cannot act on them. The framework even stores the shield's
executed action and recomputes its log-probability, ideal imitation data, and then the
mask discards it.

Together, R1 and R2 produce an agent that ``rides the shield boundary'': it calibrates its own
imprecision to whatever tolerance the shield currently allows, so its competence \emph{with} the
shield says nothing about its competence \emph{without} it.

\section{Measuring dependence: the shield-off probe}\label{sec:dep-probe}

Quantifying dependence requires an evaluation that the shaped reward and the shield-activation
rate both hide. Two instruments are added.

\paragraph{Shield-off probe.} Every \(N_{\mathrm{probe}}\) training episodes the shield is placed
in a \emph{bypass} mode in which it forwards the proposed action to the simulator unchanged, and
the deterministic policy is rolled out for a handful of episodes with no buffer writes and no
optimiser steps; the observation normaliser is frozen and restored afterwards so the probe does
not contaminate training. The resulting shield-off success, crash and off-road rates are the
\emph{true} learning curve: they answer ``can the agent drive on its own yet?''.

\paragraph{Dynamic\,/\,static reliance split.} The per-episode intervention count is separated
into a \emph{dynamic} rate (vehicles and pedestrians) and a \emph{static} rate (barriers, walls,
guardrails). On the Town04 highway the static channel dominates by roughly \(30{:}1\), so the
overall shield-activation rate is mostly a function of the map geometry rather than of the
agent's skill; the split exposes the part of reliance that is actually learnable.

\section{First approach: an LLM meta-controller}\label{sec:dep-llm}

\subsection{Design}

The first hypothesis, inspired by the use of LLMs as meta-optimisers for reward
design~\cite{ma2023eureka} and by automatic curriculum learning, is that an LLM can
\emph{wean} the agent by progressively loosening the shield as the agent improves. A single
scalar \emph{permissiveness} \(p\in[0,1]\) controls every tunable threshold \(\tau_i\) of the
adaptive shield by linearly interpolating each one between a strict and a permissive endpoint, so
that \(p=1\) recovers the production shield and \(p=0\) the maximally weaned shield
(Figure~\ref{fig:permissiveness-ramp}; the interpolation is Equation~\eqref{eq:permissiveness} in
Appendix~\ref{app:math}).

\begin{figure}[H]
  \centering
  \includegraphics[width=0.8\textwidth]{imgs/permissiveness_ramp}
  \caption{The permissiveness scalar \(p\) sets every tunable shield threshold \(\tau_i\) at once by
           linear interpolation between its strict value (at \(p=1\), the production shield) and its
           permissive value (at \(p=0\), maximally weaned). Three representative thresholds are
           shown; the LLM meta-controller proposes a single \(p\) every \(100\) episodes and the
           deterministic safety governor accepts or vetoes it.}
  \label{fig:permissiveness-ramp}
\end{figure}

Every \(100\) episodes, a local LLM (Gemma~\cite{team2024gemma} via Ollama, running on CPU so as not to contend with
CARLA for GPU memory) receives a compact snapshot of the safety and competence metrics and
proposes a target permissiveness. Crucially, the LLM never acts directly: a deterministic
\emph{safety governor} mediates every proposal under four hard rules: a safety floor that
forces the shield \emph{stricter} whenever crash or off-road rates breach a ceiling, a ratchet
that only loosens when the metrics are below stricter targets, a bounded per-decision step, and a
dwell time between changes. The governor makes ``safety first'' a property of code rather than of
the language model.

\subsection{Result: loosening is not weaning}

A \(3{,}500\)-episode run with this controller produced three independent and mutually
reinforcing observations.

\begin{enumerate}
  \item \textbf{The LLM added no adaptive intelligence.} It proposed ``loosen'' on all \(34\)
        decisions, irrespective of the crash rate; the deterministic governor performed all of
        the safety regulation (twelve \emph{hold} and two \emph{force-tighten} vetoes). A
        three-line fixed schedule would have behaved identically.
  \item \textbf{Loosening did not reduce reliance.} As Figure~\ref{fig:llm-negative} shows, the
        permissiveness fell from \(1.0\) to \(\approx0.55\) while the shield-activation rate
        remained essentially flat: the Pearson correlation between the two series is
        \(\approx 0\). Meanwhile the correlation between permissiveness and the crash rate is
        \(-0.50\): removing the shield's help immediately surfaced unsafe behaviour, exactly the
        ``boundary-riding'' dynamic of Section~\ref{sec:dep-problem}.
  \item \textbf{Reliance was geometry, not skill.} \(96.7\%\) of all interventions were static
        (guardrail proximity on the Town04 highway), a structural feature of the map that the
        tunable thresholds barely affect.
\end{enumerate}

\begin{figure}[H]
  \centering
  \includegraphics[width=0.85\textwidth]{imgs/llm_negative}
  \caption{LLM meta-controller over \(3{,}500\) episodes. The controller loosens the shield
           (permissiveness, left axis) but the shield-activation rate (right axis) does not fall:
           the agent compensates by driving less precisely, keeping its reliance constant. The
           crash rate, in contrast, rises as the shield is loosened.}
  \label{fig:llm-negative}
\end{figure}

The conclusion is unambiguous: loosening the shield's \emph{thresholds} does not teach the agent
anything; it merely moves the boundary the agent rides into, paying the safety cost of weaning
without buying any independence. The LLM, additionally, never justified its place over a
deterministic schedule.

\section{Second approach: the shield as a teacher}\label{sec:dep-teacher}

\subsection{Filling the gradient hole}

If dependence is caused by the absence of gradient on the shielded states (R2), the remedy is to
supply one. The masked PPO surrogate is left untouched, since it is \emph{correct} not to apply an
importance-weighted update to an off-policy action that the shield substituted, and a
supervised behavioural-cloning (BC) term is \emph{added} on exactly those masked steps, pulling
the policy mean towards the action the shield executed:
\begin{equation}\label{eq:bc}
  \mathcal{L}_{\mathrm{BC}}(\theta) \;=\;
  \frac{\sum_t \alpha_t\,\bigl(\tanh(\mu_\theta(s_t))_{\mathrm{steer}} - a^{\mathrm{sh}}_{t,\mathrm{steer}}\bigr)^2}
       {\sum_t \alpha_t},
  \qquad
  \mathcal{L} \;=\; \mathcal{L}_{\mathrm{PPO}} + \lambda_{\mathrm{BC}}\,\mathcal{L}_{\mathrm{BC}},
\end{equation}
where \(\mu_\theta\) is the actor mean, \(\tanh(\mu_\theta(s_t))_{\mathrm{steer}}\) the policy's
deterministic steering output, \(a^{\mathrm{sh}}_{t,\mathrm{steer}}\) the steering component of the
shield's executed action, and \(\alpha_t\in[0,1]\) the per-step shield mixing coefficient of
Equation~\eqref{eq:alpha-blend}, so that each step is cloned in proportion to how strongly the
shield intervened (\(\alpha_t = 0\) on unshielded steps contributes nothing). The total loss adds
this term to the masked PPO surrogate \(\mathcal{L}_{\mathrm{PPO}}\) of
Section~\ref{sec:design-agent-ppo} with weight \(\lambda_{\mathrm{BC}}\ge0\).
The resulting learning rule is clean: PPO on the on-policy steps, behavioural cloning of the
shield on the shield-corrected steps, and the critic on all steps. The coefficient
\(\lambda_{\mathrm{BC}}\) is annealed to zero so the shield teaches strongly early and the agent
is released to its own policy late, a curriculum analogous to learning from
intervention~\cite{ross2011dagger}. Critically, the cloning target in
Equation~\eqref{eq:bc} is restricted to the \emph{steering} dimension; the reason is analysed in
Section~\ref{sec:dep-collapse}.

\subsection{Result: the agent learns to drive unaided}

The scheme is applied as a warm start from the competent shielded policy. The effect on the
shield-off probe is dramatic (Figure~\ref{fig:probe}): off-road episodes, the dominant unaided
failure, fall from \(\approx90\%\) to \(\approx6\%\), and the shield-off success rate rises from
\(0\%\) to \(50\text{--}75\%\) within a few hundred episodes, where the baseline curve stays flat
at zero. The agent has learned the shield's lane-keeping skill.

\begin{figure}[H]
  \centering
  \includegraphics[width=0.95\textwidth]{imgs/shield_off_probe}
  \caption{Shield-off probe (the true learning curve). Without the teacher scheme (grey) the
           agent never learns to drive unaided: it leaves the road within \(\approx40\,\mathrm{m}\)
           on essentially every episode. With the steering behavioural-cloning teacher (green),
           shield-off success climbs to \(50\text{--}75\%\) and off-road collapses to a few per
           cent.}
  \label{fig:probe}
\end{figure}

A final shield ablation confirms the result on identical, seed-aligned scenarios with \(60\) NPCs
(Table~\ref{tab:dep-ablation} and Figure~\ref{fig:ablation}). The headline number is that the
\emph{shield-off} success rate (\(66\%\)) \emph{matches} the \emph{shield-on} rate (\(66\%\)):
the agent no longer needs the shield to function. What the shield still buys is safety: it
reduces the crash rate from \(16\%\) to \(12\%\) and the off-road rate from \(18\%\) to \(8\%\),
with the combined unsafe rate falling from \(34\%\) to \(20\%\), at the cost of \(12\%\)
stuck episodes induced by the shield's braking. The shield has gone from a crutch the agent
could not function without to a safety net it does not need in order to drive.

\begin{table}[H]
  \centering
  \caption{Shield ablation after weaning (\texttt{learn\_fix2}), \(50\) episodes,
           \(60\) NPCs, seed-aligned scenarios.}
  \label{tab:dep-ablation}
  \small
  \begin{tabular}{@{}lccccc@{}}
    \toprule
    \textbf{Configuration} & \textbf{Success} & \textbf{Crash} & \textbf{Off-road} &
    \textbf{Stuck} & \textbf{Shields/ep} \\
    \midrule
    Shield OFF (\texttt{none})      & \(66\%\) & \(16\%\) & \(18\%\) & \(0\%\) & \(0.0\) \\
    Shield ON (\texttt{adaptive})   & \(66\%\) & \(12\%\) & \(8\%\)  & \(12\%\) & \(65.6\) \\
    \bottomrule
  \end{tabular}
\end{table}

\begin{figure}[H]
  \centering
  \includegraphics[width=\textwidth]{imgs/shield_ablation}
  \caption{Ablation of the weaned agent (\texttt{learn\_fix2}, \(60\) NPCs). Shield-off and
           shield-on success are matched at \(66\%\); the shield's remaining value is reducing
           the combined crash\,+\,off-road rate from \(34\%\) to \(20\%\).}
  \label{fig:ablation}
\end{figure}

Figure~\ref{fig:weaned-comparison} contextualises this result against the pre-weaning
baseline (\texttt{FullSendC1\_p4}) across all three shield configurations under identical
conditions (20~NPCs, seed-aligned scenarios).

\begin{figure}[H]
  \centering
  \includegraphics[width=\textwidth]{imgs/weaned_vs_baseline}
  \caption{Baseline (\texttt{FullSendC1\_p4}, top row) versus weaned (\texttt{learn\_fix2},
           bottom row), 50 seed-aligned episodes, 20~NPCs. Left: outcome distribution per
           shield type. Centre: crash and off-road rates. Right: average distance and shield
           interventions per episode. The baseline agent requires the adaptive shield to
           drive at all (0\% success without it); the weaned agent reaches 74\% success
           with no shield and matches that rate under all three configurations.}
  \label{fig:weaned-comparison}
\end{figure}

\section{Failure analysis: why only the steering}\label{sec:dep-collapse}

Restricting the cloning target in Equation~\eqref{eq:bc} to the steering dimension is not a
detail but the difference between success and collapse. When the throttle/brake dimension is
included, the policy collapses to a stationary ``parking'' equilibrium
(Figure~\ref{fig:collapse}): the mean speed falls to near zero and the success rate to zero.

The mechanism is that \emph{the shield's throttle output is always an emergency brake}, as seen in
the lateral-recovery and emergency-brake commands of Section~\ref{sec:design-adaptive-shield}. Unlike
its steering, which demonstrates good lane keeping, the shield never demonstrates good
\emph{longitudinal} driving; it only ever demonstrates slamming the brakes. Cloning that signal
teaches the agent to brake, and because behavioural cloning is a \emph{direct} supervised
gradient on the policy mean it overpowers the reward-mediated incentives that would otherwise
keep the agent moving. A narrower variant that cloned the brake only on \emph{dynamic} (vehicle)
interventions was also tried; it too collapsed, as soon as the curriculum raised the traffic
density the now-frequent vehicle interventions taught systematic over-braking
(Figure~\ref{fig:collapse}). The lesson generalises: \textbf{the shield's steering is a good
teacher, but its throttle is never one}, because it encodes a last-resort reflex rather than a
driving policy.

\begin{figure}[H]
  \centering
  \includegraphics[width=0.85\textwidth]{imgs/throttle_collapse}
  \caption{Imitating the shield's throttle collapses the policy to parking. Cloning the steering
           only (green) is stable; cloning the throttle on every shielded step (orange) or only on
           dynamic interventions (red) drives the mean speed to zero, the latter precisely when
           the curriculum increases the traffic density.}
  \label{fig:collapse}
\end{figure}

\section{Discussion}\label{sec:dep-discussion}

The two approaches share a target but differ in what they treat as the cause of dependence. The
LLM meta-controller treats it as a \emph{scheduling} problem, asking how aggressively the
shield should filter, and the experiment of Section~\ref{sec:dep-llm} shows that this is the wrong
variable: loosening the filter neither teaches the agent nor reduces its reliance, it only
exposes it to more risk. The behavioural-cloning teacher treats dependence as a \emph{learning}
problem, since the agent never received gradient on the states the shield handled, and resolves
it by supplying that gradient. Equation~\eqref{eq:bc} is, in effect, the missing half of the
masked objective.

Two corollaries are worth recording. First, the safety governor of
Section~\ref{sec:dep-llm} is a sound and reusable component independently of the LLM: it never
permitted an unsafe run, and the ``proposer + deterministic veto'' pattern is a safe way to put a
learned controller in a safety-critical loop even when, as here, the proposer turns out to be
weak. Second, the residual gap after weaning is collision avoidance (the unsafe rate
remains at \(34\%\) without the shield against \(20\%\) with it); since the shield is a poor
longitudinal
teacher, closing that gap is better posed as a reward-shaping problem (a time-to-collision or
following-distance penalty that lets the agent learn its \emph{own} gentle avoidance) than as a
further imitation target, and is left as future work.

\clearpage

% DRAFT scaffold -- uncomment once imgs/ figures exist and referencias.bib has the new keys
% (see docs/shielded_motif_design.md §8, §10):
% ============================================================
%   Chapter 8  \textemdash\  Learning Collision Avoidance from
%                            LLM-Derived Rewards
% ------------------------------------------------------------
%   DRAFT SCAFFOLD (design-first). Not yet wired into main.tex.
%   To include: uncomment the % ============================================================
%   Chapter 8  \textemdash\  Learning Collision Avoidance from
%                            LLM-Derived Rewards
% ------------------------------------------------------------
%   DRAFT SCAFFOLD (design-first). Not yet wired into main.tex.
%   To include: uncomment the % ============================================================
%   Chapter 8  \textemdash\  Learning Collision Avoidance from
%                            LLM-Derived Rewards
% ------------------------------------------------------------
%   DRAFT SCAFFOLD (design-first). Not yet wired into main.tex.
%   To include: uncomment the \input{tex/08_llm_reward_exposure}
%   line in main.tex AFTER the figures listed below exist in
%   Memoria/imgs/ and the bib keys in §10 of
%   docs/shielded_motif_design.md are added to referencias.bib.
%
%   Planned figures (imgs/, generate from run SQLite as in Ch.7):
%     - encounter_rate_curve   (exposure director lifts encounter_rate ~3% -> ~30%)
%     - reward_field           (r_phi over front-range x closing-speed slice)
%     - arms_crash_bars        (C0/C1/C2/C3 shield-OFF crash @60 NPC)
%     - shaping_learning_curve (shield-OFF crash vs updates, per arm)
%   Planned tables:
%     - tab:cov-arms           (factorial arms x outcomes)
% ============================================================

\chapter{Learning Collision Avoidance from LLM-Derived Rewards}\label{ch:llmreward}

Chapter~\ref{ch:dependence} closed on an explicit open problem. Behavioural cloning of the
shield's steering removed shield \emph{dependence} \textemdash\ the weaned agent drives as well
unaided as it does shielded \textemdash\ but it left a residual \emph{collision} gap: without the
shield the agent still crashes on roughly one episode in six, because the shield's throttle is an
emergency brake and is therefore a poor longitudinal teacher (Section~\ref{sec:dep-collapse}).
That chapter argued the gap is ``better posed as a reward-shaping problem \dots\ that lets the
agent learn its \emph{own} gentle avoidance''. This chapter takes up that problem, and in doing so
revisits \textemdash\ and finally vindicates \textemdash\ the use of a Large Language Model in the
training loop, after the negative result of Section~\ref{sec:dep-llm}.

The chapter is organised around a single thesis: the earlier LLM failure and the failure of a
hand-designed time-to-collision (TTC) penalty share \emph{one} root cause, and that cause dictates
the design that finally works. Section~\ref{sec:cov-diagnosis} establishes the diagnosis;
Section~\ref{sec:cov-method} derives the method (an LLM-learned dense reward trained on an
LLM-generated exposure curriculum); Sections~\ref{sec:cov-exposure}--\ref{sec:cov-shaping} detail
its three components; Section~\ref{sec:cov-results} reports the experiments; and
Section~\ref{sec:cov-discussion} places the result in the \emph{experience\,>\,gradient\,>\,
incentive} hierarchy that organises the whole of Chapters~\ref{ch:dependence}--\ref{ch:llmreward}.

\section{Diagnosis: one root cause behind two negative results}\label{sec:cov-diagnosis}

Two attempts to use a language model or a designed reward to close the avoidance gap failed, and
they failed identically.

\paragraph{The permissiveness dial (Section~\ref{sec:dep-llm}).} The LLM controlled a single
scalar every hundred episodes; its output was \emph{not causally coupled} to the metric it was
meant to move (the correlation between permissiveness and shield-activation rate was
\(\approx 0\)). A fixed schedule would have done the same.

\paragraph{The TTC penalty.} A hand-designed time-to-collision / following-distance penalty was
added to the reward shaper and its weight swept across four runs. The penalty \emph{fired in only
\(\approx 3\%\) of episodes}, its magnitude when fired was flat across weights, and the shield-off
crash rate did not fall \textemdash\ it merely redistributed failures from off-road to collision.
% TODO: cite/forward-ref the TTC sweep table if it is written up elsewhere.

\paragraph{The common cause.} Both signals were starved. On the Town04 highway loop, globally
spawned traffic places a vehicle in the ego's front cone in only about \(3\%\) of episodes, so any
collision-specific signal \textemdash\ whether an LLM dial or a TTC penalty \textemdash\ almost
never has anything to act on. Three constraints follow, and the method of
Section~\ref{sec:cov-method} is the minimal design that satisfies all three:
\begin{enumerate}
  \item \textbf{Bandwidth.} The learning signal must be \emph{dense} (per step), like the
        behavioural-cloning term of Equation~\eqref{eq:bc}, not a per-window scalar.
  \item \textbf{Exposure.} A collision signal starves at \(3\%\) coverage; the relevant
        experience must be \emph{manufactured}.
  \item \textbf{Stationarity.} A live additive reward rescales the running return
        normaliser the critic is calibrated against (Section~\ref{sec:design-agent-ppo}); the
        signal must be \emph{policy-invariant}.
\end{enumerate}

\section{Method: dense LLM reward over generated exposure}\label{sec:cov-method}

The proposed method, which we call \emph{Shielded Motif under generated exposure}, has three
components that map one-to-one onto the three constraints. An \textbf{exposure director}
(Section~\ref{sec:cov-exposure}) manufactures vehicle encounters, supplying the missing
experience. A \textbf{preference-trained reward model} (Section~\ref{sec:cov-reward}) supplies a
dense, per-step safety signal learned from a language model's judgements \textemdash\ offline,
so the model never runs inside the control loop. A \textbf{potential-based shaping}
integration (Section~\ref{sec:cov-shaping}) delivers that signal without altering the optimal
policy. Throughout, the safety shield and the deterministic safety governor of
Chapter~\ref{ch:dependence} are retained unchanged: the shield still prevents unsafe
\emph{exploration} on the harder, manufactured scenarios, and the ``proposer~+~deterministic
veto'' pattern \textemdash\ already shown to be sound independently of the LLM
(Section~\ref{sec:dep-discussion}) \textemdash\ now guards the exposure curriculum.

The design follows the Motif paradigm of intrinsic motivation from artificial-intelligence
feedback~\cite{klissarov2023motif}, in which an LLM expresses preferences over pairs of
observation captions that are distilled into a reward model; it inherits the preference-learning
formulation of~\cite{christiano2017preferences} and the policy-invariance guarantee
of~\cite{ng1999shaping}. It is related to, but distinct from, language-conditioned reward
generation for driving~\cite{huang2025vlmrl} and to LLM-based safety-critical scenario
generation~\cite{zhang2024chatscene,nie2025seeking}, from which the exposure director borrows its
intent-conditioned design.

\section{Component I \textemdash\ the exposure director}\label{sec:cov-exposure}

% Narrative to write:
%  - route-traffic injection along the ego waypoint chain (leads + cut-ins),
%    absolute desired speeds (not relative to the speed limit), governor-gated.
%  - failure-driven scenario generation: shield-off crash traces -> LLM infers a
%    behavioural intent -> structured scenario spec; intent->recipe memory bank
%    (after nie2025seeking, zhang2024chatscene).
%  - progress KPI: encounter_rate = fraction of episodes with min_front_dynamic < 0.7.
%  - safety: a dedicated governor with exposure ceilings (crash 0.20 / target 0.12)
%    that CUTS traffic on a crash breach.
%  - figure: encounter_rate_curve.

\section{Component II \textemdash\ the preference-trained reward model}\label{sec:cov-reward}

% Narrative to write:
%  - offline pipeline: sample observations (focus on the dynamic-LIDAR channel
%    obs[240:480] + front-cone scalars), template each into a structured caption,
%    LLM ranks PAIRS into preferences (reuse the schema-constrained, deterministic,
%    fallback-guarded decoding of the abandoned shield tuner).
%  - reward model r_phi: MLP 240(+k)->64->64->1; Bradley-Terry loss below.
%  - DISTINCTION from a crash-label safety critic: dense graded preference signal
%    available on near-miss steps that never crashed, vs sparse binary crash labels.

The reward model \(r_\phi\) is trained from preferences with the Bradley--Terry
objective~\cite{christiano2017preferences}: for an annotated pair of scenes \((s_a, s_b)\) in
which the language model judges \(s_a\) the safer situation,
\begin{equation}\label{eq:bt}
  \mathcal{L}(\phi) \;=\; -\,\mathbb{E}_{(s_a,s_b)}\;
  c_{ab}\,\log \sigma\!\bigl(r_\phi(s_a) - r_\phi(s_b)\bigr),
\end{equation}
where \(\sigma\) is the logistic function and \(c_{ab}\in[0,1]\) is the language model's stated
confidence, used to down-weight uncertain pairs. The annotation is performed entirely offline, so
the language model never contends with CARLA for GPU memory and adds no latency to the \(20\,
\mathrm{Hz}\) control loop \textemdash\ directly removing the operational difficulty that the
in-loop controller of Section~\ref{sec:dep-llm} suffered.

\section{Component III \textemdash\ policy-invariant shaping}\label{sec:cov-shaping}

The learned reward is delivered as potential-based shaping~\cite{ng1999shaping}, using \(r_\phi\)
as the potential:
\begin{equation}\label{eq:potential}
  F_t \;=\; \gamma\,r_\phi(s_{t+1}) - r_\phi(s_t),
  \qquad
  r^{\mathrm{shaped}}_t \;=\; r_t \;+\; \beta\,F_t .
\end{equation}
Because \(F_t\) is a difference of potentials, the set of optimal policies is provably
unchanged~\cite{ng1999shaping}: the term can only \emph{re-time} credit toward the states that
precede danger, it cannot create a new optimum. This is the formal guarantee that the
park-attractor of Section~\ref{sec:dep-collapse} \textemdash\ in which a direct supervised
gradient overpowered the reward and drove the policy to a stationary equilibrium \textemdash\
cannot recur here, and it is the reason the design respects the reward-stationarity constraint
that the original meta-controller was deliberately built to honour. The coefficient \(\beta\) is
frozen per run and annealed toward a small floor, strong early to install the avoidance skill.

\section{Experiments}\label{sec:cov-results}

% Pre-registered factorial design (see docs/shielded_motif_design.md, table tab:cov-arms):
%  C0 control (= learn_fix2), C1 exposure-only, C2 reward-only (predicted to STARVE),
%  C3 full. All warm-start learn_fix2 with --log_std_anchor_coef 1e-3 (mandatory).
%  Primary metric: shield-OFF crash @60 NPC (main_eval --num_npc 60 --seed 100).
%  Success: C3 shield-OFF crash < 16% (the learn_fix2 control).
%  Decisive/falsification test: if exposure lifts encounter_rate to ~30% and crash
%  still will not move, that is a clean, causally-isolated third negative result.

\begin{table}[H]
  \centering
  \caption{Pre-registered factorial arms. Each isolates one leg of the method; C2 is predicted
           to \emph{starve} (a reward with no exposure), which would itself confirm the
           exposure diagnosis of Section~\ref{sec:cov-diagnosis}.}
  \label{tab:cov-arms}
  \small
  \begin{tabular}{@{}llll@{}}
    \toprule
    \textbf{Arm} & \textbf{Exposure} & \textbf{Pref.\ reward} & \textbf{Purpose} \\
    \midrule
    C0 control       & global NPC only    & off & reproduce \texttt{learn\_fix2} \\
    C1 exposure-only & LLM director       & off & does exposure alone move crash? \\
    C2 reward-only   & global NPC only    & on  & predicted to starve (3\% coverage) \\
    C3 full          & LLM director       & on  & the proposed method \\
    \bottomrule
  \end{tabular}
\end{table}

% TODO figures once runs land:
% \begin{figure}[H]\centering\includegraphics[width=0.85\textwidth]{imgs/arms_crash_bars}
%   \caption{Shield-off crash rate at 60 NPCs across the four arms.}\label{fig:cov-arms}\end{figure}

\section{Discussion}\label{sec:cov-discussion}

% Narrative to write:
%  - the experience > gradient > incentive hierarchy: TTC (incentive) failed for lack of
%    exposure; shield BC (gradient) fixed steering but not throttle; only when an LLM
%    manufactures the experience does a learned dense reward teach self-driven avoidance.
%  - salvage story: the abandoned meta-stack (schema-constrained decoding, deterministic
%    governor, async runner) is reused as trustworthy offline annotation infrastructure;
%    the negative result of Ch.7 is what FORCED this (correct) design, not a lucky guess.
%  - honest limitation: if the decisive test fails, report the bounded negative result.

\clearpage

%   line in main.tex AFTER the figures listed below exist in
%   Memoria/imgs/ and the bib keys in §10 of
%   docs/shielded_motif_design.md are added to referencias.bib.
%
%   Planned figures (imgs/, generate from run SQLite as in Ch.7):
%     - encounter_rate_curve   (exposure director lifts encounter_rate ~3% -> ~30%)
%     - reward_field           (r_phi over front-range x closing-speed slice)
%     - arms_crash_bars        (C0/C1/C2/C3 shield-OFF crash @60 NPC)
%     - shaping_learning_curve (shield-OFF crash vs updates, per arm)
%   Planned tables:
%     - tab:cov-arms           (factorial arms x outcomes)
% ============================================================

\chapter{Learning Collision Avoidance from LLM-Derived Rewards}\label{ch:llmreward}

Chapter~\ref{ch:dependence} closed on an explicit open problem. Behavioural cloning of the
shield's steering removed shield \emph{dependence} \textemdash\ the weaned agent drives as well
unaided as it does shielded \textemdash\ but it left a residual \emph{collision} gap: without the
shield the agent still crashes on roughly one episode in six, because the shield's throttle is an
emergency brake and is therefore a poor longitudinal teacher (Section~\ref{sec:dep-collapse}).
That chapter argued the gap is ``better posed as a reward-shaping problem \dots\ that lets the
agent learn its \emph{own} gentle avoidance''. This chapter takes up that problem, and in doing so
revisits \textemdash\ and finally vindicates \textemdash\ the use of a Large Language Model in the
training loop, after the negative result of Section~\ref{sec:dep-llm}.

The chapter is organised around a single thesis: the earlier LLM failure and the failure of a
hand-designed time-to-collision (TTC) penalty share \emph{one} root cause, and that cause dictates
the design that finally works. Section~\ref{sec:cov-diagnosis} establishes the diagnosis;
Section~\ref{sec:cov-method} derives the method (an LLM-learned dense reward trained on an
LLM-generated exposure curriculum); Sections~\ref{sec:cov-exposure}--\ref{sec:cov-shaping} detail
its three components; Section~\ref{sec:cov-results} reports the experiments; and
Section~\ref{sec:cov-discussion} places the result in the \emph{experience\,>\,gradient\,>\,
incentive} hierarchy that organises the whole of Chapters~\ref{ch:dependence}--\ref{ch:llmreward}.

\section{Diagnosis: one root cause behind two negative results}\label{sec:cov-diagnosis}

Two attempts to use a language model or a designed reward to close the avoidance gap failed, and
they failed identically.

\paragraph{The permissiveness dial (Section~\ref{sec:dep-llm}).} The LLM controlled a single
scalar every hundred episodes; its output was \emph{not causally coupled} to the metric it was
meant to move (the correlation between permissiveness and shield-activation rate was
\(\approx 0\)). A fixed schedule would have done the same.

\paragraph{The TTC penalty.} A hand-designed time-to-collision / following-distance penalty was
added to the reward shaper and its weight swept across four runs. The penalty \emph{fired in only
\(\approx 3\%\) of episodes}, its magnitude when fired was flat across weights, and the shield-off
crash rate did not fall \textemdash\ it merely redistributed failures from off-road to collision.
% TODO: cite/forward-ref the TTC sweep table if it is written up elsewhere.

\paragraph{The common cause.} Both signals were starved. On the Town04 highway loop, globally
spawned traffic places a vehicle in the ego's front cone in only about \(3\%\) of episodes, so any
collision-specific signal \textemdash\ whether an LLM dial or a TTC penalty \textemdash\ almost
never has anything to act on. Three constraints follow, and the method of
Section~\ref{sec:cov-method} is the minimal design that satisfies all three:
\begin{enumerate}
  \item \textbf{Bandwidth.} The learning signal must be \emph{dense} (per step), like the
        behavioural-cloning term of Equation~\eqref{eq:bc}, not a per-window scalar.
  \item \textbf{Exposure.} A collision signal starves at \(3\%\) coverage; the relevant
        experience must be \emph{manufactured}.
  \item \textbf{Stationarity.} A live additive reward rescales the running return
        normaliser the critic is calibrated against (Section~\ref{sec:design-agent-ppo}); the
        signal must be \emph{policy-invariant}.
\end{enumerate}

\section{Method: dense LLM reward over generated exposure}\label{sec:cov-method}

The proposed method, which we call \emph{Shielded Motif under generated exposure}, has three
components that map one-to-one onto the three constraints. An \textbf{exposure director}
(Section~\ref{sec:cov-exposure}) manufactures vehicle encounters, supplying the missing
experience. A \textbf{preference-trained reward model} (Section~\ref{sec:cov-reward}) supplies a
dense, per-step safety signal learned from a language model's judgements \textemdash\ offline,
so the model never runs inside the control loop. A \textbf{potential-based shaping}
integration (Section~\ref{sec:cov-shaping}) delivers that signal without altering the optimal
policy. Throughout, the safety shield and the deterministic safety governor of
Chapter~\ref{ch:dependence} are retained unchanged: the shield still prevents unsafe
\emph{exploration} on the harder, manufactured scenarios, and the ``proposer~+~deterministic
veto'' pattern \textemdash\ already shown to be sound independently of the LLM
(Section~\ref{sec:dep-discussion}) \textemdash\ now guards the exposure curriculum.

The design follows the Motif paradigm of intrinsic motivation from artificial-intelligence
feedback~\cite{klissarov2023motif}, in which an LLM expresses preferences over pairs of
observation captions that are distilled into a reward model; it inherits the preference-learning
formulation of~\cite{christiano2017preferences} and the policy-invariance guarantee
of~\cite{ng1999shaping}. It is related to, but distinct from, language-conditioned reward
generation for driving~\cite{huang2025vlmrl} and to LLM-based safety-critical scenario
generation~\cite{zhang2024chatscene,nie2025seeking}, from which the exposure director borrows its
intent-conditioned design.

\section{Component I \textemdash\ the exposure director}\label{sec:cov-exposure}

% Narrative to write:
%  - route-traffic injection along the ego waypoint chain (leads + cut-ins),
%    absolute desired speeds (not relative to the speed limit), governor-gated.
%  - failure-driven scenario generation: shield-off crash traces -> LLM infers a
%    behavioural intent -> structured scenario spec; intent->recipe memory bank
%    (after nie2025seeking, zhang2024chatscene).
%  - progress KPI: encounter_rate = fraction of episodes with min_front_dynamic < 0.7.
%  - safety: a dedicated governor with exposure ceilings (crash 0.20 / target 0.12)
%    that CUTS traffic on a crash breach.
%  - figure: encounter_rate_curve.

\section{Component II \textemdash\ the preference-trained reward model}\label{sec:cov-reward}

% Narrative to write:
%  - offline pipeline: sample observations (focus on the dynamic-LIDAR channel
%    obs[240:480] + front-cone scalars), template each into a structured caption,
%    LLM ranks PAIRS into preferences (reuse the schema-constrained, deterministic,
%    fallback-guarded decoding of the abandoned shield tuner).
%  - reward model r_phi: MLP 240(+k)->64->64->1; Bradley-Terry loss below.
%  - DISTINCTION from a crash-label safety critic: dense graded preference signal
%    available on near-miss steps that never crashed, vs sparse binary crash labels.

The reward model \(r_\phi\) is trained from preferences with the Bradley--Terry
objective~\cite{christiano2017preferences}: for an annotated pair of scenes \((s_a, s_b)\) in
which the language model judges \(s_a\) the safer situation,
\begin{equation}\label{eq:bt}
  \mathcal{L}(\phi) \;=\; -\,\mathbb{E}_{(s_a,s_b)}\;
  c_{ab}\,\log \sigma\!\bigl(r_\phi(s_a) - r_\phi(s_b)\bigr),
\end{equation}
where \(\sigma\) is the logistic function and \(c_{ab}\in[0,1]\) is the language model's stated
confidence, used to down-weight uncertain pairs. The annotation is performed entirely offline, so
the language model never contends with CARLA for GPU memory and adds no latency to the \(20\,
\mathrm{Hz}\) control loop \textemdash\ directly removing the operational difficulty that the
in-loop controller of Section~\ref{sec:dep-llm} suffered.

\section{Component III \textemdash\ policy-invariant shaping}\label{sec:cov-shaping}

The learned reward is delivered as potential-based shaping~\cite{ng1999shaping}, using \(r_\phi\)
as the potential:
\begin{equation}\label{eq:potential}
  F_t \;=\; \gamma\,r_\phi(s_{t+1}) - r_\phi(s_t),
  \qquad
  r^{\mathrm{shaped}}_t \;=\; r_t \;+\; \beta\,F_t .
\end{equation}
Because \(F_t\) is a difference of potentials, the set of optimal policies is provably
unchanged~\cite{ng1999shaping}: the term can only \emph{re-time} credit toward the states that
precede danger, it cannot create a new optimum. This is the formal guarantee that the
park-attractor of Section~\ref{sec:dep-collapse} \textemdash\ in which a direct supervised
gradient overpowered the reward and drove the policy to a stationary equilibrium \textemdash\
cannot recur here, and it is the reason the design respects the reward-stationarity constraint
that the original meta-controller was deliberately built to honour. The coefficient \(\beta\) is
frozen per run and annealed toward a small floor, strong early to install the avoidance skill.

\section{Experiments}\label{sec:cov-results}

% Pre-registered factorial design (see docs/shielded_motif_design.md, table tab:cov-arms):
%  C0 control (= learn_fix2), C1 exposure-only, C2 reward-only (predicted to STARVE),
%  C3 full. All warm-start learn_fix2 with --log_std_anchor_coef 1e-3 (mandatory).
%  Primary metric: shield-OFF crash @60 NPC (main_eval --num_npc 60 --seed 100).
%  Success: C3 shield-OFF crash < 16% (the learn_fix2 control).
%  Decisive/falsification test: if exposure lifts encounter_rate to ~30% and crash
%  still will not move, that is a clean, causally-isolated third negative result.

\begin{table}[H]
  \centering
  \caption{Pre-registered factorial arms. Each isolates one leg of the method; C2 is predicted
           to \emph{starve} (a reward with no exposure), which would itself confirm the
           exposure diagnosis of Section~\ref{sec:cov-diagnosis}.}
  \label{tab:cov-arms}
  \small
  \begin{tabular}{@{}llll@{}}
    \toprule
    \textbf{Arm} & \textbf{Exposure} & \textbf{Pref.\ reward} & \textbf{Purpose} \\
    \midrule
    C0 control       & global NPC only    & off & reproduce \texttt{learn\_fix2} \\
    C1 exposure-only & LLM director       & off & does exposure alone move crash? \\
    C2 reward-only   & global NPC only    & on  & predicted to starve (3\% coverage) \\
    C3 full          & LLM director       & on  & the proposed method \\
    \bottomrule
  \end{tabular}
\end{table}

% TODO figures once runs land:
% \begin{figure}[H]\centering\includegraphics[width=0.85\textwidth]{imgs/arms_crash_bars}
%   \caption{Shield-off crash rate at 60 NPCs across the four arms.}\label{fig:cov-arms}\end{figure}

\section{Discussion}\label{sec:cov-discussion}

% Narrative to write:
%  - the experience > gradient > incentive hierarchy: TTC (incentive) failed for lack of
%    exposure; shield BC (gradient) fixed steering but not throttle; only when an LLM
%    manufactures the experience does a learned dense reward teach self-driven avoidance.
%  - salvage story: the abandoned meta-stack (schema-constrained decoding, deterministic
%    governor, async runner) is reused as trustworthy offline annotation infrastructure;
%    the negative result of Ch.7 is what FORCED this (correct) design, not a lucky guess.
%  - honest limitation: if the decisive test fails, report the bounded negative result.

\clearpage

%   line in main.tex AFTER the figures listed below exist in
%   Memoria/imgs/ and the bib keys in §10 of
%   docs/shielded_motif_design.md are added to referencias.bib.
%
%   Planned figures (imgs/, generate from run SQLite as in Ch.7):
%     - encounter_rate_curve   (exposure director lifts encounter_rate ~3% -> ~30%)
%     - reward_field           (r_phi over front-range x closing-speed slice)
%     - arms_crash_bars        (C0/C1/C2/C3 shield-OFF crash @60 NPC)
%     - shaping_learning_curve (shield-OFF crash vs updates, per arm)
%   Planned tables:
%     - tab:cov-arms           (factorial arms x outcomes)
% ============================================================

\chapter{Learning Collision Avoidance from LLM-Derived Rewards}\label{ch:llmreward}

Chapter~\ref{ch:dependence} closed on an explicit open problem. Behavioural cloning of the
shield's steering removed shield \emph{dependence} \textemdash\ the weaned agent drives as well
unaided as it does shielded \textemdash\ but it left a residual \emph{collision} gap: without the
shield the agent still crashes on roughly one episode in six, because the shield's throttle is an
emergency brake and is therefore a poor longitudinal teacher (Section~\ref{sec:dep-collapse}).
That chapter argued the gap is ``better posed as a reward-shaping problem \dots\ that lets the
agent learn its \emph{own} gentle avoidance''. This chapter takes up that problem, and in doing so
revisits \textemdash\ and finally vindicates \textemdash\ the use of a Large Language Model in the
training loop, after the negative result of Section~\ref{sec:dep-llm}.

The chapter is organised around a single thesis: the earlier LLM failure and the failure of a
hand-designed time-to-collision (TTC) penalty share \emph{one} root cause, and that cause dictates
the design that finally works. Section~\ref{sec:cov-diagnosis} establishes the diagnosis;
Section~\ref{sec:cov-method} derives the method (an LLM-learned dense reward trained on an
LLM-generated exposure curriculum); Sections~\ref{sec:cov-exposure}--\ref{sec:cov-shaping} detail
its three components; Section~\ref{sec:cov-results} reports the experiments; and
Section~\ref{sec:cov-discussion} places the result in the \emph{experience\,>\,gradient\,>\,
incentive} hierarchy that organises the whole of Chapters~\ref{ch:dependence}--\ref{ch:llmreward}.

\section{Diagnosis: one root cause behind two negative results}\label{sec:cov-diagnosis}

Two attempts to use a language model or a designed reward to close the avoidance gap failed, and
they failed identically.

\paragraph{The permissiveness dial (Section~\ref{sec:dep-llm}).} The LLM controlled a single
scalar every hundred episodes; its output was \emph{not causally coupled} to the metric it was
meant to move (the correlation between permissiveness and shield-activation rate was
\(\approx 0\)). A fixed schedule would have done the same.

\paragraph{The TTC penalty.} A hand-designed time-to-collision / following-distance penalty was
added to the reward shaper and its weight swept across four runs. The penalty \emph{fired in only
\(\approx 3\%\) of episodes}, its magnitude when fired was flat across weights, and the shield-off
crash rate did not fall \textemdash\ it merely redistributed failures from off-road to collision.
% TODO: cite/forward-ref the TTC sweep table if it is written up elsewhere.

\paragraph{The common cause.} Both signals were starved. On the Town04 highway loop, globally
spawned traffic places a vehicle in the ego's front cone in only about \(3\%\) of episodes, so any
collision-specific signal \textemdash\ whether an LLM dial or a TTC penalty \textemdash\ almost
never has anything to act on. Three constraints follow, and the method of
Section~\ref{sec:cov-method} is the minimal design that satisfies all three:
\begin{enumerate}
  \item \textbf{Bandwidth.} The learning signal must be \emph{dense} (per step), like the
        behavioural-cloning term of Equation~\eqref{eq:bc}, not a per-window scalar.
  \item \textbf{Exposure.} A collision signal starves at \(3\%\) coverage; the relevant
        experience must be \emph{manufactured}.
  \item \textbf{Stationarity.} A live additive reward rescales the running return
        normaliser the critic is calibrated against (Section~\ref{sec:design-agent-ppo}); the
        signal must be \emph{policy-invariant}.
\end{enumerate}

\section{Method: dense LLM reward over generated exposure}\label{sec:cov-method}

The proposed method, which we call \emph{Shielded Motif under generated exposure}, has three
components that map one-to-one onto the three constraints. An \textbf{exposure director}
(Section~\ref{sec:cov-exposure}) manufactures vehicle encounters, supplying the missing
experience. A \textbf{preference-trained reward model} (Section~\ref{sec:cov-reward}) supplies a
dense, per-step safety signal learned from a language model's judgements \textemdash\ offline,
so the model never runs inside the control loop. A \textbf{potential-based shaping}
integration (Section~\ref{sec:cov-shaping}) delivers that signal without altering the optimal
policy. Throughout, the safety shield and the deterministic safety governor of
Chapter~\ref{ch:dependence} are retained unchanged: the shield still prevents unsafe
\emph{exploration} on the harder, manufactured scenarios, and the ``proposer~+~deterministic
veto'' pattern \textemdash\ already shown to be sound independently of the LLM
(Section~\ref{sec:dep-discussion}) \textemdash\ now guards the exposure curriculum.

The design follows the Motif paradigm of intrinsic motivation from artificial-intelligence
feedback~\cite{klissarov2023motif}, in which an LLM expresses preferences over pairs of
observation captions that are distilled into a reward model; it inherits the preference-learning
formulation of~\cite{christiano2017preferences} and the policy-invariance guarantee
of~\cite{ng1999shaping}. It is related to, but distinct from, language-conditioned reward
generation for driving~\cite{huang2025vlmrl} and to LLM-based safety-critical scenario
generation~\cite{zhang2024chatscene,nie2025seeking}, from which the exposure director borrows its
intent-conditioned design.

\section{Component I \textemdash\ the exposure director}\label{sec:cov-exposure}

% Narrative to write:
%  - route-traffic injection along the ego waypoint chain (leads + cut-ins),
%    absolute desired speeds (not relative to the speed limit), governor-gated.
%  - failure-driven scenario generation: shield-off crash traces -> LLM infers a
%    behavioural intent -> structured scenario spec; intent->recipe memory bank
%    (after nie2025seeking, zhang2024chatscene).
%  - progress KPI: encounter_rate = fraction of episodes with min_front_dynamic < 0.7.
%  - safety: a dedicated governor with exposure ceilings (crash 0.20 / target 0.12)
%    that CUTS traffic on a crash breach.
%  - figure: encounter_rate_curve.

\section{Component II \textemdash\ the preference-trained reward model}\label{sec:cov-reward}

% Narrative to write:
%  - offline pipeline: sample observations (focus on the dynamic-LIDAR channel
%    obs[240:480] + front-cone scalars), template each into a structured caption,
%    LLM ranks PAIRS into preferences (reuse the schema-constrained, deterministic,
%    fallback-guarded decoding of the abandoned shield tuner).
%  - reward model r_phi: MLP 240(+k)->64->64->1; Bradley-Terry loss below.
%  - DISTINCTION from a crash-label safety critic: dense graded preference signal
%    available on near-miss steps that never crashed, vs sparse binary crash labels.

The reward model \(r_\phi\) is trained from preferences with the Bradley--Terry
objective~\cite{christiano2017preferences}: for an annotated pair of scenes \((s_a, s_b)\) in
which the language model judges \(s_a\) the safer situation,
\begin{equation}\label{eq:bt}
  \mathcal{L}(\phi) \;=\; -\,\mathbb{E}_{(s_a,s_b)}\;
  c_{ab}\,\log \sigma\!\bigl(r_\phi(s_a) - r_\phi(s_b)\bigr),
\end{equation}
where \(\sigma\) is the logistic function and \(c_{ab}\in[0,1]\) is the language model's stated
confidence, used to down-weight uncertain pairs. The annotation is performed entirely offline, so
the language model never contends with CARLA for GPU memory and adds no latency to the \(20\,
\mathrm{Hz}\) control loop \textemdash\ directly removing the operational difficulty that the
in-loop controller of Section~\ref{sec:dep-llm} suffered.

\section{Component III \textemdash\ policy-invariant shaping}\label{sec:cov-shaping}

The learned reward is delivered as potential-based shaping~\cite{ng1999shaping}, using \(r_\phi\)
as the potential:
\begin{equation}\label{eq:potential}
  F_t \;=\; \gamma\,r_\phi(s_{t+1}) - r_\phi(s_t),
  \qquad
  r^{\mathrm{shaped}}_t \;=\; r_t \;+\; \beta\,F_t .
\end{equation}
Because \(F_t\) is a difference of potentials, the set of optimal policies is provably
unchanged~\cite{ng1999shaping}: the term can only \emph{re-time} credit toward the states that
precede danger, it cannot create a new optimum. This is the formal guarantee that the
park-attractor of Section~\ref{sec:dep-collapse} \textemdash\ in which a direct supervised
gradient overpowered the reward and drove the policy to a stationary equilibrium \textemdash\
cannot recur here, and it is the reason the design respects the reward-stationarity constraint
that the original meta-controller was deliberately built to honour. The coefficient \(\beta\) is
frozen per run and annealed toward a small floor, strong early to install the avoidance skill.

\section{Experiments}\label{sec:cov-results}

% Pre-registered factorial design (see docs/shielded_motif_design.md, table tab:cov-arms):
%  C0 control (= learn_fix2), C1 exposure-only, C2 reward-only (predicted to STARVE),
%  C3 full. All warm-start learn_fix2 with --log_std_anchor_coef 1e-3 (mandatory).
%  Primary metric: shield-OFF crash @60 NPC (main_eval --num_npc 60 --seed 100).
%  Success: C3 shield-OFF crash < 16% (the learn_fix2 control).
%  Decisive/falsification test: if exposure lifts encounter_rate to ~30% and crash
%  still will not move, that is a clean, causally-isolated third negative result.

\begin{table}[H]
  \centering
  \caption{Pre-registered factorial arms. Each isolates one leg of the method; C2 is predicted
           to \emph{starve} (a reward with no exposure), which would itself confirm the
           exposure diagnosis of Section~\ref{sec:cov-diagnosis}.}
  \label{tab:cov-arms}
  \small
  \begin{tabular}{@{}llll@{}}
    \toprule
    \textbf{Arm} & \textbf{Exposure} & \textbf{Pref.\ reward} & \textbf{Purpose} \\
    \midrule
    C0 control       & global NPC only    & off & reproduce \texttt{learn\_fix2} \\
    C1 exposure-only & LLM director       & off & does exposure alone move crash? \\
    C2 reward-only   & global NPC only    & on  & predicted to starve (3\% coverage) \\
    C3 full          & LLM director       & on  & the proposed method \\
    \bottomrule
  \end{tabular}
\end{table}

% TODO figures once runs land:
% \begin{figure}[H]\centering\includegraphics[width=0.85\textwidth]{imgs/arms_crash_bars}
%   \caption{Shield-off crash rate at 60 NPCs across the four arms.}\label{fig:cov-arms}\end{figure}

\section{Discussion}\label{sec:cov-discussion}

% Narrative to write:
%  - the experience > gradient > incentive hierarchy: TTC (incentive) failed for lack of
%    exposure; shield BC (gradient) fixed steering but not throttle; only when an LLM
%    manufactures the experience does a learned dense reward teach self-driven avoidance.
%  - salvage story: the abandoned meta-stack (schema-constrained decoding, deterministic
%    governor, async runner) is reused as trustworthy offline annotation infrastructure;
%    the negative result of Ch.7 is what FORCED this (correct) design, not a lucky guess.
%  - honest limitation: if the decisive test fails, report the bounded negative result.

\clearpage

\chapter{Ethical assessment of the project}\label{sec:ethics}

The development of autonomous driving systems cannot be separated from the ethical
weight of the domain in which they operate. A reinforcement learning agent that
controls a vehicle is, in a very literal sense, making decisions that affect human
safety, public infrastructure, and the welfare of other road users. This section
reflects on the ethical issues raised by the design of \textit{CALA\_SafeRL}, contrasts
the option of deploying a standard, unconstrained reinforcement learning agent against
the Shielded-RL approach actually adopted, and justifies the latter as the decision
that best satisfies the engineer's professional and ethical responsibilities. The
discussion closes by situating the project within the broader frameworks of digital
human rights and the United Nations Sustainable Development Goals (SDGs).

\section{Identification of ethical issues}\label{sec:ethics-issues}

Three ethical concerns are inherent to any project that trains a driving policy through
trial and error:

\begin{itemize}
  \item \textbf{Risk to human life and physical integrity.} Reinforcement learning is, by
        construction, an exploratory process: the agent improves by experiencing the
        consequences of its actions, including the negative ones. In a real or
        high-fidelity simulated traffic environment, ``negative consequences'' translate
        into collisions, off-road excursions, or dangerous manoeuvres near pedestrians and
        other vehicles. Even if the immediate setting is a simulator, the policies trained
        in it are intended to generalise toward real-world deployment, so any tolerance for
        unsafe exploratory behaviour during training is, in principle, a tolerance for
        unsafe behaviour that could eventually reach public roads.

  \item \textbf{Algorithmic accountability.} A policy represented as a neural network is
        difficult to audit. When such a system causes harm, it is not always possible to
        trace the decision back to an interpretable cause, which complicates the
        attribution of responsibility among the developer, the operator, and the
        manufacturer. An engineering team that deploys a purely learned controller without
        any verifiable safety layer is, in effect, delegating safety-critical decisions to
        a component whose failure modes cannot be fully characterised in advance.

  \item \textbf{Black-box decision making.} Beyond accountability after the fact, there is
        a forward-looking concern: end users, regulators, and the general public are
        increasingly reluctant to trust systems whose behaviour cannot be bounded or
        explained. A driving agent that "mostly works" but offers no guarantee about its
        worst-case behaviour is difficult to certify and difficult to trust, regardless of
        how well it performs on average.
\end{itemize}

These three issues are not abstract concerns external to the project; they are direct
consequences of the choice of algorithm (PPO, a model-free reinforcement learning
method) and of the choice of action interface (continuous steering and throttle/brake
commands sent directly to the simulated vehicle).

\section{Critical analysis of design options}\label{sec:ethics-options}

Two broad design philosophies were considered during the early stages of the project,
and the trade-off between them is fundamentally an ethical one rather than a purely
technical one.

\paragraph{Option A: an unconstrained reinforcement learning agent.}
The first option is to train the PPO agent directly against the CARLA simulator, with no
intermediate safety layer (the \texttt{--shield\_type none} configuration evaluated in
this work, Section~\ref{ch:experimentation}). This option has clear technical
advantages: the agent's gradient signal is unfiltered, exploration is unconstrained, and
in principle the policy can discover driving strategies that a hand-designed corrector
might not anticipate. From a purely performance-driven perspective, this option could be
argued to maximise \textit{learning speed} and \textit{asymptotic performance}.

However, this option places the entire burden of safety on the learning process itself,
during a phase in which the policy is, by definition, not yet competent. Early in
training the agent behaves close to randomly; under Option A, that randomness is applied
directly to the vehicle's controls, and the resulting collisions and off-road events are
treated merely as numerical penalties in a reward function. There is no mechanism that
prevents an unsafe action from being executed --- safety is, at best, an emergent
property that the agent might converge towards, not a property that is guaranteed at any
point during training or deployment. The evaluation data gathered for this project
(Chapter~\ref{ch:experimentation}) confirms that, even after substantial training, a
non-trivial fraction of episodes under the unshielded configuration end in a crash or an
off-road event, which would be unacceptable if those episodes corresponded to interactions
with real vehicles, pedestrians, or infrastructure.

\paragraph{Option B: Shielded-RL with a deterministic safety layer.}
The second option, and the one adopted in this project, interposes a \textit{safety
shield} between the agent and the simulator. The shield --- in its most developed form,
the \texttt{CarlaAdaptiveHorizonShield} described in Chapter~\ref{ch:design} ---
monitors the agent's proposed action, predicts the vehicle's short-horizon trajectory
using a kinematic bicycle model, and deterministically overrides the action whenever the
prediction indicates an unacceptable risk (collision, lane departure, or deadlock). The
agent still learns from experience, including from the cases in which its action was
corrected, but the action that actually reaches the simulator is bounded by an explicit,
inspectable, rule-based safety envelope.

This option does not maximise raw learning speed or asymptotic reward in the same
unconstrained sense as Option A: the shield occasionally prevents the agent from taking
the action it would have chosen, and the credit-assignment machinery
(Chapter~\ref{ch:design}) must be designed carefully so that the agent is not
penalised for the shield's intervention. What Option B does provide, by construction, is
a \textit{guaranteed upper bound on the severity of unsafe behaviour} at every stage of
training and deployment, together with a fully interpretable component (the shield) whose
decisions can be logged, audited, and explained independently of the neural network.

\paragraph{The ethical trade-off.}
Stated in its simplest form, the choice between Option A and Option B is a choice between
\textit{optimising for performance under an implicit assumption that safety will emerge},
and \textit{optimising for performance under an explicit constraint that safety is
guaranteed by design}. The first option treats safety as a metric to be learned; the
second treats safety as a precondition that is never violated, regardless of what the
learned component has or has not yet learned. From an engineering ethics standpoint,
these are not equivalent options separated by a small performance gap --- they represent
two fundamentally different attitudes toward the question of who, or what, is responsible
for preventing harm.

\section{Justification of the selected decision}\label{sec:ethics-justification}

The Shielded-RL architecture (Option B) was selected, and this decision is justified
along the three axes specified for this assessment.

\paragraph{Engineering professional ethics.}
Professional codes of conduct for engineers --- including those referenced by the
\textit{Colegio de Ingenieros} bodies and by international bodies such as the IEEE and
ACM --- consistently place the protection of public health, safety, and welfare above
all other professional obligations, including those toward employers, clients, or the
advancement of the technology itself. In the context of this project, that principle
translates directly into a design requirement: \textit{no action that the simulated
vehicle executes should be allowed to cause a collision or an off-road excursion if a
safe alternative exists and can be computed in real time}. The
\texttt{CarlaAdaptiveHorizonShield} operationalises this requirement. Its risk
classification (safe / warning / critical, Chapter~\ref{ch:design}) and its
trajectory-prediction logic exist precisely so that the "first do no harm" principle is
enforced at the level of the control loop, not merely hoped for as a statistical outcome
of training. Choosing to build and integrate this component, rather than relying on the
learned policy alone, is a direct application of the duty to prioritise public safety
over project convenience or raw performance metrics.

\paragraph{Social impact and the role of deterministic guardrails in building trust.}
Public acceptance of autonomous systems --- whether autonomous vehicles or, more broadly,
AI systems embedded in everyday infrastructure --- depends heavily on the existence of
\textit{verifiable, deterministic guardrails} that bound the system's worst-case
behaviour. This is the same principle that motivates safety filters and content
guardrails around large language models: a system that is occasionally impressive but
unboundedly unpredictable does not earn public trust, whereas a system whose
worst-case behaviour is provably constrained, even if its average behaviour is more
modest, can be reasoned about, certified, and ultimately trusted. The shield architecture
in this project plays exactly this role for the driving agent. Because the shield's
corrective logic is rule-based and inspectable (kinematic predictions, explicit distance
and time-to-collision thresholds), its behaviour can be explained to a non-specialist
in terms of physical quantities (``the vehicle was predicted to leave its lane within
\(N\) steps, so its steering was corrected''), rather than in terms of opaque network
activations. This interpretability is, in itself, a contribution to the social
acceptability of learned autonomous controllers: it demonstrates that performance and
verifiable safety are not mutually exclusive design goals, and that a learning system
can be embedded inside --- rather than substituted for --- a deterministic safety
framework.

\paragraph{Alignment with Deusto's values and the Sustainable Development Goals.}
This project is also consistent with the broader institutional commitments of the
University of Deusto, in particular the principles articulated in the
\textit{Declaración Deusto de Derechos Humanos en Entornos Digitales}, which calls for
the design of digital and automated systems that actively safeguard human rights ---
including the right to life and physical integrity --- rather than treating their
protection as an afterthought. A driving system whose safety properties are designed in
from the outset, and whose decisions remain auditable, is a direct application of that
principle to the domain of autonomous mobility.

The project's goals also align with several of the United Nations Sustainable
Development Goals:

\begin{itemize}
  \item \textbf{SDG 3 (Good Health and Well-being):} road traffic collisions are a
        leading cause of injury and death worldwide. By developing and evaluating
        mechanisms that reduce the rate of crashes and off-road events during the
        learning process of an autonomous agent (Chapter~\ref{ch:experimentation}), this
        project contributes, even at a research scale, to the broader goal of reducing
        traffic-related harm.

  \item \textbf{SDG 9 (Industry, Innovation, and Infrastructure):} the project advances
        the state of the art in safe reinforcement learning for autonomous driving,
        proposing and evaluating a concrete architecture (the adaptive horizon shield)
        that other researchers and practitioners can build upon, fostering responsible
        innovation in intelligent transport infrastructure.

  \item \textbf{SDG 11 (Sustainable Cities and Communities):} autonomous vehicles are
        widely expected to play a role in future urban mobility systems. Safety
        mechanisms such as the one developed in this project are a prerequisite for the
        responsible integration of such vehicles into shared urban spaces, where
        pedestrians, cyclists, and other vehicles must coexist safely with autonomous
        agents.
\end{itemize}

\paragraph{Summary.}
Taken together, these considerations support the conclusion that the Shielded-RL
architecture is not merely a technical refinement of the baseline PPO agent, but the
ethically appropriate design choice for this project. It reflects the professional duty
to prioritise public safety over unconstrained optimisation, it embodies the kind of
deterministic, interpretable guardrail that is increasingly recognised as a precondition
for public trust in autonomous systems, and it situates the project within the human
rights and sustainability commitments that frame the work of engineers trained at the
University of Deusto.

% ============================================================
%   Chapter 7  \textemdash\  Conclusions and Future Work
% ============================================================

\chapter{Conclusions and Future Work}\label{ch:conclusions}

This chapter closes the document. Section~\ref{sec:concl-contributions} recapitulates the
contributions, Section~\ref{sec:concl-objectives} traces them against the specific objectives
laid out in Chapter~\ref{ch:introduction}, Section~\ref{sec:concl-lessons} reflects on the
lessons learned during the project, and Section~\ref{sec:concl-future} outlines research
directions that go beyond the scope of a Bachelor's thesis.

\section{Summary of contributions}\label{sec:concl-contributions}

The work presented in this document produced the following concrete contributions:

\begin{enumerate}
  \item \textbf{CALA-SafeRL framework.} A modular Safe RL framework for autonomous driving on
        CARLA 0.9.16, organised as a chain of interchangeable \texttt{gymnasium.Wrapper} layers.
        The wrapper chain isolates the PPO agent, the safety shield, the reward shaper and the
        simulator from each other, which makes each component independently testable and
        replaceable.
  \item \textbf{Semantic 739-dimensional observation space.} A sensor pre-processing stage
        packages three LIDAR channels (combined, dynamic-only, static-only) of \(240\) rays each
        together with \(19\) geometric and route features (including lane-marking type flags)
        derived from CARLA's Waypoint API. The decomposition of the LIDAR into semantic channels
        mitigates the zig-zag pathology introduced by static guardrails.
  \item \textbf{PPO agent with modern stabilisers.} An implementation of PPO with Generalised
        Advantage Estimation, clipped surrogate updates, value clipping, approximate-KL early
        stopping, cosine learning-rate annealing, return normalisation and a Welford running-mean
        estimator that is applied only to the vector block of the observation.
  \item \textbf{Basic safety shield.} A lightweight, interpretable shield that combines LIDAR
        thresholding with lane-keeping corrections derived from the Waypoint API. The shield
        obeys the minimal-interference principle: its override action is zero whenever the
        situation is safe and grows smoothly with the severity of the violation.
  \item \textbf{Adaptive-horizon safety shield.} A richer shield that classifies the current
        situation into three risk levels, predicts the vehicle's future trajectory with a
        kinematic bicycle model calibrated for the Tesla Model 3, and selects the least invasive
        corrective action from a priority-ordered candidate set. The shield supports a
        pedestrian emergency, a lane-change relaxation rule and a threshold-floor safeguard that
        prevents the \emph{critical} level from rejecting all candidate actions.
  \item \textbf{Dense reward shaper.} A reward wrapper with twelve additive components plus an
        alive bonus, all of them gated by speed when appropriate to avoid the degenerate
        ``stop-and-collect'' optimum. The shaper also implements a lane-change transition guard
        and a shield grace period that stabilise the interaction between the agent, the shield
        and the reward.
  \item \textbf{Performance-driven curriculum with rollback.} A curriculum manager that divides
        the interval \([0, N_{\max}]\) of NPC vehicles into five uniform stages and advances or
        rolls back according to \(100\)-episode rolling-window statistics. Stage transitions
        update the NPC count in place, without reconnecting to CARLA.
  \item \textbf{Live metrics pipeline.} A SQLite-based logger in WAL mode that exposes two
        tables (\emph{per-episode} and \emph{per-update}) and a Streamlit dashboard that reads
        them while training is still in progress. Unknown metric keys are rendered as
        first-class citizens, which makes the dashboard forward-compatible with future
        experiment extensions.
  \item \textbf{Reproducibility and tooling.} Unit tests for the reward shaper that do not
        require a running simulator, a \texttt{ruff}-enforced code style, seed-separated
        training and evaluation protocols, a checkpoint format that includes the observation
        normaliser state, and a reproducibility guide in Appendix~\ref{app:reproducibility}.
\end{enumerate}

\section{Fulfilment of objectives}\label{sec:concl-objectives}

Table~\ref{tab:objectives-traceability} traces the specific objectives of
Section~\ref{sec:intro-objectives} against the chapters in which each of them is addressed and
against the contributions enumerated above.

\begin{table}[h]
  \centering
  \caption{Traceability matrix between specific objectives and the rest of the document.}
  \label{tab:objectives-traceability}
  \small
  \begin{tabular}{@{}llll@{}}
    \toprule
    \textbf{Objective} & \textbf{Chapter(s)} & \textbf{Contribution} & \textbf{Status} \\
    \midrule
    SO1 \textendash\ Gymnasium CARLA wrapper
         & \ref{ch:design}, \ref{ch:implementation}
         & 1, 2
         & Achieved. \\
    SO2 \textendash\ PPO agent with stabilisers
         & \ref{ch:design}, \ref{ch:implementation}
         & 3
         & Achieved. \\
    SO3 \textendash\ Basic shield
         & \ref{ch:design}, \ref{ch:implementation}, \ref{ch:experimentation}
         & 4
         & Achieved. \\
    SO4 \textendash\ Adaptive-horizon shield
         & \ref{ch:design}, \ref{ch:implementation}, \ref{ch:experimentation}
         & 5
         & Achieved. \\
    SO5 \textendash\ Dense reward shaper
         & \ref{ch:design}, \ref{ch:implementation}
         & 6
         & Achieved. \\
    SO6 \textendash\ Curriculum with rollback
         & \ref{ch:design}, \ref{ch:implementation}, \ref{ch:experimentation}
         & 7
         & Achieved. \\
    SO7 \textendash\ Live metrics pipeline
         & \ref{ch:design}, \ref{ch:implementation}
         & 8
         & Achieved. \\
    \bottomrule
  \end{tabular}
\end{table}

The quantitative evaluation that anchors \emph{``Achieved''} for SO3, SO4 and SO6 is deferred to
the final reports of Chapter~\ref{ch:experimentation}; until the TBD entries of that chapter
are filled in, the empirical side of the argument rests on the acceptance criteria of
Table~\ref{tab:acceptance-criteria}.

\section{Lessons learned}\label{sec:concl-lessons}

The development of the framework surfaced a number of design lessons that, in retrospect, would
have saved considerable time if they had been internalised earlier in the process.

\begin{description}
  \item[\textbf{Place the shield above the reward shaper.}] An earlier prototype inverted this
        ordering and the shaper could not read the executed action. The reward-shaping signals
        that depend on the shield (intervention penalty, smoothness suppression during
        override, grace period) were all computed incorrectly until the chain was reordered.
  \item[\textbf{Never normalise the LIDAR.}] Applying the running-mean normaliser to the
        \(720\) LIDAR dimensions was the single biggest source of training instability.
        When the curriculum advances from stage~0 (no NPCs, the dynamic channel is constant at
        \(1.0\)) to stage~1, the running standard deviation of the dynamic channel is close to
        zero and the normalised inputs become extreme outliers. Restricting the normaliser to
        the \(15\)-dimensional vector block eliminated the gradient spikes without any further
        change to the training loop.
  \item[\textbf{Gate the alive bonus by speed.}] An unconditional alive bonus, even a small one,
        opens a local optimum in which the agent stops, does nothing, and earns
        \(0.15\)/step forever. Gating the alive bonus by the speed factor \(g_t\) of
        Equation~\eqref{eq:speed-gate} closes this trap without hurting exploration.
  \item[\textbf{Give the agent a grace period after a shield override.}] A brake-based shield
        intervention drops the vehicle's speed, which would immediately trigger the idle penalty
        and push the agent into a stationary equilibrium. The shield-grace mechanism of
        Section~\ref{sec:design-reward} breaks this feedback loop. Critically, the grace counter
        is re-armed only when the shield fires \emph{and} the counter is already at zero;
        otherwise, an always-active shield would perpetually reset the grace period and silence
        the idle penalty forever.
  \item[\textbf{Prefer in-place updates to reconnection.}] Every early attempt at changing the
        curriculum stage by rebuilding the environment caused CARLA to deadlock. Updating the
        attribute \texttt{num\_npc\_vehicles} on the base environment and letting the change
        take effect at the next \texttt{reset} is both simpler and robust.
  \item[\textbf{Attribute credit to the executed action, not to the proposed one.}] Storing the
        action that the shield let through, and recomputing its log probability under the
        current policy, is the single correction that made shielded PPO actually learn.
        Without it, the gradient estimator pushes the policy towards an action that never
        reached the environment, and training stagnates.
\end{description}

\section{Future work}\label{sec:concl-future}

Several extensions of this framework are natural next steps. They are listed below in order of
increasing research effort.

\begin{itemize}
  \item \textbf{Dynamic bicycle model.} The kinematic model used by the adaptive shield ignores
        tyre slip and lateral dynamics. At higher speeds or in tight turns this approximation
        becomes inaccurate and the predicted trajectory can diverge from what the simulator
        actually produces. Replacing it with a dynamic bicycle model or with a simple
        lateral-dynamics model~\cite{rajamani2012vehicle} would tighten the safety guarantees
        without changing the rest of the architecture.
  \item \textbf{Reachability-based shield.} A principled alternative to the candidate-search
        strategy is Hamilton\textendash Jacobi reachability
        analysis~\cite{bansal2017hj,fisac2019safety}, which precomputes the set of unsafe
        states and projects the agent's action onto the largest safe control invariant set.
        This yields formal safety guarantees but imposes a non-trivial offline computation cost.
  \item \textbf{Camera-based perception.} The observation space currently does not include RGB
        imagery or semantic-segmentation maps. Adding a convolutional perception head trained
        jointly with the policy, or reusing the frozen backbone of an imitation-learning
        pipeline~\cite{chitta2023transfuser}, would enable the agent to perceive traffic lights,
        road signs and crosswalks that are not visible from the semantic LIDAR alone.
  \item \textbf{Urban benchmarks.} The evaluation has been concentrated on Town04, a highway
        map. Urban maps such as Town01\textendash Town03 and Town10HD feature intersections,
        traffic lights and dense pedestrian traffic that would stress-test the pedestrian-
        emergency branch of the adaptive shield. Adapting the curriculum to this richer
        scenario is itself a research question.
  \item \textbf{Multi-agent safe RL.} The current setup has a single RL-controlled ego vehicle
        surrounded by rule-based NPCs. Replacing the NPCs with independent RL agents and
        composing their individual shields into a global safety filter is an open research
        direction at the intersection of Safe RL and multi-agent RL.
  \item \textbf{Constrained policy optimisation.} PPO is used here as a black-box optimiser of
        the shaped reward. A closer integration with Safe RL would replace the reward shaping
        by explicit safety constraints and solve the resulting Constrained MDP with CPO or a
        Lagrangian variant~\cite{achiam2017cpo,stooke2020responsive}. This would remove the
        hand-tuning of reward weights at the cost of a more complex optimiser.
  \item \textbf{Benchmarking against other shields.} The MetaDrive and HighwayEnv ecosystems
        offer a variety of shielding baselines. Porting those baselines to CARLA and comparing
        them on a shared benchmark would provide an external reference for the adaptive shield
        proposed here.
\end{itemize}

The framework has been designed with these extensions in mind. Each of them should be achievable
by introducing a new wrapper in the chain of Figure~\ref{fig:wrapper-chain} or by replacing the
existing \texttt{Shield} or \texttt{RewardShaper} classes, without modifying the remaining
components.


% ============================================================
%   BACKMATTER  (bibliography and appendices)
% ============================================================
\backmatter

\bibliografia{referencias}

\appendix
% ============================================================
%   Appendix A  \textemdash\  Hyperparameters and Default Configuration
% ============================================================

\chapter{Hyperparameters and Default Configuration}\label{app:hyperparameters}

This appendix collects the default values of every hyperparameter exposed by the command-line
interfaces of \texttt{main\_train.py} and \texttt{main\_eval.py}, together with the constructor
defaults of the reward shaper, the two shields and the curriculum manager.

\section{PPO agent}\label{app:hp-ppo}

\begin{table}[h]
  \centering
  \caption{Default PPO hyperparameters.}
  \label{tab:hp-ppo}
  \small
  \begin{tabular}{@{}llc@{}}
    \toprule
    \textbf{Parameter} & \textbf{CLI flag / constructor} & \textbf{Default} \\
    \midrule
    Learning rate (initial)        & \texttt{--lr}                & \(1\times 10^{-4}\) \\
    Learning rate (final)          & internal                     & \(1\times 10^{-5}\) \\
    LR schedule                    & internal                     & Cosine annealing \\
    Discount factor \(\gamma\)     & constructor                  & \(0.99\) \\
    GAE \(\lambda\)                & constructor                  & \(0.95\) \\
    PPO clip \(\epsilon\)          & constructor                  & \(0.20\) \\
    Value-loss coefficient         & \texttt{--value\_loss\_coef} & \(0.25\) \\
    Entropy coefficient            & \texttt{--entropy\_coef}     & \(0.02\) \\
    Rollout length                 & \texttt{--update\_timestep}  & \(2\,048\) \\
    PPO epochs per update          & \texttt{--k\_epochs}         & \(10\) \\
    Approximate-KL target          & \texttt{--kl\_target}        & \(0.08\) (0 disables) \\
    Max gradient \(\ell_2\)-norm   & constructor                  & \(1.0\) \\
    Value clipping \(\xi\)         & constructor                  & disabled \\
    Hidden size                    & constructor                  & \(256\) \\
    LIDAR encoder dims             & constructor                  & \(720 \to 256 \to 64\) \\
    \(\log\sigma\) initial value   & constructor                  & \(-0.70\) \\
    \(\log\sigma\) bounds          & constructor                  & \([-3.0, -0.2]\) \\
    Obs.\ normaliser clip          & \texttt{RunningMeanStd}      & \([-5, 5]\) \\
    Obs.\ normalisation enabled    & \texttt{--obs-norm}          & enabled (vector block only) \\
    Return normalisation           & internal                     & enabled (std only) \\
    \bottomrule
  \end{tabular}
\end{table}

\section{Environment}\label{app:hp-env}

\begin{table}[h]
  \centering
  \caption{Default environment hyperparameters.}
  \label{tab:hp-env}
  \small
  \begin{tabular}{@{}llc@{}}
    \toprule
    \textbf{Parameter} & \textbf{CLI flag / constructor} & \textbf{Default} \\
    \midrule
    CARLA host                & \texttt{--host}             & \texttt{localhost} \\
    CARLA port                & \texttt{--port}             & \(2\,000\) \\
    Traffic Manager port      & \texttt{--tm\_port}         & \(8\,000\) \\
    Map                       & \texttt{--map}              & \texttt{Town04} \\
    Number of NPC vehicles    & \texttt{--num\_npc}         & \(40\) \\
    Weather preset            & \texttt{--weather}          & \texttt{ClearNoon} \\
    Target speed              & \texttt{--target\_speed\_kmh} & \(50\,\mathrm{km/h}\) \\
    Success distance          & \texttt{--success\_distance}  & \(250\,\mathrm{m}\) \\
    Synchronous mode          & constructor                 & \texttt{True} \\
    Simulator time step       & constructor                 & \(0.05\,\mathrm{s}\) (\(20\,\mathrm{Hz}\)) \\
    Number of LIDAR rays      & constructor                 & \(240\) (per channel) \\
    LIDAR range               & constructor                 & \(50\,\mathrm{m}\) \\
    Max steps per episode     & \texttt{--max\_steps}       & \(1\,000\) \\
    Training seed             & \texttt{--seed}             & \(42\) \\
    Evaluation seed           & constructor (eval)          & \(100\) \\
    Base success reward       & constructor                 & \(+30\) \\
    Base crash penalty        & constructor                 & \(-10\) \\
    Base off-road penalty     & constructor                 & \(-20\) \\
    \bottomrule
  \end{tabular}
\end{table}

\section{Reward shaper}\label{app:hp-reward}

\begin{table}[h]
  \centering
  \caption{Default weights of the reward shaper. Symbols refer to the equations of
           Section~\ref{sec:design-reward}.}
  \label{tab:hp-reward}
  \small
  \begin{tabular}{@{}lllc@{}}
    \toprule
    \textbf{Symbol} & \textbf{Description} & \textbf{CLI flag} & \textbf{Default} \\
    \midrule
    \(w_{\mathrm{v}}\)       & Speed-bonus scale           & \texttt{--speed\_weight}               & \(0.08\) \\
    \(w_{\mathrm{sm}}\)      & Smoothness penalty scale    & \texttt{--smoothness\_weight}          & \(0.10\) \\
    \(w_{\mathrm{lc}}\)      & Lane-centring scale         & \texttt{--lane\_centering\_weight}     & \(0.15\) \\
    \(w_{\mathrm{h}}\)       & Heading-alignment scale     & constructor                            & \(0.04\) \\
    \(w_{\mathrm{inv}}\)     & Lane-invasion penalty       & \texttt{--lane\_invasion\_penalty}     & \(0.25\) \\
    \(w_{\mathrm{edge}}\)    & Edge-warning scale          & constructor                            & \(0.30\) \\
    \(w_{\mathrm{prog}}\)    & Progress-milestone bonus    & constructor                            & \(0.30\) \\
    \(w_{\mathrm{wh}}\)      & Wrong-heading penalty       & constructor                            & \(0.50\) \\
    \(w_{\mathrm{road}}\)    & Off-road penalty (shaper)   & \texttt{--off\_road\_penalty}          & \(3.00\) \\
    \(w_{\mathrm{sh}}\)      & Shield-intervention penalty & \texttt{--shield\_intervention\_penalty}& \(0.00\) \\
    \(w_{\mathrm{no-sh}}\)   & Shield-not-active bonus     & constructor                            & \(0.02\) \\
    \(w_{\mathrm{idle}}\)    & Idle penalty scale          & \texttt{--idle\_penalty\_weight}       & \(0.04\) \\
    \(w_{\mathrm{dr}}\)      & Drift penalty scale         & constructor                            & \(0.08\) \\
    \(c_{\mathrm{alive}}\)   & Alive bonus                 & constructor                            & \(0.15\) \\
    \(v_{\min}\)             & Minimum moving speed        & \texttt{--min\_moving\_speed\_kmh}     & \(5.0\,\mathrm{km/h}\) \\
    \(v_g\)                  & Full speed-gate speed       & constructor                            & \(10.0\,\mathrm{km/h}\) \\
    \(c_{\mathrm{curv}}\)    & Curvature speed-target scale & constructor                           & \(0.4\) \\
    \(\eta\)                 & Speed-limit tolerance       & constructor                            & \(0.05\) \\
    \(G\)                    & Shield grace-period length  & constructor                            & \(10\) steps \\
    \(\Delta_{\mathrm{prog}}\)& Progress-milestone length  & constant (\texttt{PROGRESS\_MILESTONE\_M}) & \(25\,\mathrm{m}\) \\
    \bottomrule
  \end{tabular}
\end{table}

\section{Safety shields}\label{app:hp-shields}

\begin{table}[h]
  \centering
  \caption{Default hyperparameters of both shields. ``Basic only'' and ``Adaptive only'' columns
           indicate parameters that exist solely in one of the shields.}
  \label{tab:hp-shields}
  \small
  \begin{tabular}{@{}lllc@{}}
    \toprule
    \textbf{Parameter} & \textbf{CLI flag} & \textbf{Scope} & \textbf{Default} \\
    \midrule
    Front LIDAR threshold    & \texttt{--front\_threshold}   & both            & \(0.15\) \\
    Side LIDAR threshold     & \texttt{--side\_threshold}    & both            & \(0.04\) \\
    Lateral offset threshold & \texttt{--lateral\_threshold} & both            & \(0.65\) (train), \(0.82\) (eval) \\
    Heading threshold        & constructor                   & basic only      & \(0.60\) \\
    Steering gain \(k_s\)    & constructor                   & both            & \(2.5\) \\
    Brake gain \(k_b\)       & constructor                   & both            & \(8.0\) \\
    Pedestrian emergency dist.& constant                     & adaptive only   & \(4.0\,\mathrm{m}\) \\
    Risk-level horizons (\(H\))& internal                    & adaptive only   & \(1, 5, 10\) \\
    Threshold multipliers (\(\kappa\))& internal              & adaptive only   & \(1.0, 1.5, 2.0\) \\
    Frontal-risk boundaries  & internal                      & adaptive only   & \(0.50,\, 0.20\) \\
    Lateral-risk boundaries  & internal                      & adaptive only   & \(0.70,\, 0.85\) \\
    Lane-change lateral relaxation & internal                & adaptive only   & \(1.2\) \\
    Minimum lateral floor    & internal                      & adaptive only   & \(0.45\) \\
    \bottomrule
  \end{tabular}
\end{table}

\section{Curriculum manager}\label{app:hp-curriculum}

\begin{table}[h]
  \centering
  \caption{Default hyperparameters of the curriculum manager.}
  \label{tab:hp-curriculum}
  \small
  \begin{tabular}{@{}llc@{}}
    \toprule
    \textbf{Parameter} & \textbf{CLI flag / constructor} & \textbf{Default} \\
    \midrule
    Curriculum enabled       & \texttt{--curriculum} / \texttt{--no-curriculum} & enabled \\
    Maximum NPC count        & \texttt{--num\_npc}          & \(40\) \\
    Number of stages         & internal                     & \(5\) (uniform in \([0, N_{\max}]\)) \\
    Min.\ episodes per stage & constructor                  & \(100\) \\
    Rollback patience        & constructor                  & \(50\) \\
    Rollback counter decay   & constructor                  & \(-1\) per non-collapsing episode \\
    \bottomrule
  \end{tabular}
\end{table}

% ============================================================
%   Appendix B  --  Observation Space Breakdown
% ============================================================

\chapter{Observation Space Breakdown}\label{app:observation}

This appendix describes every one of the \(739\) dimensions of the observation vector
\(s_t\). Indices are zero-based and match the slicing conventions used in
\texttt{src/CARLA/Env/carla\_env.py}.

\section{LIDAR block (dimensions 0\textendash 719)}\label{app:obs-lidar}

The first \(720\) dimensions contain three semantic LIDAR channels of \(240\) rays each. Each ray
is normalised to \([0, 1]\), where \(1\) means a clear ray over the whole \(50\,\mathrm{m}\)
range and values close to \(0\) mean an obstacle within \({\approx}\,1\,\mathrm{m}\). Ray
\(k\in\{0,\ldots,239\}\) points towards the azimuth angle \(\phi_k = k \cdot 1.5^\circ\)
measured counter-clockwise from the vehicle's forward direction.

\begin{table}[H]
  \centering
  \caption{LIDAR block of the observation vector.}
  \label{tab:obs-lidar}
  \small
  \begin{tabular}{@{}lll@{}}
    \toprule
    \textbf{Indices} & \textbf{Channel} & \textbf{Description} \\
    \midrule
    \([0, 240)\)   & Combined  & All obstacles (vehicles, pedestrians, static, road edges). \\
    \([240, 480)\) & Dynamic   & Only vehicles and pedestrians. \\
    \([480, 720)\) & Static    & Only walls, barriers, poles and similar static objects. \\
    \bottomrule
  \end{tabular}
\end{table}

\section{Lane features (dimensions 720\textendash 727)}\label{app:obs-lane}

The next \(8\) dimensions are lane features computed from CARLA's Waypoint API at the
vehicle's current position.

\begin{table}[H]
  \centering
  \caption{Lane features.}
  \label{tab:obs-lane}
  \small
  \begin{tabular}{@{}cllp{6cm}@{}}
    \toprule
    \textbf{Idx.} & \textbf{Name} & \textbf{Range} & \textbf{Meaning} \\
    \midrule
    720 & \texttt{lateral\_offset\_norm}    & \([-1, 1]\) & Signed lateral offset of the
        vehicle with respect to the lane centre, normalised by the lane half-width. \\
    721 & \texttt{heading\_error\_norm}     & \([-1, 1]\) & Signed heading error between the
        vehicle and the lane tangent, normalised by \(180^\circ\). \\
    722 & \texttt{on\_edge\_warning}        & \([0, 1]\) & Continuous warning that grows
        quadratically as the closest lane edge approaches. \\
    723 & \texttt{lane\_width\_norm}        & \([0, 1]\) & Lane width at the current
        position, normalised by a reference width of \(4.5\,\mathrm{m}\). \\
    724 & \texttt{dist\_left\_edge\_norm}   & \([0, 1]\) & Distance to the left lane edge,
        normalised by the lane width. \\
    725 & \texttt{dist\_right\_edge\_norm}  & \([0, 1]\) & Distance to the right lane edge,
        normalised by the lane width. \\
    726 & \texttt{lane\_change\_left}        & \(\{0, 1\}\) & \(1\) if a left-side lane change is
        permitted at the current position (only the left direction is exposed). \\
    727 & \texttt{road\_curvature\_norm}    & \([-1, 1]\) & Estimated signed road curvature,
        normalised by a reference curvature. \\
    \bottomrule
  \end{tabular}
\end{table}

\section{Lane-marking type (dimensions 728\textendash 731)}\label{app:obs-marking}

Four binary flags classify the lane-marking type on each side of the vehicle, derived from
CARLA's \texttt{LaneMarkingType} at the current position.

\begin{table}[H]
  \centering
  \caption{Lane-marking type features.}
  \label{tab:obs-marking}
  \small
  \begin{tabular}{@{}cllp{6cm}@{}}
    \toprule
    \textbf{Idx.} & \textbf{Name} & \textbf{Range} & \textbf{Meaning} \\
    \midrule
    728 & \texttt{solid\_left}   & \(\{0, 1\}\) & \(1\) if the left lane marking is solid
        (crossing is forbidden). \\
    729 & \texttt{solid\_right}  & \(\{0, 1\}\) & \(1\) if the right lane marking is solid. \\
    730 & \texttt{dashed\_left}  & \(\{0, 1\}\) & \(1\) if the left marking is dashed
        (lane change permitted). \\
    731 & \texttt{dashed\_right} & \(\{0, 1\}\) & \(1\) if the right marking is dashed. \\
    \bottomrule
  \end{tabular}
\end{table}

These flags allow the policy to distinguish legal lane changes (dashed marking) from illegal
solid-line crossings without relying on CARLA's lane-invasion event alone.

\section{Vehicle state (dimensions 732\textendash 733)}\label{app:obs-vehicle}

\begin{table}[H]
  \centering
  \caption{Vehicle state features.}
  \label{tab:obs-vehicle}
  \small
  \begin{tabular}{@{}cllp{6cm}@{}}
    \toprule
    \textbf{Idx.} & \textbf{Name} & \textbf{Range} & \textbf{Meaning} \\
    \midrule
    732 & \texttt{speed\_norm} & \([0, 1]\) & Vehicle speed normalised by a dynamic reference:
        \(\max(\text{speed\_limit} \times 1.5,\; 10\,\mathrm{km/h})\). \\
    733 & \texttt{steering}    & \([-1, 1]\) & Last steering command actually applied to the
        vehicle (after shield correction). \\
    \bottomrule
  \end{tabular}
\end{table}

\section{Route information (dimensions 734\textendash 738)}\label{app:obs-route}

The final \(5\) dimensions provide look-ahead information along the route graph.

\begin{table}[H]
  \centering
  \caption{Route features.}
  \label{tab:obs-route}
  \small
  \begin{tabular}{@{}cllp{6cm}@{}}
    \toprule
    \textbf{Idx.} & \textbf{Name} & \textbf{Range} & \textbf{Meaning} \\
    \midrule
    734 & \texttt{next\_wp\_angle\_norm} & \([-1, 1]\) & Signed angle to the next waypoint
        \(5\,\mathrm{m}\) ahead, normalised by \(180^\circ\). \\
    735 & \texttt{wp\_angle\_20m\_norm} & \([-1, 1]\) & Signed angle to the waypoint
        \(20\,\mathrm{m}\) ahead, normalised by \(180^\circ\) (look-ahead for smooth turning). \\
    736 & \texttt{progress\_norm}       & \([0, 1]\) & Fraction of the success distance
        traversed so far during the current episode. \\
    737 & \texttt{speed\_limit\_norm}   & \([0, 1]\) & Current speed limit normalised by the
        reference speed. \\
    738 & \texttt{speed\_ratio}         & \([0, 2]\) & Ratio of the vehicle speed to the
        current speed limit (clipped at \(2\)). \\
    \bottomrule
  \end{tabular}
\end{table}

\section{Summary}\label{app:obs-summary}

The full observation vector has the structure
\begin{equation*}
  s_t \;=\; \underbrace{\big[\, \text{LIDAR}^{\text{comb}}_t \,\big]}_{240}
         \;\|\; \underbrace{\big[\, \text{LIDAR}^{\text{dyn}}_t \,\big]}_{240}
         \;\|\; \underbrace{\big[\, \text{LIDAR}^{\text{stat}}_t \,\big]}_{240}
         \;\|\; \underbrace{\big[\, \text{Lane}_t \,\big]}_{8}
         \;\|\; \underbrace{\big[\, \text{Marking}_t \,\big]}_{4}
         \;\|\; \underbrace{\big[\, \text{Veh}_t \,\big]}_{2}
         \;\|\; \underbrace{\big[\, \text{Route}_t \,\big]}_{5}
         \;\in\; \mathbb{R}^{739}.
\end{equation*}
Only the rightmost \(19\) dimensions (lane, marking, vehicle and route) are normalised online
by the running-mean estimator of Section~\ref{sec:design-obsnorm}; the \(720\) LIDAR dimensions
are already in \([0,1]\) by construction and are left untouched.

% ============================================================
%   Appendix C  \textemdash\  Reproducibility Guide
% ============================================================

\chapter{Reproducibility Guide}\label{app:reproducibility}

This appendix lists the exact commands needed to reproduce every experiment reported in
Chapter~\ref{ch:experimentation}, together with the SQL queries used to extract aggregate
statistics from the per-run SQLite databases.

\section{Environment setup}\label{app:repro-setup}

\begin{enumerate}
  \item Install the CARLA 0.9.16 binary distribution. On Windows, launch the server with the
        command shown in Figure~\ref{fig:repro-launch}.
  \item Activate the project virtual environment (for example, the shared venv at
        \texttt{C:\textbackslash Users\textbackslash Isaac\textbackslash Documents\textbackslash
        TFG\textbackslash Shielded-RL\textbackslash TFG-RL}).
  \item Verify that the Python dependencies are installed (PyTorch with CUDA support,
        Gymnasium, NumPy, Matplotlib, Streamlit; SQLite is part of the Python standard
        library).
\end{enumerate}

\begin{figure}[H]
  \centering
  \includegraphics[width=0.9\textwidth]{imgs/code/repro_launch}
  \caption{CARLA server launch command. The \texttt{-RenderOffScreen} flag disables the GPU
           renderer for headless training; \texttt{-carla-port=2000} sets the TCP port.}
  \label{fig:repro-launch}
\end{figure}

\section{Training commands}\label{app:repro-training}

Figure~\ref{fig:repro-training} shows the commands used to reproduce the three training
configurations of Chapter~\ref{ch:experimentation}.

\begin{figure}[H]
  \centering
  \includegraphics[width=0.9\textwidth]{imgs/code/repro_training}
  \caption{Training commands for the baseline (no shield), basic shield, adaptive shield and
           curriculum-ablation configurations. Each run writes metrics to
           \texttt{runs/<model>\_<shield>\_<timestamp>/metrics.sqlite} and checkpoints to
           \texttt{data/models/<model>\_<shield>\_\{best,final\}.pth}.}
  \label{fig:repro-training}
\end{figure}

\section{Evaluation commands}\label{app:repro-evaluation}

Figure~\ref{fig:repro-evaluation} shows the evaluation commands, both in headless mode (metrics
only) and with the Matplotlib dashboard.

\begin{figure}[H]
  \centering
  \includegraphics[width=0.9\textwidth]{imgs/code/repro_evaluation}
  \caption{Evaluation commands. The shield configuration must match the one used during
           training, as a policy trained with a shield has learned to rely on its interventions.}
  \label{fig:repro-evaluation}
\end{figure}

\section{Live dashboard}\label{app:repro-dashboard}

Figure~\ref{fig:repro-dashboard} shows how to launch the Streamlit dashboard. The dashboard
reads every \texttt{runs/*/metrics.sqlite} found under the project root and allows the user
to select the run(s) to visualise. Because the database is opened in WAL mode, the dashboard
can be launched while a training run is still in progress.

\begin{figure}[H]
  \centering
  \includegraphics[width=0.6\textwidth]{imgs/code/repro_dashboard}
  \caption{Streamlit dashboard launch command.}
  \label{fig:repro-dashboard}
\end{figure}

\section{Query recipes}\label{app:repro-queries}

The following queries document how the empirical numbers in the tables of
Chapter~\ref{ch:experimentation} are computed from the SQLite databases.

\begin{figure}[H]
  \centering
  \includegraphics[width=0.9\textwidth]{imgs/code/repro_sql_rates}
  \caption{SQL queries for the final 100-episode rolling reward and safety rates (crash and
           off-road). Replace the \texttt{key} strings with the metric of interest.}
  \label{fig:repro-sql-rates}
\end{figure}

\begin{figure}[H]
  \centering
  \includegraphics[width=0.9\textwidth]{imgs/code/repro_sql_risklevels}
  \caption{SQL query for the fraction of steps spent in each risk level (adaptive shield only).}
  \label{fig:repro-sql-risklevels}
\end{figure}

\begin{figure}[H]
  \centering
  \includegraphics[width=0.9\textwidth]{imgs/code/repro_sql_interventions}
  \caption{SQL query for the average number of shield interventions per episode by threat
           category (adaptive shield only).}
  \label{fig:repro-sql-interventions}
\end{figure}

\begin{figure}[H]
  \centering
  \includegraphics[width=0.9\textwidth]{imgs/code/repro_sql_curriculum}
  \caption{SQL query for curriculum stage transitions. The \texttt{LAG} window function
           identifies advance and rollback events by comparing consecutive NPC counts.}
  \label{fig:repro-sql-curriculum}
\end{figure}

\section{Checkpoint inspection}\label{app:repro-ckpt}

The utility class \texttt{ModelManager} in \texttt{utils/} exposes helpers to list, compare and
inspect checkpoints. Figure~\ref{fig:repro-modelmanager} shows a quick sanity-check invocation.

\begin{figure}[H]
  \centering
  \includegraphics[width=0.75\textwidth]{imgs/code/repro_modelmanager}
  \caption{Checkpoint inspection with \texttt{ModelManager}. The \texttt{print\_model\_info}
           call reports the network architecture, checkpoint format version, training seed, and
           the state of the running-mean normaliser.}
  \label{fig:repro-modelmanager}
\end{figure}

\section{Unit tests}\label{app:repro-tests}

The reward shaper tests (Figure~\ref{fig:repro-pytest}) run without a CARLA server. They exercise
each component of the shaped reward of Equation~\eqref{eq:shaped-reward} against a crafted
\texttt{info} dictionary and ensure that disabling a weight zeroes out the corresponding
component. They are run locally before every commit that touches \texttt{src/reward\_shaper.py}.

\begin{figure}[H]
  \centering
  \includegraphics[width=0.7\textwidth]{imgs/code/repro_pytest}
  \caption{Running the reward-shaper unit tests. The suite completes in under one second and
           requires no running CARLA server.}
  \label{fig:repro-pytest}
\end{figure}

% ============================================================
%   Appendix E  --  Mathematical Formulation
% ============================================================

\chapter{Mathematical Formulation}\label{app:math}

For readability, several quantities in the main text are presented through figures and tables
rather than equations. This appendix reproduces their exact closed forms for reference, with every
symbol defined. In each case the figure cited in the main text is the primary, intuition-carrying
presentation, and the equations below are its algebraic counterpart.

\section{Risk-level classification (adaptive shield)}\label{app:math-risk}

The frontal-distance criterion of the adaptive shield (Section~\ref{sec:design-adaptive-shield},
Figure~\ref{fig:risk-levels} and Table~\ref{tab:risk-params}) maps the minimum of the dynamic-only
LIDAR front sector \(d^{\mathrm{dyn}}_t\in[0,1]\) to a discrete risk level:
\begin{equation}
  \mathrm{level}^{\mathrm{frontal}}_t \;=\;
     \begin{cases}
       \mathrm{safe}     & \text{if } d^{\mathrm{dyn}}_t > 0.50,\\
       \mathrm{warning}  & \text{if } 0.20 < d^{\mathrm{dyn}}_t \leq 0.50,\\
       \mathrm{critical} & \text{if } d^{\mathrm{dyn}}_t \leq 0.20,
     \end{cases}
  \label{eq:shield-risk-front}
\end{equation}
where a higher \(d^{\mathrm{dyn}}_t\) denotes a more distant in-path obstacle (\(1\) is a clear
lane). The lateral-position criterion applies the analogous thresholds of
Table~\ref{tab:risk-params} to the absolute lateral offset \(|\ell_t|\) and the absolute heading
error \(|\psi_t|\); the overall risk level is the worst (most severe) of the two criteria.

\section{Lane-keep gate (reward shaper)}\label{app:math-gate}

The lane-keep gate of Section~\ref{sec:design-reward}, visualised in
Figure~\ref{fig:lane-keep-gate}, multiplies the forward-progress bonuses by
\begin{equation}
  \mu_t \;=\;
  \underbrace{\max(1 - |\ell_t|,\, 0)}_{\text{centering factor}}
  \;\times\;
  \underbrace{\max\!\left(1 - \frac{|\psi_t|}{20^\circ},\, 0\right)}_{\text{heading factor}}
  \;\in\;[0,1],
  \label{eq:lane-keep-gate}
\end{equation}
where \(\ell_t\) is the normalised lateral offset (\(0\) at the lane centre, \(\pm1\) at the edges)
and \(\psi_t\) is the heading error in degrees. The gate equals \(1\) only when the vehicle is both
centred and aligned, and \(0\) as soon as \(|\ell_t|\geq1\) or \(|\psi_t|\geq20^\circ\). During a
permitted lane change (\(\ell_t>0.5\) with CARLA reporting the change allowed) \(\mu_t\) is floored
at \(0.3\).

\section{Shaped reward (reward shaper)}\label{app:math-reward}

The total shaped reward of Section~\ref{sec:design-reward}, summarised in
Figure~\ref{fig:reward-overview} and Table~\ref{tab:reward-terms}, is the sum of the sparse base
reward and the dense bonus and penalty terms:
\begin{equation}\label{eq:shaped-reward}
  \begin{split}
  \tilde{r}_t \;=\;{} & r^{\mathrm{base}}_t
    + r^{\mathrm{prog}}_t + r^{\mathrm{accel}}_t + r^{\mathrm{speed}}_t
    + r^{\mathrm{lane}}_t + r^{\mathrm{head}}_t + r^{\mathrm{mile}}_t \\
    & - p^{\mathrm{smooth}}_t - p^{\mathrm{inv}}_t - p^{\mathrm{solid}}_t
      - p^{\mathrm{edge}}_t - p^{\mathrm{wrong\text{-}h}}_t
      - p^{\mathrm{idle}}_t - p^{\mathrm{drift}}_t
      - p^{\mathrm{lc}}_t - p^{\mathrm{shield}}_t.
  \end{split}
\end{equation}
Each symbol is defined in Table~\ref{tab:reward-terms}; the default weights are listed in
Appendix~\ref{app:hyperparameters}, Table~\ref{tab:hp-reward}.

\section{Shield permissiveness (LLM meta-controller)}\label{app:math-permissiveness}

The single permissiveness scalar \(p\in[0,1]\) of Section~\ref{sec:dep-llm}, illustrated in
Figure~\ref{fig:permissiveness-ramp}, sets every tunable threshold \(\tau_i\) of the adaptive
shield by linear interpolation between a strict and a permissive endpoint:
\begin{equation}\label{eq:permissiveness}
  \tau_i(p) \;=\; \tau_i^{\mathrm{strict}} \;+\; (1-p)\,\bigl(\tau_i^{\mathrm{perm}} - \tau_i^{\mathrm{strict}}\bigr),
\end{equation}
where \(\tau_i^{\mathrm{strict}}\) and \(\tau_i^{\mathrm{perm}}\) are the strict and permissive
endpoints of threshold \(i\). Thus \(p=1\) recovers the production shield (every
\(\tau_i=\tau_i^{\mathrm{strict}}\)) and \(p=0\) the maximally weaned shield (every
\(\tau_i=\tau_i^{\mathrm{perm}}\)).

% ============================================================
%   Appendix D  --  Transparent Use of AI
% ============================================================

\chapter{Transparent use of AI}\label{app:ai-usage}

\begin{figure}[H]
  \centering
  \includegraphics[width=0.6\textwidth]{imgs/ai-label_banner-assisted-by-ai.png}
  \caption{Disclosure banner indicating that parts of this project's development and
           documentation were produced with the assistance of AI tools, under continuous
           human supervision.}
  \label{fig:ai-banner}
\end{figure}

In line with the growing institutional expectation that academic work disclose the use
of generative AI tools, this appendix documents how, where, and why such tools were
employed during the development of \textit{CALA\_SafeRL}, the analysis of experimental
results, and the writing of this memoir.

\paragraph{Human oversight statement.}
It must be stated explicitly that \textbf{every output produced with the assistance of
an AI tool, whether source code, data analysis, or written text, was reviewed, tested,
and revised by the author before being incorporated into the project or this document}.
AI tools were used strictly in an assistive capacity, as accelerators for routine or
exploratory work, debugging aids, and proofreading support. At no point did an AI tool
act as an autonomous decision-maker with respect to the project's scientific
conclusions, the design of the safety shield, or the interpretation of results. Code
suggested by an assistant was run and inspected by the author; quantitative claims were
cross-checked against the underlying SQLite metric databases before being reported. This
principle of continuous human oversight is reiterated at the end of this appendix.

\section{Claude Code as a programming assistant}\label{app:ai-claude}

\textit{Claude Code}~\cite{anthropic2025claudecode} is an agentic command-line coding
assistant developed by Anthropic, used as the primary programming assistant during
implementation and experimentation. It operated directly inside the repository, with
access to the source tree, the test suite, and the per-run metrics database.

\paragraph{Debugging the wrapper chain.}
The layered \texttt{gymnasium} wrapper chain connecting the PPO agent to the CARLA
simulator (Chapter~\ref{ch:design}) is a frequent source of subtle bugs: the
agent interacts with \texttt{CarlaRewardShaper}, which wraps the active safety shield,
which in turn wraps \texttt{CarlaEnv}. An error in how one layer modifies the action,
observation, or reward can surface as a confusing symptom several layers away, for
example a mismatch between the policy's log-probabilities and the action actually
executed in the simulator (\texttt{info["executed\_action"]}).

The typical workflow was to describe a symptom in plain language, ask Claude Code to
trace the relevant code path across the wrapper boundaries, and then verify the proposed
explanation manually against the source code and the relevant unit test in
\texttt{tests/} before applying and re-testing a fix. This pattern, describe, trace,
verify, fix, was repeated for most of the non-trivial bugs in the shield credit-assignment
logic and in the continuous action-space handling.

\paragraph{Continuous action space and curriculum learning.}
Debugging the saturation of \texttt{log\_std} (the effective policy standard deviation
becoming pinned at its upper bound during early training) required correlating
per-update telemetry (\texttt{log\_std\_steering\_raw}, \texttt{log\_std\_throttle\_raw},
\texttt{log\_std\_saturated\_fraction}) with the entropy-coefficient schedule and the
\texttt{tanh}-squashed parameterisation in \texttt{ActorCritic.log\_std()}. Claude Code
was used to write throwaway scripts that queried \texttt{metrics.sqlite} and plotted
these quantities across runs, which sped up forming and discarding hypotheses. The
resulting decision, to replace the hard clamp with the smooth squashed parameterisation
described in Chapter~\ref{ch:design}, was the author's; the assistant's role was
limited to producing the exploratory tooling.

A similar use applied to the multi-phase curriculum in \texttt{train\_curriculum.py},
where each phase warm-starts from the previous phase's checkpoint. Claude Code helped
compare checkpoint metadata (\texttt{obs\_normalizer}, \texttt{ret\_rms},
\texttt{returns\_acc}) across phase boundaries, ruling out checkpoint-compatibility
issues and focusing the investigation on the hyperparameter changes themselves.

\paragraph{Large-scale analysis of PPO training metrics.}
Each training run produces per-episode and per-update SQLite databases
(Chapter~\ref{ch:design}), and many model variants were trained across shield
configurations, curriculum settings, and reward-shaping parameters. Claude Code drafted
SQL queries and analysis scripts to aggregate crash rate, off-road rate, mean reward,
and shielded fraction across runs, producing the comparison tables used in
Chapter~\ref{ch:experimentation}. As before, every aggregate figure reported in this
memoir was independently spot-checked by the author against the raw databases.

\section{Ollama with Gemma as a local, privacy-preserving assistant}\label{app:ai-ollama}

A second class of tasks benefited from running entirely on local hardware, using
\textit{Ollama} to host a \textit{Gemma}~\cite{team2024gemma} model.

\paragraph{Motivation.}
Two reasons motivated a local deployment. First, the per-run \texttt{metrics.sqlite}
databases contain project-specific experimental data, and keeping their analysis on the
local machine avoids any ambiguity about data retention by a third party. Second, the
\texttt{dataVisualizer.py} dashboard (Chapter~\ref{ch:design}) provides live
analysis while training is in progress, and a local model avoids the latency and
connectivity requirements of a remote API during long unattended runs.

\paragraph{Local analysis of training metrics.}
The dashboard optionally invokes the local Gemma model to narrate the Plotly charts
rendered from the live database, for example summarising whether a run's crash-rate
curve has plateaued or shows the rollback behaviour implemented by
\texttt{CurriculumManager}. The prompts include only numerical summaries already
computed locally (rolling crash/off-road rates, mean rewards, curriculum stage
transitions), not raw trajectory data. These summaries were treated as a first pass to
flag runs for closer inspection, never as a substitute for the quantitative analysis in
Chapter~\ref{ch:experimentation}.

\paragraph{Collision-type categorisation.}
A second, ongoing use of the local Gemma deployment is the categorisation of collision
events for the planned LLM-based reward-exposure work
(\texttt{tex/08\_llm\_reward\_exposure.tex}, in preparation). This part of the
project is still work in progress: the goal is to classify recorded collision episodes
into semantically meaningful categories, for example distinguishing a collision with an
oncoming vehicle during a forced lane change from one with a stationary obstacle near the
road edge, in order to inform the preference-based reward signal described in that
chapter. Running this locally allows the pipeline to be iterated on quickly without
per-call costs and keeps the per-collision context on the local machine. The categories
proposed by the model are a starting point for the preference-annotation pipeline and
remain subject to manual spot-checking, a validation step that is part of the work still
to be completed.

\section{Gemini as a writing and proofreading assistant}\label{app:ai-gemini}

\textit{Gemini}~\cite{geminiteam2025gemini}, developed by Google, was used during the
writing of this memoir as a structural and proofreading assistant, confined to the
presentation layer of the document.

Draft sections, written first by the author based directly on the implementation and on
the results in Chapter~\ref{ch:experimentation}, were passed to Gemini with requests
such as checking clarity and terminology consistency, or suggesting a clearer paragraph
structure without changing the technical content. Suggestions concerned sentence-level
clarity, consistent use of terms such as ``safety shield'', the ordering of paragraphs,
and minor grammatical corrections.

Every suggestion was reviewed by the author, who verified that no rewording altered the
technical meaning of a statement, introduced unsupported claims, or removed hedging
language around still-open questions, such as the work-in-progress status of the
collision-categorisation pipeline described in Section~\ref{app:ai-ollama}. Gemini
therefore improved readability and consistency, but did not influence the technical
content, argumentation, or conclusions of this memoir.

\section{Summary of human oversight}\label{app:ai-summary}

Three AI tools were used during this project, each in a distinct and bounded role:
Claude Code~\cite{anthropic2025claudecode} as a programming and data-analysis assistant
on the source code and metric databases; a locally-hosted Gemma
model~\cite{team2024gemma}, served through Ollama, for live metric narration and, in
ongoing work, collision categorisation; and Gemini~\cite{geminiteam2025gemini} as a
proofreading and structural-editing assistant for this memoir. In every case the role
was assistive, accelerating routine or exploratory tasks without making design
decisions, determining experimental conclusions, or producing final text or code without
human review. \textbf{All code, all quantitative results, and all written content in
this project and in this memoir were ultimately validated, tested, and approved by the
author}, who takes full responsibility for their correctness and content.


\end{document}
